% Money and banking(continue) — Economics Upper Sixth — Cours signé OnBuch+
% Official MINESEC syllabus. Compiler avec : tectonic ECO04-money-and-banking-continue.tex
\newcommand{\DOCMATIERE}{Economics}
\newcommand{\DOCNIVEAU}{Upper Sixth}
\newcommand{\DOCMODULE}{Upper Sixth — Economics}
\newcommand{\DOCLECON}{4}
\newcommand{\DOCTITRE}{Money and banking(continue)}
% OnBuch+ — préambule « cours signés » ENRICHI (compilé avec tectonic / XeLaTeX)
% Système visuel « Pop » d'OnBuch+ : fond crème (jamais blanc), bordures encre
% épaisses, ombres « dures » décalées, titres Archivo Black en capitales.
% Version MATHÉMATIQUES : graphes (pgfplots/tikz), boîtes pédagogiques
% (définition, propriété, méthode, attention, le savais-tu, exercice résolu,
% exercice, TP, bilan), page de garde et signature OnBuch+ ancrée en bas.
%
% Macros attendues dans main.tex AVANT \input{preamble} :
%   \DOCMATIERE  \DOCNIVEAU  \DOCMODULE  \DOCLECON (numéro)  \DOCTITRE
\documentclass[11pt]{article}

\usepackage{fontspec}
\usepackage[english]{babel}
\usepackage[a4paper,margin=2cm,top=3.4cm,bottom=2.8cm,headheight=1.4cm,footskip=1.3cm]{geometry}
\usepackage{amsmath,amssymb,amsfonts}
\usepackage{mathtools} % \xmapsto, \coloneqq... souvent utilisés par les modèles
\usepackage{cancel} % simplifications barrées (bilans de forces)
% \abs{z} (module, valeur absolue) et \norm{u} (norme), utilisés par les modèles
\providecommand{\abs}{}\let\abs\relax\DeclarePairedDelimiter{\abs}{\lvert}{\rvert}
\providecommand{\norm}{}\let\norm\relax\DeclarePairedDelimiter{\norm}{\lVert}{\rVert}
\usepackage[table]{xcolor} % [table] : fournit aussi \rowcolors (alternance de lignes)
\usepackage{pagecolor}
\usepackage{tikz}
\usetikzlibrary{shapes.symbols,circuits.logic.US,circuits.logic.IEC,arrows.meta,positioning,calc,shapes.geometric,shapes.misc,shapes.arrows,decorations.pathmorphing,decorations.pathreplacing,decorations.markings,patterns,shadows,mindmap,trees,fit,backgrounds,matrix,intersections,angles,quotes}
\usepackage{pgfplots}
\usepackage[version=4]{mhchem} % équations nucléaires \ce{^{235}_{92}U}
\usepackage[european,siunitx]{circuitikz} % schémas électriques
\pgfplotsset{compat=1.18}
\usepgfplotslibrary{fillbetween,groupplots} % aires entre courbes (calcul intégral)
\usepackage{enumitem}
\usepackage{fancyhdr}
% [most] charge toutes les bibliothèques tcolorbox (listings, minted, raster,
% documentation...) et gonfle inutilement la mémoire TeX sur un document long
% (~300 boîtes) : on ne charge que ce qui est réellement utilisé.
\usepackage[skins,breakable]{tcolorbox}
\usepackage{graphicx}
\usepackage{colortbl}
\usepackage{booktabs}
\usepackage{tabularx}
\usepackage{diagbox} % case d'en-tête diagonale des tableaux à double entrée
% Colonnes extensibles centrée / gauche / droite (C, L, R), souvent utilisées par les modèles
\newcolumntype{C}{>{\centering\arraybackslash}X}
\newcolumntype{L}{>{\raggedright\arraybackslash}X}
\newcolumntype{R}{>{\raggedleft\arraybackslash}X}
\usepackage{array}
\usepackage{multicol}
\usepackage{siunitx}
% parse-numbers=false : accepte aussi \SI{2,5 \times 10^{-3}}{...}
\sisetup{locale=UK,output-decimal-marker={.},group-minimum-digits=4,parse-numbers=false}
\DeclareSIUnit\ohm{\ensuremath{\symup{\Omega}}} % U+2126 absent de la police
\DeclareSIUnit\division{div} % réglages d'oscilloscope (ms/div, V/div)
\DeclareSIUnit\year{yr}\DeclareSIUnit\annum{yr}\DeclareSIUnit\Ma{Ma}\DeclareSIUnit\tr{tr} % tours (vitesse de rotation, ex. \SI{1500}{\tr\per\minute})
\usepackage{qrcode}
% \degree : commande fréquemment utilisée par les modèles pour un angle ou une
% température en dehors de \SI{}{} (non fournie par les paquets chargés).
\providecommand{\degree}{^{\circ}}
% \qed (fin de démonstration), utilisé par les modèles sans amsthm
\providecommand{\qed}{\hfill$\square$}
% \barre : double barre des valeurs interdites dans un tableau de variations (tkz-tab)
\providecommand{\barre}{\ensuremath{\|}}
% environnement proof (amsthm non chargé) utilisé par les modèles
\ifdefined\proof\else\newenvironment{proof}[1][Proof]{\par\noindent\textit{#1.}\ }{\qed\par}\fi
\usepackage{lastpage}
\usepackage{hyperref}
\hypersetup{hidelinks}

% ============================================================
% Palette « Pop » OnBuch+ — mêmes hex que la classe P (plus_ui.dart)
% ============================================================
\definecolor{poporange}{HTML}{F59321}
\definecolor{popdark}{HTML}{E07A0C}
\definecolor{popcream}{HTML}{FFF6E8}
\definecolor{popink}{HTML}{1C1714}
\definecolor{poppurple}{HTML}{7A5AE0}
\definecolor{popgreen}{HTML}{2E9E6A}
\definecolor{poppink}{HTML}{E0457B}
\definecolor{popblue}{HTML}{2F7DE1}
\definecolor{popgold}{HTML}{C98A12}
\definecolor{popmuted}{HTML}{8A7F76}
\definecolor{popline}{HTML}{EADFD2}
\definecolor{popgray}{HTML}{8A7F76}
\colorlet{popgrey}{popgray}
% Alias tolérants (le modèle nomme parfois les couleurs différemment)
\colorlet{popwhite}{white}
\colorlet{popblack}{popink}
\colorlet{popred}{poppink}
\colorlet{popyellow}{popgold}
% Teintes claires pour fonds de tableaux / zones
\colorlet{popblueL}{popblue!10!white}
\colorlet{poporangeL}{poporange!14!white}
\colorlet{popgreenL}{popgreen!12!white}
\colorlet{poppurpleL}{poppurple!12!white}
\colorlet{poppinkL}{poppink!12!white}
\colorlet{popgoldL}{popgold!14!white}

\pagecolor{popcream}

% ============================================================
% Typographie
% ============================================================
\defaultfontfeatures{Ligatures=TeX}
\setmainfont{PlusJakartaSans-Medium.ttf}[Path=../../../fonts/,BoldFont=PlusJakartaSans-Bold.ttf]
% Maths en sans-serif (Fira Math) avec chiffres et lettres droites de Plus
% Jakarta, pour que les formules se fondent dans le texte.
\usepackage[math-style=ISO,bold-style=ISO]{unicode-math}
\setmathfont{FiraMath-Regular.otf}[Path=../../../fonts/]
\setmathfont{PlusJakartaSans-Medium.ttf}[Path=../../../fonts/,range={up/num,up/{Latin,latin},\mathup}]
% (icomma retiré : décimales avec point en anglais)
% Caractères mathématiques Unicode absents des polices de texte (Plus Jakarta,
% Archivo Black) : rendus par leur équivalent mathématique, en texte comme en
% formule — sinon ils s'affichent en carré vide (ex. « Arithmétique dans ℤ »).
\usepackage{newunicodechar}
\newunicodechar{ℝ}{\ensuremath{\mathbb{R}}}
\newunicodechar{ℤ}{\ensuremath{\mathbb{Z}}}
\newunicodechar{ℂ}{\ensuremath{\mathbb{C}}}
\newunicodechar{ℕ}{\ensuremath{\mathbb{N}}}
\newunicodechar{ℚ}{\ensuremath{\mathbb{Q}}}
\newunicodechar{𝔻}{\ensuremath{\mathbb{D}}}
\newunicodechar{α}{\ensuremath{\alpha}}
\newunicodechar{β}{\ensuremath{\beta}}
\newunicodechar{γ}{\ensuremath{\gamma}}
\newunicodechar{δ}{\ensuremath{\delta}}
\newunicodechar{ε}{\ensuremath{\varepsilon}}
\newunicodechar{θ}{\ensuremath{\theta}}
\newunicodechar{λ}{\ensuremath{\lambda}}
\newunicodechar{μ}{\ensuremath{\mu}}
\newunicodechar{σ}{\ensuremath{\sigma}}
\newunicodechar{τ}{\ensuremath{\tau}}
\newunicodechar{φ}{\ensuremath{\varphi}}
\newunicodechar{ψ}{\ensuremath{\psi}}
\newunicodechar{ω}{\ensuremath{\omega}}
\newunicodechar{Σ}{\ensuremath{\Sigma}}
% majuscules produites par \MakeUppercase dans les titres (α-aminés -> Α)
\newunicodechar{Α}{\ensuremath{\alpha}}
\newunicodechar{Β}{\ensuremath{\beta}}
\newunicodechar{Γ}{\ensuremath{\Gamma}}
\newunicodechar{Δ}{\ensuremath{\Delta}}
\newunicodechar{Ε}{\ensuremath{\varepsilon}}
\newunicodechar{Θ}{\ensuremath{\Theta}}
\newunicodechar{Λ}{\ensuremath{\Lambda}}
\newunicodechar{Μ}{\ensuremath{\mu}}
\newunicodechar{Τ}{\ensuremath{\tau}}
\newunicodechar{Φ}{\ensuremath{\Phi}}
\newunicodechar{Ψ}{\ensuremath{\Psi}}
\newunicodechar{Ω}{\ensuremath{\Omega}}
\newunicodechar{↦}{\ensuremath{\mapsto}}
\newunicodechar{∈}{\ensuremath{\in}}
\newunicodechar{∉}{\ensuremath{\notin}}
\newunicodechar{∧}{\ensuremath{\wedge}}
\newunicodechar{⇔}{\ensuremath{\Leftrightarrow}}
\newunicodechar{⟺}{\ensuremath{\Longleftrightarrow}}
\newunicodechar{⇒}{\ensuremath{\Rightarrow}}
\newunicodechar{⟹}{\ensuremath{\Longrightarrow}}
\newunicodechar{∘}{\ensuremath{\circ}}
\newunicodechar{∪}{\ensuremath{\cup}}
\newunicodechar{∩}{\ensuremath{\cap}}
\newunicodechar{≡}{\ensuremath{\equiv}}
\newunicodechar{⊂}{\ensuremath{\subset}}
\newunicodechar{⊥}{\ensuremath{\perp}}
\newunicodechar{∃}{\ensuremath{\exists}}
\newunicodechar{∀}{\ensuremath{\forall}}
\newunicodechar{∛}{\ensuremath{\sqrt[3]{\,}}}
\newunicodechar{∜}{\ensuremath{\sqrt[4]{\,}}}
\newunicodechar{⃗}{\ensuremath{{}^{\to}}}
\newunicodechar{ⁿ}{\ifmmode^{n}\else\textsuperscript{\ensuremath{n}}\fi}
\newunicodechar{ˣ}{\ifmmode^{x}\else\textsuperscript{\ensuremath{x}}\fi}
\newunicodechar{ᵇ}{\ifmmode^{b}\else\textsuperscript{\ensuremath{b}}\fi}
\newunicodechar{ʳ}{\ifmmode^{r}\else\textsuperscript{\ensuremath{r}}\fi}
\newunicodechar{ᵘ}{\ifmmode^{u}\else\textsuperscript{\ensuremath{u}}\fi}
\newunicodechar{ʸ}{\ifmmode^{y}\else\textsuperscript{\ensuremath{y}}\fi}
\newunicodechar{ᵃ}{\ifmmode^{a}\else\textsuperscript{\ensuremath{a}}\fi}
\newunicodechar{ᵉ}{\ifmmode^{e}\else\textsuperscript{\ensuremath{e}}\fi}
\newunicodechar{ᵏ}{\ifmmode^{k}\else\textsuperscript{\ensuremath{k}}\fi}
\newunicodechar{ᵖ}{\ifmmode^{p}\else\textsuperscript{\ensuremath{p}}\fi}
\newunicodechar{ᴺ}{\ifmmode^{N}\else\textsuperscript{\ensuremath{N}}\fi}
\newunicodechar{ᶜ}{\ifmmode^{c}\else\textsuperscript{\ensuremath{c}}\fi}
\newunicodechar{ʰ}{\ifmmode^{h}\else\textsuperscript{\ensuremath{h}}\fi}
\newunicodechar{ᵠ}{\ifmmode^{q}\else\textsuperscript{\ensuremath{q}}\fi}
\newunicodechar{ₙ}{\ifmmode_{n}\else\textsubscript{\ensuremath{n}}\fi}
\newunicodechar{ₐ}{\ifmmode_{a}\else\textsubscript{\ensuremath{a}}\fi}
\newunicodechar{ᵢ}{\ifmmode_{i}\else\textsubscript{\ensuremath{i}}\fi}
\newunicodechar{ᵣ}{\ifmmode_{r}\else\textsubscript{\ensuremath{r}}\fi}
\newunicodechar{ₖ}{\ifmmode_{k}\else\textsubscript{\ensuremath{k}}\fi}
\newunicodechar{ᵦ}{\ifmmode_{\beta}\else\textsubscript{\ensuremath{\beta}}\fi}
\newunicodechar{ₓ}{\ifmmode_{x}\else\textsubscript{\ensuremath{x}}\fi}
\newfontfamily\popdisplay{ArchivoBlack-Regular.ttf}[Path=../../../fonts/]
\newfontfamily\poplogo{SpaceGrotesk-Bold.ttf}[Path=../../../fonts/]
\newfontfamily\popsemi{PlusJakartaSans-SemiBold.ttf}[Path=../../../fonts/]
\newfontfamily\popxbold{PlusJakartaSans-ExtraBold.ttf}[Path=../../../fonts/]
\newfontfamily\popxboldls{PlusJakartaSans-ExtraBold.ttf}[Path=../../../fonts/,LetterSpace=3.0]

\setlist[enumerate]{leftmargin=1.7em,itemsep=5pt,topsep=4pt}
\setlist[itemize]{leftmargin=1.7em,itemsep=4pt,topsep=4pt,label={\color{poporange}\textbullet}}
\setlength{\parindent}{0pt}
\setlength{\parskip}{5pt}
\renewcommand{\arraystretch}{1.25}

% pgfplots : style Pop par défaut
\pgfplotsset{
  every axis/.append style={axis line style={popink,line width=1pt},tick style={popink},
    grid style={popline,line width=0.5pt},label style={font=\small},tick label style={font=\footnotesize}},
  popcurve/.style={poporange,line width=1.8pt,no markers,smooth},
}
% Même style disponible dans un \draw TikZ simple (hors pgfplots)
\tikzset{popcurve/.style={poporange,line width=1.8pt}}
\tikzset{popgrid/.style={popline,very thin}}
\colorlet{popcurve}{poporange} % le nom du style est parfois utilisé comme couleur

% ============================================================
% Pill / squiggle
% ============================================================
\newcommand{\poppill}[3]{%
  \tikz[baseline=-0.55ex]\node[draw=popink,line width=1.2pt,rounded corners=2.6pt,%
    fill=#2,inner xsep=6.5pt,inner ysep=3pt]{%
    \popxboldls\fontsize{7}{7}\selectfont\color{#3}\MakeUppercase{#1}};%
}
\newcommand{\popsquiggle}[1]{%
  \tikz[baseline=-0.4ex]{
    \draw[#1,line width=2.6pt,line cap=round]
      plot [smooth] coordinates {(0,0.09) (0.34,0.01) (0.68,0.21) (1.02,0.01) (1.36,0.10)};
  }%
}
% Mot-clé surligné (marqueur orange sous le texte)
\newcommand{\cle}[1]{{\popxbold\color{popdark}#1}}

% ============================================================
% En-tête / pied
% ============================================================
\pagestyle{fancy}
\fancyhf{}
\renewcommand{\headrulewidth}{0pt}
\renewcommand{\footrulewidth}{0pt}
\lhead{\raisebox{-0.15ex}{\poplogo\fontsize{13}{13}\selectfont\color{popink}OnBuch\color{poporange}+}}
\rhead{\poppill{\DOCMATIERE\ \textbullet\ \DOCNIVEAU\ \textbullet\ Chapter \DOCLECON}{white}{popink}}
\lfoot{\popxbold\fontsize{7.6}{7.6}\selectfont\color{popmuted}\MakeUppercase{Signed course OnBuch+}\ \textbullet\ {\popsemi\DOCTITRE}}
\rfoot{\tikz[baseline=-0.6ex]\node[rounded corners=8.5pt,fill=popink,minimum height=17pt,minimum width=17pt,inner xsep=5pt,inner sep=0]{\popxbold\fontsize{8}{8}\selectfont\color{popcream}\thepage\,/\,\pageref*{LastPage}};}

% ============================================================
% Titres
% ============================================================
\newcommand{\chaptitre}[2]{%
  \begin{tcolorbox}[enhanced,colback=poporange,colframe=popink,boxrule=2.6pt,arc=11pt,
      drop shadow=popink,left=15pt,right=15pt,top=12pt,bottom=11pt,before skip=3pt,after skip=18pt]
    \poppill{#2}{popink}{popcream}\par\vspace{9pt}
    {\raggedright\hyphenpenalty=10000\exhyphenpenalty=10000\popdisplay\fontsize{22}{24}\selectfont\color{popcream}\MakeUppercase{#1}\par}
  \end{tcolorbox}
}
% Section numérotée : pastille orange avec le numéro + titre Archivo
\newcounter{popsec}
\newcommand{\coursec}[1]{%
  \stepcounter{popsec}\setcounter{popsub}{0}%
  \par\vspace{16pt}\noindent
  \begin{minipage}{\linewidth}
  \tikz[baseline=-0.7ex]\node[draw=popink,line width=1.6pt,fill=poporange,rounded corners=4pt,minimum size=22pt,inner sep=0]{\popdisplay\fontsize{12}{12}\selectfont\color{popcream}\thepopsec};%
  \hspace{8pt}{\popdisplay\fontsize{15}{17}\selectfont\color{popink}\MakeUppercase{#1}}\par
  \vspace{4pt}\noindent{\color{poporange}\rule{36pt}{3.6pt}}
  \end{minipage}\par\nopagebreak\vspace{8pt}
}
\newcounter{popsub}
\newcommand{\courssub}[1]{%
  \stepcounter{popsub}%
  \par\vspace{9pt}\noindent{\popxbold\fontsize{11.5}{13}\selectfont\color{popink}{\color{poporange}\thepopsec.\thepopsub}\hspace{6pt}#1}\par\nopagebreak\vspace{4pt}
}

% ============================================================
% Boîtes pédagogiques — même recette : fond blanc, bordure épaisse colorée,
% ombre dure encre, badge plein. Titre optionnel [#1].
% ============================================================
\tcbset{popbase/.style={breakable,enhanced,colback=white,boxrule=2.3pt,arc=9pt,
  shadow={3pt}{-3pt}{0pt}{popink},left=14pt,right=14pt,top=11pt,bottom=11pt,before skip=11pt,after skip=15pt,
  fontupper=\popsemi\fontsize{10.6}{14.5}\selectfont\color{popink},
  fontlower=\popsemi\fontsize{10.6}{14.5}\selectfont\color{popink}}}
% Coche dessinée (le glyphe ✓ n'existe pas dans les polices du document)
\newcommand{\popcheck}[1]{\tikz[baseline=-0.5ex]\draw[#1,line width=1.6pt,line cap=round,line join=round](-0.13,0.0)--(-0.04,-0.09)--(0.13,0.1);}
\providecommand{\coche}{\popcheck{popgreen}}
\providecommand{\resultat}[1]{\par\smallskip\noindent\fcolorbox{popgreen}{popgreenL}{\parbox{\dimexpr\linewidth-2\fboxsep-2\fboxrule\relax}{\textbf{Result.} #1}}\par\smallskip}
\newcommand{\popboxhead}[3]{\poppill{#1}{#2}{white}\ifx\relax#3\relax\else\hspace{6pt}{\popxbold\fontsize{10.6}{12}\selectfont\color{popink}#3}\fi\par\vspace{7pt}}

\newtcolorbox{aretenir}[1][]{popbase,colframe=popgreen,before upper={\popboxhead{Key points}{popgreen}{#1}}}
\newtcolorbox{exemplebox}[1][]{popbase,colframe=poporange,before upper={\popboxhead{Exemple}{poporange}{#1}}}
\newtcolorbox{definition}[1][]{popbase,colframe=popblue,before upper={\popboxhead{Definition}{popblue}{#1}}}
\newtcolorbox{propriete}[1][]{popbase,colframe=poppurple,before upper={\popboxhead{Property}{poppurple}{#1}}}
\newtcolorbox{methode}[1][]{popbase,colframe=popgold,before upper={\popboxhead{Method}{popgold}{#1}}}
\newtcolorbox{attention}[1][]{popbase,colframe=poppink,before upper={\popboxhead{Warning}{poppink}{#1}}}
\newtcolorbox{experience}[1][]{popbase,colframe=popblue,colback=popblueL,before upper={\popboxhead{Experiment}{popblue}{#1}}}
% « Le savais-tu ? » : carte violette pleine (contexte, culture, Cameroun)
\newtcolorbox{savaistu}[1][]{popbase,colback=poppurple,colframe=popink,
  fontupper=\popsemi\fontsize{10.4}{14.2}\selectfont\color{white},
  before upper={\poppill{Did you know?}{white}{poppurple}\ifx\relax#1\relax\else\hspace{6pt}{\popxbold\color{white}#1}\fi\par\vspace{7pt}}}
% Exercice résolu : énoncé en haut, \tcblower puis la solution sur fond crème
\newcounter{popexo}\newcounter{popexr}
\newtcolorbox{exoresolu}[1][]{popbase,colframe=poporange,colbacklower=popcream,
  segmentation style={popink,line width=1.2pt,dashed},
  before upper={\stepcounter{popexr}\popboxhead{Worked example \thepopexr}{poporange}{#1}},
  before lower={\poppill{Solution}{popink}{popcream}\par\vspace{6pt}}}
% Exercice (non résolu en place) — la correction est dans la partie Corrigés
\newtcolorbox{exercice}[1][]{popbase,colframe=popink,
  before upper={\stepcounter{popexo}\popboxhead{Exercise \thepopexo}{popink}{#1}}}
% Corrigé
\newtcolorbox{corrige}[1][]{popbase,colframe=popgreen,colback=popgreenL,
  before upper={\popboxhead{Solution}{popgreen}{#1}}}
% Objectifs / prérequis / bilan
\newtcolorbox{objectifs}{popbase,colframe=popink,colback=poporangeL,
  before upper={\popboxhead{Chapter objectives}{popink}{}}}
\newtcolorbox{prerequis}{popbase,colframe=popink,colback=popblueL,
  before upper={\popboxhead{Prerequisites}{popblue}{}}}
\newtcolorbox{bilan}{popbase,colframe=popink,colback=white,boxrule=2.8pt,
  before upper={\popboxhead{Fiche bilan}{poporange}{L'essentiel à connaître pour l'examen}}}
% Figure encadrée avec légende
\newtcolorbox{popfigure}[1][]{enhanced,breakable=false,colback=white,colframe=popline,boxrule=1.4pt,arc=8pt,
  left=10pt,right=10pt,top=10pt,bottom=8pt,before skip=10pt,after skip=12pt,halign=center,
  lower separated=false,halign lower=center,
  fontlower=\popsemi\fontsize{9}{11.5}\selectfont\color{popmuted},
  #1}
\newcommand{\legende}[1]{\par\vspace{6pt}{\popsemi\fontsize{9}{11.5}\selectfont\color{popmuted}#1}}

% ============================================================
% Signature OnBuch+
% ============================================================
\newcommand{\poplogoinline}[1]{{\poplogo\fontsize{#1}{#1}\selectfont\color{popink}OnBuch\color{poporange}+}}
\newcommand{\signatureonbuch}{%
  \begin{tcolorbox}[enhanced,colback=popink,colframe=popink,boxrule=2.2pt,arc=10pt,
      drop shadow=poporange,left=14pt,right=14pt,top=10pt,bottom=10pt,before skip=0pt,after skip=0pt]
    \tikz[baseline=-0.7ex]\node[circle,fill=popgreen,draw=popcream,line width=1.3pt,minimum size=22pt,inner sep=0]{\popcheck{white}};%
    \hspace{9pt}\begin{minipage}[c]{0.62\linewidth}
      {\popxbold\fontsize{12}{13}\selectfont\color{popcream}Signed\ \textbullet\ {\poplogo OnBuch\color{poporange}+}}\par\vspace{2pt}
      {\raggedright\popsemi\fontsize{8}{10.5}\selectfont\color{popline}\MakeUppercase{OnBuch+ teaching team}\\\MakeUppercase{Follows the official MINESEC syllabus}\par}
    \end{minipage}\hfill
    {\poplogo\fontsize{22}{22}\selectfont\color{popcream}OnBuch\color{poporange}+}
  \end{tcolorbox}%
}

% ============================================================
% Page de garde
% #1 titre · #2 matière · #3 niveau · #4 module · #5 n° leçon · #6 durée ·
% #7 description / objectifs (texte ou itemize)
% ============================================================
\newcommand{\pagegarde}[7]{%
  \thispagestyle{empty}
  \newgeometry{a4paper,margin=1.7cm,top=1.6cm,bottom=1.4cm}%
  \begin{tikzpicture}[remember picture,overlay]
    \fill[poporange!18] ([xshift=-3.2cm,yshift=-3.6cm]current page.north east) circle (4.6cm);
    \fill[poppurple!14] ([xshift=2.4cm,yshift=7.6cm]current page.south west) circle (3.8cm);
  \end{tikzpicture}%
  {\poplogo\fontsize{30}{30}\selectfont\color{popink}OnBuch\color{poporange}+}\hfill
  \raisebox{0.6ex}{\poppill{Signed course \textbullet\ Premium}{popink}{popcream}}\par
  \vspace{1pt}\popsquiggle{poppurple}\par
  \vspace{8pt}
  \begin{tcolorbox}[enhanced,colback=poporange,colframe=popink,boxrule=2.8pt,arc=13pt,
      drop shadow=popink,left=18pt,right=18pt,top=12pt,bottom=13pt,after skip=10pt]
    \poppill{#2 \textbullet\ #3}{popink}{popcream}\hspace{5pt}\poppill{#4}{white}{popink}\par\vspace{8pt}
    {\popdisplay\fontsize{14}{14}\selectfont\color{popink}CHAPTER #5}\par\vspace{4pt}
    {\raggedright\hyphenpenalty=10000\exhyphenpenalty=10000\popdisplay\fontsize{25}{27}\selectfont\color{popcream}\MakeUppercase{#1}\par}
    \vspace{8pt}
    {\popxbold\fontsize{9.5}{9.5}\selectfont\color{popink}\MakeUppercase{Suggested time: #6}}
  \end{tcolorbox}
  \begin{tcolorbox}[enhanced,colback=white,colframe=popink,boxrule=2.2pt,arc=10pt,
      drop shadow=popink,left=16pt,right=16pt,top=10pt,bottom=10pt,after skip=8pt]
    \poppill{In this chapter}{poppurple}{white}\par\vspace{5pt}
    \popsemi\fontsize{9.3}{12.6}\selectfont\color{popink}#7
  \end{tcolorbox}
  \noindent\begin{minipage}[t]{0.487\linewidth}
    \begin{tcolorbox}[enhanced,colback=popgreen,colframe=popink,boxrule=2.2pt,arc=10pt,
        drop shadow=popink,left=11pt,right=11pt,top=10pt,bottom=10pt,height=3.1cm,valign=center]
      \tcbox[colback=white,colframe=popink,boxrule=1.3pt,arc=5pt,boxsep=0pt,nobeforeafter,
        left=3pt,right=3pt,top=3pt,bottom=3pt,valign=center]{\qrcode[height=1.9cm]{https://play.google.com/store/apps/details?id=com.luvvix.onbuch&hl=en}}%
      \hspace{8pt}\begin{minipage}[c]{0.5\linewidth}
        {\popdisplay\fontsize{11}{12.5}\selectfont\color{white}\MakeUppercase{Download OnBuch}}\par\vspace{4pt}
        {\raggedright\popsemi\fontsize{8.4}{11}\selectfont\color{white}Free \textbullet\ Google Play\\AI tutor, past papers, results\par}
      \end{minipage}
    \end{tcolorbox}
  \end{minipage}\hfill
  \begin{minipage}[t]{0.487\linewidth}
    \begin{tcolorbox}[enhanced,colback=poppurple,colframe=popink,boxrule=2.2pt,arc=10pt,
        drop shadow=popink,left=11pt,right=11pt,top=10pt,bottom=10pt,height=3.1cm,valign=center]
      \tikz[baseline=-0.7ex]\node[circle,fill=white,draw=popink,line width=1.3pt,minimum size=1.6cm,inner sep=0]{\popdisplay\fontsize{24}{24}\selectfont\color{poppurple}+};%
      \hspace{8pt}\begin{minipage}[c]{0.6\linewidth}
        {\popdisplay\fontsize{11}{12.5}\selectfont\color{white}\MakeUppercase{Go OnBuch+}}\par\vspace{4pt}
        {\raggedright\popsemi\fontsize{8.4}{11}\selectfont\color{white}All signed courses, unlimited AI tutor, PDF sheets — OnBuch+ tab in the app\par}
      \end{minipage}
    \end{tcolorbox}
  \end{minipage}
  \vspace{8pt}
  \popcontenu
  \vfill
  \signatureonbuch
  \restoregeometry
  \newpage
}

% Rangée « Ce cours contient » (page de garde)
\newcommand{\poptile}[3]{% #1 couleur #2 gros texte #3 légende
  \begin{tcolorbox}[enhanced,colback=#1,colframe=popink,boxrule=2pt,arc=9pt,shadow={2.5pt}{-2.5pt}{0pt}{popink},
      left=6pt,right=6pt,top=9pt,bottom=9pt,halign=center,valign=center,height=1.75cm,width=\linewidth,nobeforeafter]
    {\popdisplay\fontsize{12.5}{13}\selectfont\color{white}\MakeUppercase{#2}}\par\vspace{3pt}
    {\centering\popsemi\fontsize{7.8}{9.5}\selectfont\color{white}#3\par}
  \end{tcolorbox}}
\newcommand{\popcontenu}{%
  \par\noindent{\popxboldls\fontsize{7.5}{7.5}\selectfont\color{popmuted}THIS SIGNED COURSE CONTAINS}\par\vspace{7pt}
  \noindent\begin{minipage}[t]{0.235\linewidth}\poptile{popblue}{Course}{complete, illustrated, step by step}\end{minipage}\hfill
  \begin{minipage}[t]{0.235\linewidth}\poptile{poporange}{Methods}{and worked examples}\end{minipage}\hfill
  \begin{minipage}[t]{0.235\linewidth}\poptile{poppink}{Exercises}{exam style + solutions}\end{minipage}\hfill
  \begin{minipage}[t]{0.235\linewidth}\poptile{popgreen}{Summary sheet}{the essentials to remember}\end{minipage}\par}

% Sommaire
\newtcolorbox{sommaire}{enhanced,colback=white,colframe=popink,boxrule=2.2pt,arc=10pt,
  drop shadow=popink,left=18pt,right=18pt,top=13pt,bottom=13pt,before skip=6pt,after skip=20pt,
  before upper={\poppill{Contents}{poppurple}{white}\par\vspace{9pt}\popsemi\fontsize{10.6}{14.5}\selectfont\color{popink}}}

% Signature de fin de cours — poussée tout en bas de la dernière page
\newcommand{\findecours}{%
  \par\vfill\nopagebreak
  \begin{center}\popsquiggle{poporange}\end{center}\vspace{4pt}
  \signatureonbuch
}
\AtBeginDocument{\def\llbracket{\mathopen{[\mkern-3.5mu[}}\def\rrbracket{\mathclose{]\mkern-3.5mu]}}} % intervalles d'entiers (glyphe ⟦ absent de la police)
% Glyphes absents de Fira Math : redéfinis à partir de glyphes disponibles
\AtBeginDocument{%
  \def\setminus{\mathbin{\backslash}}\let\smallsetminus\setminus
  \def\vdots{\mathord{\vbox{\baselineskip3.2pt\lineskiplimit0pt\kern4pt\hbox{.}\hbox{.}\hbox{.}}}}%
  \def\ddots{\mathinner{\mkern1mu\raise6.5pt\hbox{.}\mkern2mu\raise3.6pt\hbox{.}\mkern2mu\raise0.7pt\hbox{.}\mkern1mu}}%
  \def\triangle{\mathord{\scalebox{0.9}{$\symup{\Delta}$}}}%
  \def\nabla{\mathord{\raisebox{\depth}{\scalebox{1}[-1]{$\symup{\Delta}$}}}}%
  \def\checkmark{\popcheck{popdark}}%
  \def\star{\mathbin{\ast}}\def\bigstar{\mathord{\ast}}%
  \def\diamond{\mathbin{\scalebox{0.6}{$\Diamond$}}}%
  \def\bigcup{\mathop{\mathchoice{\raisebox{-0.3em}{\scalebox{1.7}{$\cup$}}}{\scalebox{1.25}{$\cup$}}{\cup}{\cup}}\displaylimits}%
  \def\bigcap{\mathop{\mathchoice{\raisebox{-0.3em}{\scalebox{1.7}{$\cap$}}}{\scalebox{1.25}{$\cap$}}{\cap}{\cap}}\displaylimits}%
  \def\lesseqqgtr{\mathrel{\lessgtr}}\def\bigoplus{\mathop{\oplus}\displaylimits}%
}
\providecommand{\ssi}{if and only if} % abréviation fréquente
\AtBeginDocument{\providecommand{\wideparen}{\overparen}} % arc de cercle

% Fonctions usuelles que les modèles utilisent sans que amsmath les fournisse
\providecommand{\arsinh}{\operatorname{arsinh}}\providecommand{\arcosh}{\operatorname{arcosh}}
\providecommand{\artanh}{\operatorname{artanh}}\providecommand{\arsech}{\operatorname{arsech}}
\providecommand{\arcsch}{\operatorname{arcsch}}\providecommand{\arcoth}{\operatorname{arcoth}}
\providecommand{\arcsinh}{\operatorname{arsinh}}\providecommand{\arccosh}{\operatorname{arcosh}}
\providecommand{\arctanh}{\operatorname{artanh}}\providecommand{\sech}{\operatorname{sech}}
\providecommand{\csch}{\operatorname{csch}}\providecommand{\arcsec}{\operatorname{arcsec}}
\providecommand{\arccsc}{\operatorname{arccsc}}\providecommand{\arccot}{\operatorname{arccot}}
\providecommand{\sgn}{\operatorname{sgn}}\providecommand{\lcm}{\operatorname{lcm}}
\providecommand{\Var}{\operatorname{Var}}\providecommand{\Cov}{\operatorname{Cov}}
\providecommand{\permil}{\ensuremath{{}^{0}\!/_{00}}}\providecommand{\textperthousand}{\ensuremath{{}^{0}\!/_{00}}}

%% FIGFIX-BEGIN v5 — correctifs globaux des figures (docs/onbuch-generation/tools/figfix)
\usepackage{adjustbox}\usepackage{etoolbox}\tcbuselibrary{raster}
% toute figure TikZ/pgfplots plus large que la ligne est réduite à la largeur disponible
\BeforeBeginEnvironment{tikzpicture}{\begin{adjustbox}{max width=\linewidth}}
\AfterEndEnvironment{tikzpicture}{\end{adjustbox}}
% légendes pgfplots lisibles par-dessus les courbes
\pgfplotsset{every axis/.append style={legend style={fill=white,fill opacity=0.92,text opacity=1,draw=black!15,inner sep=3pt}}}
% étiquettes d'arêtes (syntaxe edge["..."]) avec fond blanc
\tikzset{every edge quotes/.append style={fill=white,inner sep=1.2pt}}
%% FIGFIX-END
\begin{document}

\pagegarde{Money and banking(continue)}{Economics}{Upper Sixth}{Upper Sixth — Economics}{4}{7 periods}{This chapter completes the study of money and banking by examining how commercial banks operate and create credit, how the central bank (BEAC) regulates the monetary system, and how microfinance institutions extend financial services to low-income Cameroonians. It combines theoretical lessons with practical case studies rooted in the Cameroonian and CEMAC context, preparing students for data-response and essay questions at the GCE Advanced Level.
\vspace{4pt}
\begin{multicols}{2}\small
\begin{itemize}
\item By the end of this chapter I can draw and interpret the balance sheet of a commercial bank, distinguishing assets from liabilities.
\item By the end of this chapter I can explain the mechanism of credit creation by commercial banks and compute the credit multiplier.
\item By the end of this chapter I can identify and discuss the problems faced by commercial banks in Cameroon.
\item By the end of this chapter I can define a central bank and state its functions and instruments of monetary policy.
\item By the end of this chapter I can analyse the role of BEAC as the central bank of the CEMAC zone, including Cameroon.
\item By the end of this chapter I can describe microfinance activities and their importance for financial inclusion in Cameroon.
\end{itemize}\end{multicols}}

\begin{sommaire}
\begin{multicols}{2}
\begin{enumerate}
\item Section 1: The Balance Sheet of a Commercial Bank (TU 29 — Lesson 91)
\item Section 2: Credit Creation by Commercial Banks and Their Problems in Cameroon (TU 30 — Practical Work 20; TU 31 — Further Studies 5)
\item Section 3: The Central Bank and the Case of BEAC (TU 32 — Lesson 92; TU 33 — Guided Work 5)
\item Section 4: Microfinance Activities in Cameroon (TU 34 — Lesson 93)
\item Section 5: Appraisal of a Loan File by a Credit Union (TU 35 — Guided Work 6) and Chapter Synthesis
\item Integration activity
\item Methods and worked examples
\item Exercises
\item Solutions to the exercises
\item Summary sheet
\end{enumerate}
\end{multicols}
\end{sommaire}

\begin{objectifs}
\begin{itemize}
\item By the end of this chapter I can draw and interpret the balance sheet of a commercial bank, distinguishing assets from liabilities.
\item By the end of this chapter I can explain the mechanism of credit creation by commercial banks and compute the credit multiplier.
\item By the end of this chapter I can identify and discuss the problems faced by commercial banks in Cameroon.
\item By the end of this chapter I can define a central bank and state its functions and instruments of monetary policy.
\item By the end of this chapter I can analyse the role of BEAC as the central bank of the CEMAC zone, including Cameroon.
\item By the end of this chapter I can describe microfinance activities and their importance for financial inclusion in Cameroon.
\item By the end of this chapter I can outline the elements of a good appraisal of a loan file by a credit union.
\item By the end of this chapter I can solve GCE-style numerical and essay questions on banking and monetary policy.
\end{itemize}
\end{objectifs}

\begin{prerequis}
\begin{itemize}
\item Definition, functions and qualities of money (Lower Sixth / ECO03)
\item Types of money: fiduciary and scriptural (bank) money
\item The demand for and supply of money; quantity theory basics
\item Inflation: meaning, causes and measurement in Cameroon
\item Basic arithmetic with percentages and ratios
\end{itemize}
\end{prerequis}

\newpage

\coursec{Section 1: The Balance Sheet of a Commercial Bank (TU 29 — Lesson 91)}

Aminatou, a trader at Mokolo Market in Yaoundé, deposits her daily sales of $150\,000$ FCFA at her commercial bank every evening before going home. Two weeks later, her neighbour obtains a loan of $1\,000\,000$ FCFA from the very same bank to buy a taxi. Where did the bank find that money? Meanwhile, the BEAC in Yaoundé announces a rise in its key interest rate, and suddenly loans become more expensive across the whole CEMAC zone, while the credit union in her quarter refuses a loan file that seemed convincing. Behind these everyday scenes lie the \cle{balance sheet} of commercial banks, the mysterious power of credit creation, the supervisory role of the central bank and the growing world of microfinance. This chapter opens the doors of the banking system so that you can explain, calculate and evaluate what really happens with money in Cameroon.

Our first problem is simple to state but fundamental: \emph{how can we describe, at a glance, everything a bank owns and everything it owes?} The answer is the balance sheet.

\courssub{Recall: What Is a Commercial Bank?}

Before opening the accounts, let us recall what a commercial bank actually does. When Aminatou hands her $150\,000$ FCFA to the cashier, the bank is performing its first function: \cle{deposit-taking}. When her neighbour drives away in his new taxi, the bank has performed its second function: \cle{lending}. And when Aminatou pays her wholesaler by cheque or by mobile transfer instead of carrying cash across town, the bank is performing its third function: providing \cle{payment services}.

\begin{definition}[Commercial bank]
A \cle{commercial bank} is a financial institution whose main activities are (i) collecting deposits from the public, (ii) granting loans and overdrafts to households and firms, and (iii) providing means of payment (cheques, transfers, bank cards, mobile banking). Its objective is to make a profit from the difference between the interest it charges on loans and the interest it pays on deposits — this difference is called the \cle{interest margin} or \emph{spread}.
\end{definition}

In Cameroon, well-known commercial banks include \textbf{Afriland First Bank}, \textbf{Société Générale Cameroun (SGBC)}, \textbf{BICEC} (Banque Internationale du Cameroun pour l'Épargne et le Crédit) and \textbf{Ecobank Cameroun}. All of them are licensed and supervised by the regional central bank, the \textbf{BEAC} (Bank of Central African States), which we shall study in Section 3.

\courssub{The Meaning and Structure of a Balance Sheet}

Imagine you freeze time at midnight on 31 December and list, on one sheet of paper, everything the bank \emph{owns} and everything the bank \emph{owes}. That sheet is the balance sheet.

\begin{definition}[Balance sheet]
A \cle{balance sheet} is a financial statement that shows, at a given date, the \cle{assets} of a bank (what it owns or is owed), its \cle{liabilities} (what it owes to others) and its \cle{net worth} (the residual belonging to shareholders).
\end{definition}

\begin{definition}[Assets and liabilities]
\begin{itemize}
\item An \cle{asset} is anything the bank \textbf{owns} or anything \textbf{owed to the bank} (cash in the vault, loans granted to customers, buildings, government securities).
\item A \cle{liability} is anything the bank \textbf{owes to others} (customers' deposits, money borrowed from the central bank or from other banks).
\item The \cle{net worth} (or equity, or own funds) is what would remain for the shareholders if all liabilities were repaid: share capital plus accumulated reserves and retained profits.
\end{itemize}
\end{definition}

\begin{propriete}[The accounting identity]
In every balance sheet, by construction:
\[ \text{Total Assets} \;=\; \text{Total Liabilities} \;+\; \text{Net Worth}. \]
This identity always holds because net worth is \emph{defined} as the residual: $\text{Net Worth} = \text{Assets} - \text{Liabilities}$. A balance sheet that does not balance contains an error. (The identity itself is definitional and does not require a proof in the examination, but you must be able to \emph{use} it, for example to find a missing net worth figure.)
\end{propriete}

\begin{methode}[Classifying any balance sheet item]
To decide where an item belongs, ask yourself one question:
\begin{enumerate}
\item \textbf{Does the bank own it, or is it owed to the bank?} $\rightarrow$ it is an \textbf{asset} (e.g. a loan granted to a customer: the customer owes the bank).
\item \textbf{Does the bank owe it to someone else?} $\rightarrow$ it is a \textbf{liability} (e.g. a customer's deposit: the bank owes the customer).
\item \textbf{Is it the shareholders' own contribution or accumulated profit?} $\rightarrow$ it is part of \textbf{net worth}.
\end{enumerate}
\end{methode}

\begin{attention}[Deposits are NOT the bank's property]
The most common mistake at the GCE is to treat customers' deposits as something the bank \emph{owns}. Wrong! When Aminatou deposits $150\,000$ FCFA, the bank must repay her \textbf{on demand}. The deposit is therefore a \textbf{liability} of the bank. Symmetrically, a loan granted to her neighbour is an \textbf{asset} of the bank, because the neighbour owes the money to the bank. Always reason \emph{from the bank's point of view}.
\end{attention}

\courssub{The Liabilities Side: Where Does the Bank's Money Come From?}

The liabilities side answers the question: \emph{who supplied the funds the bank is using?} We distinguish four main sources.

\begin{itemize}
\item \textbf{Sight deposits (demand deposits).} Money placed on current accounts, withdrawable at any time by cheque, card or transfer. This is typically the largest liability of a commercial bank. Aminatou's deposit belongs here.
\item \textbf{Time and savings deposits.} Money deposited for a fixed period or on savings accounts (e.g. a \emph{compte à terme} at SGBC). The bank pays interest on them but can count on them more durably.
\item \textbf{Borrowed funds.} Banks themselves borrow: either from the \textbf{central bank} (refinancing or rediscounting) or from other banks on the \textbf{interbank market}.
\item \textbf{Shareholders' equity and reserves (net worth).} The capital contributed by the owners plus retained profits. This is the bank's safety cushion: it absorbs losses before depositors are affected.
\end{itemize}

\courssub{The Assets Side: How Does the Bank Use the Funds?}

The assets side answers the question: \emph{what has the bank done with the funds?} We move from the most liquid to the least liquid (and generally most profitable) uses.

\begin{itemize}
\item \textbf{Cash in vault and reserves at the central bank.} Notes and coins in the tills, plus the account the bank holds at the BEAC. These earn little or nothing but allow the bank to meet withdrawals immediately. Part of these reserves is \emph{compulsory} (required reserves imposed by the central bank).
\item \textbf{Loans and overdrafts to customers.} Credits granted to households (the neighbour's taxi loan) and to firms (a loan to a cocoa exporter in Douala). This is usually the most profitable asset, but also the riskiest and least liquid: a loan cannot be recalled instantly.
\item \textbf{Investment securities.} Treasury bills and government bonds bought by the bank. They earn interest and can be sold if the bank needs cash — a compromise between liquidity and profitability.
\item \textbf{Fixed assets.} Buildings, branches, computers, vehicles. Necessary for operations but totally illiquid and not directly profitable.
\end{itemize}

\begin{popfigure}
\begin{center}
\begin{tikzpicture}[scale=1.0]
% T-account frame
\draw[thick, popink] (-7.4,0) rectangle (7.4,-6.4);
\draw[thick, popink] (0,0) -- (0,-6.4);
\draw[thick, popink] (-7.4,-0.9) -- (7.4,-0.9);
\node[align=center, font=\bfseries, text=popblue] at (-3.7,-0.45) {ASSETS (uses of funds)};
\node[align=center, font=\bfseries, text=poporange] at (3.7,-0.45) {LIABILITIES \& NET WORTH (sources)};
% Assets items
\node[anchor=west, align=left, text width=6.6cm] at (-7.2,-1.5) {Cash in vault and reserves at BEAC \dotfill 800};
\node[anchor=west, align=left, text width=6.6cm] at (-7.2,-2.3) {Loans and overdrafts to customers \dotfill 6\,200};
\node[anchor=west, align=left, text width=6.6cm] at (-7.2,-3.1) {Investment securities (T-bills, bonds) \dotfill 1\,500};
\node[anchor=west, align=left, text width=6.6cm] at (-7.2,-3.9) {Fixed assets (buildings, equipment) \dotfill 500};
\draw[popline] (-7.2,-4.4) -- (-0.2,-4.4);
\node[anchor=west, font=\bfseries] at (-7.2,-4.9) {Total assets \dotfill 9\,000};
% Liabilities items
\node[anchor=west, align=left, text width=6.6cm] at (0.2,-1.5) {Sight deposits \dotfill 5\,000};
\node[anchor=west, align=left, text width=6.6cm] at (0.2,-2.3) {Time and savings deposits \dotfill 2\,500};
\node[anchor=west, align=left, text width=6.6cm] at (0.2,-3.1) {Borrowed funds (BEAC, interbank) \dotfill 500};
\node[anchor=west, align=left, text width=6.6cm] at (0.2,-3.9) {Share capital and reserves \dotfill 1\,000};
\draw[popline] (0.2,-4.4) -- (7.2,-4.4);
\node[anchor=west, font=\bfseries] at (0.2,-4.9) {Total liabilities + net worth \dotfill 9\,000};
\node[align=center, font=\itshape, text=popmuted] at (0,-5.9) {Figures in millions of FCFA, simplified example};
\end{tikzpicture}
\end{center}
\legende{Simplified balance sheet (T-account) of a Cameroonian commercial bank. Assets on the left equal liabilities plus net worth on the right: $9\,000 = 9\,000$ million FCFA.}
\end{popfigure}

\courssub{The Fundamental Trade-off: Liquidity, Profitability, Security}

Why doesn't the bank lend out \emph{all} of Aminatou's deposit and every other deposit, since loans earn the highest interest? Because depositors can come and withdraw their money \emph{at any time}. If the bank has no cash left, it fails. Bank asset management is therefore a permanent balancing act between three objectives that pull in opposite directions.

\begin{definition}[The three objectives of bank asset management]
\begin{itemize}
\item \cle{Liquidity}: the ability to meet withdrawal demands immediately. Cash is perfectly liquid but earns nothing.
\item \cle{Profitability}: the search for the highest return. Long-term loans are the most profitable but the least liquid.
\item \cle{Security} (solvency): the avoidance of losses. Risky loans pay more but may never be repaid; safe assets (treasury bills) pay less.
\end{itemize}
\end{definition}

\begin{propriete}[The liquidity--profitability--security trade-off]
For a bank, the more liquid and safe an asset is, the less profitable it tends to be; the more profitable an asset is, the less liquid and the riskier it tends to be. The bank must therefore \textbf{diversify} its assets: keep enough cash and reserves for daily withdrawals, hold some easily sellable securities, and lend the rest at profitable rates to reliable borrowers.
\end{propriete}

\begin{popfigure}
\begin{center}
\begin{tikzpicture}[scale=1.0]
\coordinate (A) at (90:4.2);
\coordinate (B) at (210:4.2);
\coordinate (C) at (330:4.2);
\draw[very thick, popblue] (A) -- (B) -- (C) -- cycle;
\node[align=center, font=\bfseries, text=popblue] at (90:5.1) {LIQUIDITY};
\node[align=center, font=\bfseries, text=poporange] at (210:5.1) {PROFITABILITY};
\node[align=center, font=\bfseries, text=popgreen] at (330:5.1) {SECURITY};
\node[align=center, text width=3.2cm, font=\small] at (90:2.6) {Cash, reserves at BEAC};
\node[align=center, text width=3.4cm, font=\small] at (210:2.6) {Loans, overdrafts};
\node[align=center, text width=3.2cm, font=\small] at (330:2.6) {Treasury bills, bonds};
\node[align=center, font=\itshape, text width=3.4cm] at (0,0) {The bank must balance all three};
\end{tikzpicture}
\end{center}
\legende{The liquidity--profitability--security triangle: moving towards one vertex means moving away from the other two.}
\end{popfigure}

\courssub{Key Ratios Derived from the Balance Sheet}

The balance sheet is not just a photograph; it is a diagnostic tool. Three ratios are essential at the GCE level.

\begin{definition}[Key balance sheet ratios]
\begin{itemize}
\item \textbf{Reserve ratio} $= \dfrac{\text{Reserves (cash + deposits at central bank)}}{\text{Total deposits}}$. It measures the fraction of deposits kept liquid.
\item \textbf{Liquidity ratio} $= \dfrac{\text{Liquid assets (cash, reserves, short-term securities)}}{\text{Short-term liabilities (sight deposits)}}$. It measures the ability to face withdrawals.
\item \textbf{Loan-to-deposit ratio} $= \dfrac{\text{Total loans}}{\text{Total deposits}}$. A high value means aggressive lending; a value above $100\%$ means the bank lends more than its deposits and must borrow elsewhere.
\end{itemize}
\end{definition}

\begin{exoresolu}[Reading the balance sheet of ``Mokolo Bank Ltd'']
\textbf{Statement.} Using the simplified balance sheet of Mokolo Bank Ltd (figures in millions of FCFA): cash and reserves at BEAC $= 800$; loans $= 6\,200$; securities $= 1\,500$; fixed assets $= 500$; sight deposits $= 5\,000$; time deposits $= 2\,500$; borrowed funds $= 500$; capital and reserves $= 1\,000$.
\begin{enumerate}
\item Verify the accounting identity.
\item Compute the reserve ratio, the liquidity ratio and the loan-to-deposit ratio.
\item Comment on the bank's policy.
\end{enumerate}
\tcblower
\textbf{Solution.}
\begin{enumerate}
\item Total assets $= 800 + 6\,200 + 1\,500 + 500 = 9\,000$. Total liabilities $+$ net worth $= 5\,000 + 2\,500 + 500 + 1\,000 = 9\,000$. The identity holds: $9\,000 = 9\,000$.
\item Total deposits $= 5\,000 + 2\,500 = 7\,500$.
\begin{align*}
\text{Reserve ratio} &= \frac{800}{7\,500} \approx 10.7\% \\
\text{Liquidity ratio} &= \frac{800 + 1\,500}{5\,000} = \frac{2\,300}{5\,000} = 46\% \\
\text{Loan-to-deposit ratio} &= \frac{6\,200}{7\,500} \approx 82.7\%
\end{align*}
\item The bank keeps about $10.7\%$ of deposits as reserves and holds comfortable liquid assets against sight deposits ($46\%$). Its loan-to-deposit ratio of $82.7\%$ shows an active but not reckless lending policy: most deposits are transformed into profitable loans, while liquidity remains adequate. The bank respects the liquidity--profitability--security balance.
\end{enumerate}
\end{exoresolu}

\begin{attention}[Exam traps on ratios]
\begin{itemize}
\item Do not confuse the \textbf{reserve ratio} (reserves over \emph{deposits}) with the \textbf{capital ratio} (net worth over \emph{assets}).
\item In the liquidity ratio, only \emph{short-term} liabilities (sight deposits) appear in the denominator, not time deposits.
\item Always express ratios as percentages and \textbf{interpret} them: a bare number without commentary earns only part of the marks.
\end{itemize}
\end{attention}

\begin{exercice}[Constructing a balance sheet]
Douala Star Bank presents the following items at 31 December (millions of FCFA): loans to customers $4\,000$; sight deposits $3\,500$; cash and reserves at BEAC $600$; treasury bills $900$; time deposits $1\,800$; buildings $300$; borrowing from BEAC $200$; share capital and reserves: to be determined.
\begin{enumerate}
\item Classify each item as asset, liability or net worth, and draw up the balance sheet.
\item Find the value of share capital and reserves using the accounting identity.
\item Compute the reserve ratio and the loan-to-deposit ratio, and comment briefly.
\end{enumerate}
\end{exercice}

\begin{corrige}[Exercise 1]
\begin{enumerate}
\item Assets: loans $4\,000$; cash and reserves $600$; treasury bills $900$; buildings $300$; total $= 5\,800$. Liabilities: sight deposits $3\,500$; time deposits $1\,800$; borrowing from BEAC $200$; total $= 5\,500$. Net worth: unknown $E$.
\item $E = 5\,800 - 5\,500 = 300$ million FCFA.
\item Reserve ratio $= 600 / (3\,500 + 1\,800) = 600 / 5\,300 \approx 11.3\%$. Loan-to-deposit ratio $= 4\,000 / 5\,300 \approx 75.5\%$. The bank keeps a prudent liquidity cushion and lends about three quarters of its deposits: a balanced, moderately profitable policy.
\end{enumerate}
\end{corrige}

\begin{aretenir}[Key points of Section 1]
\begin{itemize}
\item A commercial bank takes deposits, grants loans and provides payment services (Afriland, SGBC, BICEC, Ecobank in Cameroon).
\item The \cle{balance sheet} records assets (what the bank owns or is owed) and liabilities (what it owes) at a given date, with $\text{Assets} = \text{Liabilities} + \text{Net Worth}$.
\item Deposits are \textbf{liabilities}; loans are \textbf{assets} — always reason from the bank's point of view.
\item Bank management is a trade-off between \cle{liquidity}, \cle{profitability} and \cle{security}.
\item Master the reserve ratio, liquidity ratio and loan-to-deposit ratio, and always interpret your results.
\end{itemize}
\end{aretenir}

\begin{savaistu}[GCE exam tip]
At the GCE Advanced Level, examiners frequently give you a jumbled list of balance sheet items and ask you to (i) construct the balance sheet, (ii) compute ratios and (iii) comment. Train yourself to classify each item in seconds with the question ``does the bank own it or owe it?'', and never forget that the two totals must be equal — if they are not, the missing figure is the net worth, found by subtraction.
\end{savaistu}

Now that we can read a bank's accounts, a deeper mystery appears: if Mokolo Bank lends out deposits, and the taxi driver pays a car dealer who deposits the money again\ldots\ can banks actually \emph{create} money? That is the subject of Section 2.

\coursec{Section 2: Credit Creation by Commercial Banks and Their Problems in Cameroon (TU 30 — Practical Work 20; TU 31 — Further Studies 5)}

In Section 1 we studied the balance sheet of a commercial bank and saw that its main assets are reserves and loans, while its main liabilities are deposits. We also noticed something remarkable: a bank keeps only a \emph{fraction} of its deposits as reserves and lends out the rest. This simple fact has an extraordinary consequence. When a bank grants a loan, the money does not disappear from the economy — the borrower spends it, the seller deposits it in another bank, that bank lends most of it again, and so on. One initial deposit can therefore give birth to a whole chain of new deposits. This process, called \cle{credit creation}, is the subject of this section. We shall first explain how a single bank creates credit, then show how the whole banking system multiplies deposits, derive the \cle{credit multiplier}, discuss the limits of the process, and finish with a practical numerical exercise and a study of the problems faced by commercial banks in Cameroon.

\courssub{How a Single Bank Creates Credit}

Imagine a trader in Douala who deposits $1\,000\,000$ FCFA in cash at Bank A. If the central bank imposes a \cle{required reserve ratio} of $10\%$, Bank A must keep $100\,000$ FCFA as reserves but is free to lend the remaining $900\,000$ FCFA, called its \cle{excess reserves}. Suppose Bank A lends this $900\,000$ FCFA to a carpenter who uses it to buy timber. The timber seller receives the $900\,000$ FCFA and deposits it in her own bank, Bank B. Notice what has happened: the original $1\,000\,000$ FCFA deposit still exists (the trader can withdraw it), and a \emph{new} deposit of $900\,000$ FCFA has appeared at Bank B. Money has been \emph{created} — not printed, but created through lending.

\begin{definition}[Credit creation]
\cle{Credit creation} is the process by which commercial banks generate new deposits (and therefore new money) by granting loans out of their \cle{excess reserves}, that is, the part of deposits which exceeds the reserves they are legally or prudentially required to hold.
\end{definition}

The key intuition is that a single bank can only lend its \emph{excess} reserves — it cannot create money on its own beyond that, because the borrower will spend the loan and the bank will lose reserves to other banks. But the \emph{system as a whole} can create a multiple of the initial deposit, because one bank's loan becomes another bank's deposit.

\begin{attention}[A single bank versus the banking system]
A single bank in isolation can lend at most its excess reserves. It is the \emph{banking system as a whole} — Bank A, Bank B, Bank C, \dots — that multiplies deposits. Never write that ``one bank creates $1/r$ times the deposit''; that is the property of the system.
\end{attention}

\courssub{The Multiple Expansion of Deposits: A Step-by-Step Illustration}

Let us follow the chain with a required reserve ratio $r = 10\%$ and an initial new deposit of $1\,000\,000$ FCFA at Bank A. Each bank keeps $10\%$ of the deposit it receives and lends out $90\%$.

\textbf{Bank A.} Receives a deposit of $1\,000\,000$; keeps $100\,000$ as required reserves; lends $900\,000$.

\begin{center}
\begin{tabular}{|l|r|l|r|}
\hline
\rowcolor{popblueL} \multicolumn{2}{|c|}{\textbf{Assets (Bank A)}} & \multicolumn{2}{c|}{\textbf{Liabilities (Bank A)}} \\
\hline
Reserves & $100\,000$ & Deposits & $1\,000\,000$ \\
\hline
Loans & $900\,000$ & & \\
\hline
\end{tabular}
\end{center}

\textbf{Bank B.} The $900\,000$ loan is spent and redeposited at Bank B. Bank B keeps $90\,000$ ($10\%$) and lends $810\,000$.

\begin{center}
\begin{tabular}{|l|r|l|r|}
\hline
\rowcolor{popblueL} \multicolumn{2}{|c|}{\textbf{Assets (Bank B)}} & \multicolumn{2}{c|}{\textbf{Liabilities (Bank B)}} \\
\hline
Reserves & $90\,000$ & Deposits & $900\,000$ \\
\hline
Loans & $810\,000$ & & \\
\hline
\end{tabular}
\end{center}

\textbf{Bank C.} Receives $810\,000$; keeps $81\,000$; lends $729\,000$. The process continues through Bank D, Bank E, and so on, each deposit being $90\%$ of the previous one.

\begin{popfigure}
\begin{center}
\begin{tikzpicture}[node distance=0.55cm, every node/.style={font=\small},
box/.style={draw=popblue, thick, rounded corners, fill=popblueL, text width=4.6cm, align=center, inner sep=5pt},
arr/.style={->, thick, poporange}]
\node[box, align=center] (A){\textbf{Bank A}\\ Deposit: $1\,000\,000$\\ Reserves: $100\,000$\\ Loan: $900\,000$};
\node[box, right=1.1cm of A, align=center] (B){\textbf{Bank B}\\ Deposit: $900\,000$\\ Reserves: $90\,000$\\ Loan: $810\,000$};
\node[box, right=1.1cm of B, align=center] (C){\textbf{Bank C}\\ Deposit: $810\,000$\\ Reserves: $81\,000$\\ Loan: $729\,000$};
\node[right=0.9cm of C, font=\Large, popdark] (D) {$\cdots$};
\draw[arr] (A) -- node[above, font=\scriptsize]{loan spent \& redeposited} (B);
\draw[arr] (B) -- node[above, font=\scriptsize]{loan spent \& redeposited} (C);
\draw[arr] (C) -- (D);
\end{tikzpicture}
\end{center}
\legende{The credit creation chain with a $10\%$ reserve ratio: each bank keeps $10\%$ and lends $90\%$, so deposits shrink at each round while their sum grows.}
\end{popfigure}

The deposits form a geometric progression:
\[
1\,000\,000 + 900\,000 + 810\,000 + 729\,000 + \dots
= 1\,000\,000 \times \left(1 + 0.9 + 0.9^2 + 0.9^3 + \dots\right).
\]
The sum of an infinite geometric series with first term $1$ and common ratio $0.9$ is $\dfrac{1}{1-0.9} = \dfrac{1}{0.1} = 10$. Hence total deposits $= 1\,000\,000 \times 10 = 10\,000\,000$ FCFA.

\courssub{The Credit (Deposit) Multiplier}

\begin{propriete}[The credit multiplier — formula]
If the required reserve ratio is $r$ (expressed as a decimal) and banks hold no excess reserves while the public holds no extra cash, then the \cle{credit multiplier} is
\[
M = \frac{1}{r}, \qquad
\text{Total deposits} = \text{initial new deposit} \times \frac{1}{r}, \qquad
\text{New credit created} = \text{Total deposits} - \text{initial deposit}.
\]
\emph{Proof sketch (instructive, not required at GCE):} total deposits $= D_0\left(1 + (1-r) + (1-r)^2 + \dots\right) = D_0 \times \dfrac{1}{1-(1-r)} = \dfrac{D_0}{r}$.
\end{propriete}

\begin{exemplebox}[Numerical applications in FCFA]
With an initial new deposit of $1\,000\,000$ FCFA:
\begin{itemize}
\item If $r = 10\%$: $M = \frac{1}{0.1} = 10$; total deposits $= 10\,000\,000$ FCFA; new credit created $= 9\,000\,000$ FCFA.
\item If $r = 20\%$: $M = \frac{1}{0.2} = 5$; total deposits $= 5\,000\,000$ FCFA; new credit created $= 4\,000\,000$ FCFA.
\item If $r = 25\%$: $M = 4$; total deposits $= 4\,000\,000$ FCFA; new credit created $= 3\,000\,000$ FCFA.
\end{itemize}
The higher the reserve ratio, the smaller the multiplier: the central bank can therefore use $r$ as an instrument of monetary policy.
\end{exemplebox}

\begin{popfigure}
\begin{center}
\begin{tikzpicture}
\begin{axis}[
ybar, bar width=0.5cm,
width=11cm, height=6.5cm,
xlabel={Round of the process}, ylabel={Cumulative deposits (millions FCFA)},
symbolic x coords={A, B, C, D, E, Limit},
xtick=data,
ymin=0, ymax=11,
legend style={at={(0.97,0.97)}, anchor=north east, font=\small},
nodes near coords, every node near coord/.append style={font=\scriptsize},
]
\addplot[fill=popblue] coordinates {(A,1) (B,1.9) (C,2.71) (D,3.439) (E,4.0951) (Limit,10)};
\addplot[fill=poporange] coordinates {(A,1) (B,1.8) (C,2.44) (D,2.952) (E,3.3616) (Limit,5)};
\legend{$r = 10\%$ (limit $= 10$), $r = 20\%$ (limit $= 5$)}
\end{axis}
\end{tikzpicture}
\end{center}
\legende{Cumulative deposits created from an initial deposit of $1\,000\,000$ FCFA. With $r = 10\%$ the total converges to $10$ million FCFA; with $r = 20\%$ it converges to only $5$ million FCFA.}
\end{popfigure}

\begin{methode}[Computing total credit created]
\begin{enumerate}
\item Find the multiplier: $M = \dfrac{1}{r}$, where $r$ is the reserve ratio as a decimal.
\item Multiply the initial \emph{new} deposit by the multiplier to obtain total deposits.
\item Subtract the initial deposit from total deposits to obtain the \emph{newly created credit}.
\end{enumerate}
\end{methode}

\begin{attention}[Total deposits versus newly created credit]
A very common GCE mistake is to give total deposits as the answer when the question asks for \emph{credit created}. The initial deposit is \emph{not} created money — it existed before as cash. Newly created credit $=$ total deposits $-$ initial deposit. In our example, credit created is $9\,000\,000$ FCFA, not $10\,000\,000$ FCFA.
\end{attention}

\begin{attention}[The multiplier is a maximum]
The formula $M = 1/r$ assumes no cash leakages and no excess reserves. In Cameroon these leakages are large (many people hold cash outside banks, and banks often keep extra reserves), so the \emph{real} credit creation is much smaller than the theoretical maximum.
\end{attention}

\courssub{Limits to Credit Creation}

In reality, the expansion of deposits is limited by several factors:
\begin{itemize}
\item \textbf{Cash leakages (withdrawals).} If part of every loan is held by the public as cash instead of being redeposited — very common in Cameroon where many transactions are in cash — the amount returning to the banking system shrinks, and the multiplier falls.
\item \textbf{Excess reserves.} Banks may prudently keep reserves above the legal minimum, especially when borrowers seem risky; idle excess reserves cannot create credit.
\item \textbf{Lack of creditworthy borrowers.} Credit creation requires demand for loans from borrowers with viable projects and acceptable guarantees. If few firms or households qualify, banks cannot lend even when they have excess liquidity.
\item \textbf{Central bank regulations.} The central bank can raise the reserve ratio, impose credit ceilings, or tighten refinancing conditions, directly limiting banks' ability to lend.
\end{itemize}

\courssub{Practical Work: Credit Creation in the Cameroonian Banking System}

\begin{exoresolu}[Computing total credit created]
\textbf{Statement.} A new deposit of $5\,000\,000$ FCFA is made in a commercial bank in Yaoundé. The reserve requirement fixed by the central bank is $10\%$. Assuming no cash leakage and no excess reserves: (a) calculate the credit multiplier; (b) calculate the total deposits the banking system can support; (c) calculate the total credit newly created; (d) show the first three rounds of the process.
\tcblower
\textbf{Solution.}
(a) Multiplier: $M = \dfrac{1}{r} = \dfrac{1}{0.1} = 10$.

(b) Total deposits $= 5\,000\,000 \times 10 = 50\,000\,000$ FCFA.

(c) Newly created credit $= 50\,000\,000 - 5\,000\,000 = 45\,000\,000$ FCFA.

(d) First three rounds:
\begin{center}
\begin{tabular}{|l|r|r|r|}
\hline
\rowcolor{popblueL} \textbf{Bank} & \textbf{Deposit (FCFA)} & \textbf{Reserves 10\% (FCFA)} & \textbf{Loan (FCFA)} \\
\hline
Bank A & $5\,000\,000$ & $500\,000$ & $4\,500\,000$ \\
\hline
Bank B & $4\,500\,000$ & $450\,000$ & $4\,050\,000$ \\
\hline
Bank C & $4\,050\,000$ & $405\,000$ & $3\,645\,000$ \\
\hline
\end{tabular}
\end{center}
Each bank lends $90\%$ of what it receives; the sum of all rounds converges to $50\,000\,000$ FCFA of deposits, of which $45\,000\,000$ FCFA is newly created credit.
\end{exoresolu}

\begin{methode}[Exam tip (GCE)]
In the examination, show \emph{each step} of the T-account chain (at least the first three banks), state the multiplier formula, and distinguish total deposits from credit created. Marks are awarded for the \emph{process}, not only for the final figure.
\end{methode}

\courssub{Problems Faced by Commercial Banks in Cameroon (Further Studies)}

Despite their theoretical power to create credit, commercial banks in Cameroon face serious practical difficulties that limit their role in financing the economy.

\begin{popfigure}
\begin{center}
\begin{center}\tcbox[colback=popblue,colframe=popblue,coltext=white,arc=9pt,boxsep=4pt,left=6pt,right=6pt]{\bfseries\small Problems of commercial banks in Cameroon}\end{center}
\vspace{-2pt}
\begin{tcbraster}[raster columns=3,raster equal height=rows,raster column skip=6pt,raster row skip=6pt,enhanced,blankest]
\begin{tcolorbox}[enhanced,colback=popblue,colframe=popblue,coltext=white,boxrule=0.9pt,arc=3pt,left=4pt,right=4pt,top=3pt,bottom=3pt,boxsep=1pt,halign=left]\bfseries\scriptsize High default rates \& loan recovery difficulties\end{tcolorbox}
\begin{tcolorbox}[enhanced,colback=poporange,colframe=poporange,coltext=white,boxrule=0.9pt,arc=3pt,left=4pt,right=4pt,top=3pt,bottom=3pt,boxsep=1pt,halign=left]\bfseries\scriptsize Dominance of the informal sector\end{tcolorbox}
\begin{tcolorbox}[enhanced,colback=popgreen,colframe=popgreen,coltext=white,boxrule=0.9pt,arc=3pt,left=4pt,right=4pt,top=3pt,bottom=3pt,boxsep=1pt,halign=left]\bfseries\scriptsize Low banking rate (bancarisation)\end{tcolorbox}
\begin{tcolorbox}[enhanced,colback=poppurple,colframe=poppurple,coltext=white,boxrule=0.9pt,arc=3pt,left=4pt,right=4pt,top=3pt,bottom=3pt,boxsep=1pt,halign=left]\bfseries\scriptsize Liquidity shortages \& high operating costs\end{tcolorbox}
\begin{tcolorbox}[enhanced,colback=popgold,colframe=popgold,coltext=white,boxrule=0.9pt,arc=3pt,left=4pt,right=4pt,top=3pt,bottom=3pt,boxsep=1pt,halign=left]\bfseries\scriptsize Competition from microfinance \& mobile money\end{tcolorbox}
\begin{tcolorbox}[enhanced,colback=poppink,colframe=poppink,coltext=white,boxrule=0.9pt,arc=3pt,left=4pt,right=4pt,top=3pt,bottom=3pt,boxsep=1pt,halign=left]\bfseries\scriptsize Government borrowing crowding out private credit; weak collateral \& land titles\end{tcolorbox}
\end{tcbraster}

\end{center}
\legende{Mind map of the main problems faced by commercial banks in Cameroon.}
\end{popfigure}

\textbf{High default rates and loan recovery difficulties.} Many borrowers fail to repay their loans, and recovering debts through the courts is slow and costly. Non-performing loans tie up banks' funds and make them reluctant to lend.

\textbf{Dominance of the informal sector.} A large share of economic activity — small traders, artisans, transporters — operates informally, without proper accounts, tax identification or registered businesses. Banks cannot assess the risk of such borrowers and therefore refuse them credit.

\textbf{Low banking rate (bancarisation).} A large proportion of the population has no bank account, keeping savings in cash, in \emph{tontines} or in microfinance institutions. This reduces the deposit base from which banks create credit.

\textbf{Liquidity shortages and high operating costs.} Banks sometimes face liquidity pressures, while the costs of branches, staff, security and energy are high relative to their business volume, pushing up interest rates on loans.

\textbf{Competition from microfinance and mobile money.} Credit unions, savings cooperatives and mobile money services (which allow transfers and payments by telephone) attract small savers and borrowers away from traditional banks.

\textbf{Government borrowing crowding out private credit.} When the government borrows heavily from banks (for example through treasury securities), banks prefer this safe lending and less credit remains for private businesses — the \cle{crowding-out} effect.

\textbf{Weak collateral and land-title systems.} Banks lend against guarantees, especially land and buildings. Because many properties lack registered land titles, potential borrowers cannot offer acceptable collateral, and banks cannot easily seize assets in case of default.

\textbf{Consequences for the economy.} These problems mean that credit creation in Cameroon is far below its theoretical potential: investment by small and medium-sized enterprises is starved of finance, interest rates remain high, economic growth and job creation are slowed, and monetary policy transmitted through banks is weakened.

\textbf{Possible remedies.} Among the solutions often proposed are: speeding up land-title registration to widen access to collateral; creating and strengthening credit bureaux to share information on borrowers; encouraging the formalisation of small businesses; developing judicial procedures for faster loan recovery; promoting financial education to raise the banking rate; and closer cooperation between banks, microfinance institutions and mobile money operators to reach the unbanked population.

\begin{aretenir}[Key points of Section 2]
\begin{itemize}
\item Credit creation: banks create new deposits by lending their excess reserves.
\item A single bank lends only its excess reserves; the \emph{system} multiplies deposits.
\item Credit multiplier: $M = \dfrac{1}{r}$; total deposits $=$ initial deposit $\times M$; new credit $=$ total deposits $-$ initial deposit.
\item The multiplier is reduced by cash leakages, excess reserves, weak loan demand and central bank regulation.
\item Cameroonian banks face defaults, informality, low bancarisation, liquidity and cost pressures, competition from microfinance and mobile money, crowding out, and weak collateral systems — all of which limit real credit creation.
\end{itemize}
\end{aretenir}

\begin{exercice}[Credit creation and the multiplier]
A new deposit of $2\,000\,000$ FCFA is made at a bank in Bamenda. The reserve ratio is $20\%$.
\begin{enumerate}
\item Calculate the credit multiplier.
\item Calculate the total deposits the banking system can support and the total credit newly created.
\item Present the T-accounts of the first two banks in the chain.
\item Explain two reasons why, in the real Cameroonian economy, the credit actually created would be smaller than your answer to question 2.
\end{enumerate}
\end{exercice}

\begin{corrige}[Exercise 1]
1. $M = \dfrac{1}{0.2} = 5$.

2. Total deposits $= 2\,000\,000 \times 5 = 10\,000\,000$ FCFA; new credit $= 10\,000\,000 - 2\,000\,000 = 8\,000\,000$ FCFA.

3. Bank A: deposit $2\,000\,000$, reserves $400\,000$, loan $1\,600\,000$. Bank B: deposit $1\,600\,000$, reserves $320\,000$, loan $1\,280\,000$.

4. Any two of: cash held by the public (leakages); banks keeping excess reserves; lack of creditworthy borrowers; central bank restrictions. Each reduces the amount redeposited or relent, so the real multiplier is below $5$.
\end{corrige}

\coursec{Section 3: The Central Bank and the Case of BEAC (TU 32 — Lesson 92; TU 33 — Guided Work 5)}

In Section 2, we saw that commercial banks in Cameroon are powerful institutions: by granting credit, they \emph{create} money, and this power can be a source of growth but also of instability — excessive credit creation can fuel inflation, while poorly managed banks can fail and destroy depositors' savings. A natural question then arises: \emph{who watches over the watchers?} Who controls the total quantity of money in the economy, who rescues a bank in difficulty, and who guarantees that the banknote in your pocket in Douala is worth exactly the same in N'Djamena or Libreville? The answer is a single, very special institution: the \cle{central bank}. In our part of the world, that institution is the \cle{BEAC} — the Bank of Central African States — whose headquarters stands in Yaoundé.

\courssub{Definition and Distinctive Features of a Central Bank}

\begin{definition}[Central bank]
A \cle{central bank} is the apex monetary institution responsible for issuing currency and for controlling the money supply and credit in an economy. It is a public (or publicly mandated) institution that does not seek profit and does not deal directly with ordinary individuals or firms.
\end{definition}

To understand what makes a central bank \emph{different} from a commercial bank, think of a football league. Commercial banks (Afriland First Bank, Société Générale Cameroun, BICEC\ldots) are the \emph{players}: they compete for deposits and customers and try to make profit. The central bank is the \emph{referee and rule-maker}: it does not play to win, but without it the game would collapse into disorder. Three distinctive features summarise this special status:

\begin{itemize}
\item \textbf{Bank of issue:} it alone has the legal right to issue (print and mint) the currency that circulates in the economy. Every FCFA banknote you hold is a liability of the BEAC.
\item \textbf{Bankers' bank:} commercial banks keep accounts at the central bank, borrow from it when short of liquidity, and settle payments between themselves through it.
\item \textbf{Government's bank:} it keeps the accounts of the State, manages public debt operations and advises the government on financial matters.
\end{itemize}

\begin{attention}[A frequent confusion]
The central bank is \textbf{not} a commercial bank and does \textbf{not} open accounts for individuals or businesses. If a candidate writes that ``people deposit their savings at the BEAC'', the examiner will penalise it. Ordinary customers deal with commercial banks; commercial banks deal with the central bank.
\end{attention}

\begin{popfigure}
\begin{center}
\begin{tikzpicture}[scale=1, every node/.style={font=\small}]
\node[draw=popblue, fill=popblueL, rounded corners, thick, minimum width=4.6cm, minimum height=1.3cm, align=center] (beac) at (0,0) {\textbf{BEAC}\\ (Central Bank)};
\node[draw=popgreen, fill=popgreenL, rounded corners, thick, minimum width=4.2cm, minimum height=1.1cm, align=center] (gov) at (-5.4,2.6) {Governments of\\ CEMAC States};
\node[draw=popgreen, fill=popgreenL, rounded corners, thick, minimum width=4.2cm, minimum height=1.1cm, align=center] (banks) at (5.4,2.6) {Commercial\\ Banks};
\node[draw=popgreen, fill=popgreenL, rounded corners, thick, minimum width=4.2cm, minimum height=1.1cm, align=center] (eco) at (0,-3.1) {The Economy\\ (firms \& households)};
\draw[->, very thick, popdark] (beac) -- (gov) node[midway, above left, align=center, font=\scriptsize] {banker \&\\ adviser};
\draw[->, very thick, popdark] (beac) -- (banks) node[midway, above right, align=center, font=\scriptsize] {refinancing,\\ supervision};
\draw[->, very thick, popdark] (beac) -- (eco) node[midway, right, align=center, font=\scriptsize] {currency issue,\\ price stability};
\draw[->, thick, popmuted] (banks) -- (eco) node[midway, right, font=\scriptsize] {credit \& deposits};
\end{tikzpicture}
\end{center}
\legende{The central position of the BEAC: it links the government, the commercial banks and the economy.}
\end{popfigure}

\courssub{Functions of a Central Bank}

The features above translate into six great functions, each of which can be illustrated with the BEAC.

\textbf{1. Monopoly of currency issue.} The central bank alone prints banknotes and mints coins, and it guarantees their value. This monopoly avoids the chaos that would result if every bank issued its own notes. In CEMAC, only the BEAC issues the \cle{FCFA} (franc of the Financial Cooperation in Central Africa). The BEAC determines the volume of fiduciary money in circulation based on the needs of the economy and its monetary policy stance.

\textbf{2. Banker to commercial banks — lender of last resort.} Commercial banks hold accounts at the central bank, through which interbank payments are settled, and they can borrow from it when they face a temporary shortage of cash.

\begin{definition}[Lender of last resort]
The \cle{lender of last resort} is the central bank's role of lending to commercial banks facing \emph{temporary liquidity crises}, so that a solvent bank does not collapse simply because it cannot meet sudden withdrawals. This function protects depositors and prevents a single bank's difficulties from spreading panic through the whole system.
\end{definition}

At the BEAC, this role is fulfilled through the \textbf{refinancing facilities} (appels d'offres, guichets de prêt marginal, pensions livrées) where banks obtain central bank money against eligible collateral (government securities, eligible private bills).

\textbf{3. Banker and adviser to the government.} The central bank keeps the Treasury's accounts, manages the State's payments and public debt, and advises the government on economic and financial policy. The BEAC performs this role for each of the six CEMAC Treasuries. It manages the \emph{compte d'opérations} of each Treasury and ensures the execution of budgetary expenditures.

\textbf{4. Management of foreign exchange reserves.} The central bank holds and manages the country's stock of foreign currencies (dollars, euros\ldots) needed to pay for imports and to defend the external value of the currency. CEMAC's foreign reserves are \emph{pooled} at the BEAC — a key feature of the monetary union. The BEAC invests these reserves in safe, liquid assets abroad. The level of gross foreign reserves (in months of imports) is a critical indicator of the zone's external sustainability.

\textbf{5. Supervision of the banking system.} The central bank licenses banks, sets prudential rules (minimum capital, liquidity ratios, risk concentration limits) and inspects them. In CEMAC this task is carried out by the \cle{COBAC} (Commission Bancaire de l'Afrique Centrale — Central African Banking Commission), the supervisory arm attached to the BEAC. COBAC ensures compliance with the \emph{Règlement COBAC} (prudential regulations) and can sanction or withdraw licences.

\textbf{6. Conduct of monetary policy.} This is the most macroeconomic function: the central bank regulates the quantity of money and credit in order to reach the objectives of monetary policy (see below). The BEAC's Monetary Policy Committee (CPM — Comité de Politique Monétaire) meets regularly to set the stance.

\begin{savaistu}[GCE exam tip]
Whenever a question asks you to ``explain the functions of a central bank'', \textbf{always illustrate each function with a BEAC example} (issue of the FCFA, TIAO, COBAC, pooled reserves, SYSTAC payment system\ldots). Application marks are reserved for candidates who root their answers in the Cameroonian/CEMAC context.
\end{savaistu}

\courssub{Objectives and Instruments of Monetary Policy}

\textbf{Objectives.} Monetary policy pursues three great objectives, ranked by the BEAC's statutes:
\begin{itemize}
\item \cle{Price stability}: keeping inflation low and predictable — the \textbf{primary objective} of the BEAC (inflation target around 3\% per year for the zone);
\item \cle{Support to growth and employment}: ensuring that credit conditions allow firms to invest and create jobs, without prejudice to price stability;
\item \cle{External stability of the currency}: maintaining the value and convertibility of the FCFA, which is pegged to the euro at a fixed parity (1 EUR = 655.957 FCFA).
\end{itemize}

\textbf{Instruments.} To reach these objectives, the central bank disposes of five main instruments. At the BEAC, they operate within a \emph{liquidity management framework} where the central bank steers short-term market rates:

\begin{itemize}
\item \textbf{The key (policy) interest rate — TIAO:} the rate at which the central bank lends to commercial banks through its main refinancing operations (appels d'offres hebdomadaires). At the BEAC it is called the \cle{TIAO} (\emph{taux d'intérêt des appels d'offres} — tender rate). When the TIAO rises, refinancing becomes dearer for banks, which raise their own lending rates.
\item \textbf{The marginal lending rate — TIP:} the \cle{TIP} (\emph{taux d'intérêt pénal}) is the rate at which banks can obtain overnight liquidity from the BEAC on their own initiative, against eligible collateral. It acts as a ceiling for the interbank overnight rate.
\item \textbf{The deposit facility rate — TIR:} the \cle{TIR} (\emph{taux d'intérêt de référence}) is the rate paid on banks' overnight deposits at the BEAC. It acts as a floor for the interbank overnight rate. The corridor (TIP – TIR) is typically 200 basis points.
\item \textbf{Reserve requirements:} the obligatory fraction of deposits that each commercial bank must keep, unlent, at the central bank. In CEMAC, the \cle{coefficient de réserves obligatoires} is currently 7\% (subject to change by CPM decision). Raising the reserve ratio directly reduces the credit multiplier studied in Section 2. Reserves are remunerated at a rate linked to the TIAO.
\item \textbf{Open market operations:} buying or selling securities (mainly Treasury bills and bonds) to inject or withdraw liquidity from the banking system. The BEAC conducts weekly \emph{appels d'offres} (main refinancing operations) and can conduct \emph{pensions livrées} (repos) or outright purchases/sales.
\item \textbf{Refinancing / rediscount facilities:} the standing facilities through which banks obtain central bank money against eligible bills (TIP facility, rediscount of eligible commercial paper).
\item \textbf{Moral suasion:} informal pressure, persuasion and directives addressed to banks to encourage desired behaviour (e.g.\ restraining credit to speculative activities, directing credit to priority sectors like agriculture or SMEs).
\end{itemize}

\begin{propriete}[Effects of monetary policy stances]
A \cle{contractionary} (restrictive) monetary policy — \textbf{higher key rate (TIAO), higher reserve requirements} — reduces the money supply and credit, and is used to \textbf{fight inflation}. An \cle{expansionary} policy — \textbf{lower key rate (TIAO), lower reserve requirements} — increases the money supply and credit, and is used to \textbf{stimulate growth and employment}. (The mechanism itself, not a formal proof, is examinable.)
\end{propriete}

The logic behind this property is the \emph{transmission mechanism} of monetary policy, summarised below.

\begin{popfigure}
\begin{center}
\begin{tikzpicture}[scale=1, every node/.style={font=\small}]
\tikzset{box/.style={draw=popblue, fill=popblueL, rounded corners, thick, minimum width=3.1cm, minimum height=1.05cm, align=center}}
\node[box, align=center] (r) at (0,0){Key rate\\ (TIAO) $\uparrow$};
\node[box, align=center] (l) at (4.1,0){Bank lending\\ rates $\uparrow$};
\node[box, align=center] (c) at (8.2,0){Credit \& money\\ supply $\downarrow$};
\node[box, align=center] (d) at (4.1,-2.6){Aggregate\\ demand $\downarrow$};
\node[box, draw=popgreen, fill=popgreenL, align=center] (i) at (8.2,-2.6){Inflation\\ $\downarrow$};
\draw[->, very thick, popdark] (r) -- (l);
\draw[->, very thick, popdark] (l) -- (c);
\draw[->, very thick, popdark] (c) -- (d);
\draw[->, very thick, popdark] (d) -- (i);
\end{tikzpicture}
\end{center}
\legende{The transmission mechanism of a contractionary monetary policy.}
\end{popfigure}

\begin{methode}[Analysing any central bank decision]
When the GCE paper presents a decision of a central bank (e.g.\ ``the BEAC raises the TIAO''), proceed in three steps:
\begin{enumerate}
\item \textbf{Identify the instrument} used (key rate, reserve requirement, open market operation\ldots);
\item \textbf{Trace the direction of the effect} on money supply and credit (expansion or contraction), using the transmission mechanism;
\item \textbf{State the final objective} targeted (fighting inflation, supporting growth, defending the currency\ldots).
\end{enumerate}
\end{methode}

\courssub{Guided Work 5: Case Study of the BEAC}

\textbf{Presentation of the institution.} The \cle{Bank of Central African States} (BEAC, \emph{Banque des États de l'Afrique Centrale}) was created in \textbf{November 1972} as the successor of the \emph{Banque Centrale des États de l'Afrique Équatoriale et du Cameroun} (BCEAEC). Its headquarters are in \textbf{Yaoundé, Cameroon}. It is the \emph{common} central bank of the six member states of \cle{CEMAC} (Communauté Économique et Monétaire de l'Afrique Centrale — Economic and Monetary Community of Central Africa): \textbf{Cameroon, Chad, the Central African Republic, Congo (Brazzaville), Gabon and Equatorial Guinea}. It issues the \cle{FCFA}, which circulates identically in all six countries, and it cooperates with the \textbf{French Treasury}, which guarantees the convertibility of the FCFA into euros at a fixed parity — the anchor of the currency's external stability. The BEAC operates under the \emph{Statuts de la BEAC} and the \emph{Convention de garantie} with France.

\begin{attention}[The classic mistake]
Do \textbf{not} write that the BEAC ``belongs to Cameroon'' or is ``Cameroon's central bank''. The BEAC is the \emph{common} central bank of the \textbf{six CEMAC states}; its headquarters simply happen to be in Yaoundé. Confusing the seat of an institution with its ownership is a frequent error that costs marks.
\end{attention}

\begin{popfigure}
\begin{center}
\begin{tikzpicture}[scale=0.95, every node/.style={font=\scriptsize}]
% Cameroon
\draw[fill=popgreenL, draw=popdark, thick] (0,0) -- (1.6,-0.4) -- (1.8,-1.8) -- (1.2,-3.0) -- (0.2,-3.2) -- (-0.6,-1.6) -- cycle;
\node[align=center] at (0.55,-1.6) {\textbf{Cameroon}\\ (Yaoundé: HQ)};
% Chad
\draw[fill=popblueL, draw=popdark, thick] (0.3,0.2) -- (2.2,0.1) -- (2.4,1.6) -- (0.6,1.9) -- cycle;
\node[align=center] at (1.35,1.0) {\textbf{Chad}};
% CAR
\draw[fill=popblueL, draw=popdark, thick] (2.0,-0.6) -- (3.6,-0.8) -- (3.4,-2.4) -- (2.0,-2.0) -- cycle;
\node[align=center] at (2.8,-1.5) {\textbf{Central African}\\ \textbf{Republic}};
% Congo
\draw[fill=popblueL, draw=popdark, thick] (1.9,-2.3) -- (3.2,-2.6) -- (3.0,-3.9) -- (1.5,-3.4) -- cycle;
\node[align=center] at (2.35,-3.2) {\textbf{Congo}};
% Gabon
\draw[fill=popblueL, draw=popdark, thick] (0.4,-3.4) -- (1.4,-3.5) -- (1.2,-4.7) -- (0.1,-4.5) -- cycle;
\node[align=center] at (0.75,-4.05) {\textbf{Gabon}};
% Eq. Guinea
\draw[fill=popblueL, draw=popdark, thick] (-0.5,-3.5) -- (0.3,-3.6) -- (0.1,-4.4) -- (-0.6,-4.2) -- cycle;
\node[align=center] at (-0.2,-3.95) {\textbf{Eq. Guinea}};
% HQ star
\node[draw=poporange, fill=poporange, thick, star, star points=5, inner sep=1.5pt] at (0.55,-1.35) {};
\end{tikzpicture}
\end{center}
\legende{Schematic map of the six CEMAC member states sharing the BEAC (star: Yaoundé, headquarters).}
\end{popfigure}

\textbf{Governance.} The BEAC is governed by a \textbf{Governor} (appointed by the Conference of Heads of State) and a \textbf{Board of Directors} (Conseil d'Administration). The \textbf{Monetary Policy Committee (CPM)} — composed of the Governor, Vice-Governors, and national directors — sets monetary policy. The \textbf{COBAC} (Commission Bancaire de l'Afrique Centrale) is the independent supervisory body.

\textbf{BEAC in action.} The BEAC's Monetary Policy Committee meets regularly and takes decisions that directly affect Cameroonian households and firms:
\begin{itemize}
\item it sets the \cle{TIAO}, the key rate at which it refinances commercial banks through weekly tenders — in recent years (2022–2024) it has been \emph{raised} several times (e.g., from 3.50\% to 5.00\%) to combat the inflationary pressures felt across CEMAC;
\item it fixes the \cle{TIP} (marginal lending rate) and \cle{TIR} (deposit facility rate), defining the interest rate corridor;
\item it fixes \textbf{reserve requirements} (coefficient de réserves obligatoires) for banks, thereby controlling their capacity to create credit;
\item through \textbf{COBAC}, it supervises banks and microfinance institutions, sanctioning those that breach prudential rules (capital adequacy, liquidity, risk concentration);
\item it manages the \textbf{pooled foreign exchange reserves} of the six states, essential for paying for imports of fuel, medicine and machinery;
\item it operates the \textbf{SYSTAC} (Système de Télécompensation de l'Afrique Centrale), the real-time gross settlement (RTGS) system for large-value payments and the automated clearing house for retail payments.
\end{itemize}

\begin{exoresolu}[Analysing a BEAC decision]
\textbf{Statement.} Facing rising inflation in the CEMAC zone, the BEAC decides to increase its key interest rate (TIAO) by 50 basis points. (a) Identify the instrument used. (b) Explain, step by step, how this decision is expected to reduce inflation. (c) State one possible cost of this policy for the Cameroonian economy.
\tcblower
\textbf{Solution.} (a) The instrument is the \emph{key (policy) interest rate} — the TIAO, the rate at which the central bank refinances commercial banks through its main refinancing operations (appels d'offres). (b) Transmission mechanism: the rise in the TIAO makes refinancing \emph{dearer} for commercial banks; they react by raising their own lending rates (prime rate, overdraft rates); credit becomes more expensive, so households and firms borrow less; the money supply and aggregate demand grow more slowly; with less demand pressure on goods markets, \emph{inflation falls}. (c) A possible cost: dearer credit discourages \emph{investment} by Cameroonian firms (e.g., in agro-industry or manufacturing) and may slow economic growth and job creation in the short run — the classic trade-off between fighting inflation and supporting activity. If the inflation is \emph{cost-push} (e.g., imported food/fuel prices), the policy may be less effective and more costly in terms of output loss.
\end{exoresolu}

\textbf{Critical appreciation (exam enrichment).} A strong GCE essay ends with evaluation. The BEAC presents clear \emph{strengths}: a shared central bank gives six small economies a credible, stable currency; the French Treasury guarantee secures convertibility; pooling reserves protects each member better than national reserves would; the fixed peg eliminates exchange rate uncertainty within the zone and with the euro area. But there are also \emph{limits}: one single monetary policy must suit six economies with different structures and shocks (oil-dependent Gabon or Equatorial Guinea versus more diversified Cameroon); the fixed peg to the euro removes the exchange rate as an adjustment tool (no depreciation to absorb external shocks); and decisions taken for the whole zone may not match Cameroon's specific needs at a given moment (asymmetric shocks). Mentioning this balance — without exaggeration — distinguishes an excellent script.

\begin{aretenir}[Key points of Section 3]
\begin{itemize}
\item A central bank is the apex monetary institution: \textbf{bank of issue, bankers' bank, government's bank}.
\item Its six functions: currency issue, lender of last resort, banker/adviser to government, management of foreign reserves, banking supervision (via COBAC), conduct of monetary policy.
\item Instruments: key rate (TIAO), marginal lending rate (TIP), deposit facility rate (TIR), reserve requirements (currently 7\%), open market operations (appels d'offres, pensions livrées), refinancing facilities, moral suasion.
\item Contractionary policy (higher TIAO, higher reserves) $\Rightarrow$ less money/credit, fights inflation; expansionary policy (lower TIAO, lower reserves) $\Rightarrow$ more money/credit, supports growth.
\item The BEAC (1972, Yaoundé) is the \emph{common} central bank of the six CEMAC states, issuing the FCFA under a convertibility guarantee of the French Treasury (fixed parity 1 EUR = 655.957 FCFA), with COBAC as supervisor and SYSTAC as payment system.
\item Primary objective: price stability (inflation target ~3\%). Secondary: support growth, external stability.
\end{itemize}
\end{aretenir}

\begin{exercice}[Check your understanding]
\begin{enumerate}
\item Define a central bank and state its three distinctive features.
\item Explain the role of ``lender of last resort'' and why it protects depositors.
\item The BEAC lowers reserve requirements from 7\% to 5\%. Analyse the expected effects on credit creation, money supply and economic activity.
\item Give two strengths and two limits of sharing a common central bank within CEMAC.
\item Distinguish between the TIAO, TIP and TIR at the BEAC. What is the typical corridor width?
\end{enumerate}
\end{exercice}

\begin{corrige}[Exercise n°3]
Lowering reserve requirements from 7\% to 5\% is an \emph{expansionary} measure. Commercial banks are obliged to keep a smaller fraction of deposits unlent at the BEAC, so their excess reserves rise. Through the credit multiplier (Section 2), each franc of primary deposit now supports a \emph{larger} volume of loans: the theoretical multiplier rises from $1/0.07 \approx 14.3$ to $1/0.05 = 20$. Credit creation expands and the money supply increases. Cheaper and more abundant credit encourages consumption and investment by Cameroonian households and firms, raising aggregate demand, output and employment. The risk, to be mentioned for full marks, is that if the economy is already near full capacity, the extra demand may translate into \emph{inflation} rather than real growth.
\end{corrige}

\begin{corrige}[Exercise n°5]
The \textbf{TIAO} (taux d'intérêt des appels d'offres) is the main policy rate — the rate at which banks obtain liquidity through weekly tenders (main refinancing operations). The \textbf{TIP} (taux d'intérêt pénal) is the marginal lending facility rate — the rate for overnight borrowing on banks' own initiative, acting as a ceiling for the interbank overnight rate. The \textbf{TIR} (taux d'intérêt de référence) is the deposit facility rate — the rate paid on banks' overnight deposits at the BEAC, acting as a floor. The typical corridor width is \textbf{200 basis points} (TIP – TIR = 2\%), with the TIAO usually set in the middle.
\end{corrige}

Having understood how the formal banking system is steered from the top by the BEAC, we turn in the next section to institutions that bring financial services to those whom commercial banks often neglect: the \emph{microfinance} sector.

\coursec{Section 4: Microfinance Activities in Cameroon (TU 34 — Lesson 93)}

In the previous section, we studied the \cle{central bank} and examined the case of the \cle{BEAC}, the institution that issues currency, supervises commercial banks and conducts monetary policy in the CEMAC zone. Yet a striking paradox remains: despite the existence of a central bank and a network of commercial banks, the large majority of Cameroonians — small farmers in the North West, market women in Douala, motorcycle-taxi riders in Yaoundé, artisans in Bafoussam — have historically had \emph{no access} to ordinary bank credit. Commercial banks demand collateral, payslips, business plans and minimum deposits that these economic agents simply cannot provide. It is precisely to fill this gap that \cle{microfinance} emerged. This section defines microfinance, describes the categories of microfinance institutions (MFIs) operating in Cameroon, analyses their services and importance, and evaluates the problems they face.

\courssub{Definition and Origin of Microfinance}

\begin{definition}[Microfinance]
\cle{Microfinance} is the provision of small-scale financial services — savings, credit, insurance and money transfers — to low-income people who are excluded from the formal banking system. \cle{Microcredit} is the specific component of microfinance that consists of granting very small loans to poor or low-income borrowers, usually to finance income-generating activities.
\end{definition}

The intuition is simple. A commercial bank will not lend 50,000 FCFA to a woman selling roasted plantain by the roadside: the loan is too small to be profitable for the bank, and she has no land title or salary slip to offer as security. Yet that 50,000 FCFA may be exactly what she needs to buy a larger stock of plantain and charcoal, double her daily sales, and gradually lift her household out of poverty. Microfinance therefore rests on a powerful idea: \emph{the poor are bankable} — they can save and they can repay, provided financial services are designed for their reality.

The modern microfinance movement is generally associated with the experience of the Grameen Bank founded by Muhammad Yunus in Bangladesh, which demonstrated that poor borrowers, especially women organised in small groups, repay loans at remarkably high rates when lending is based on \cle{group guarantees} rather than physical collateral. The idea spread worldwide, and in Cameroon it took root through cooperative movements, church-linked credit unions and village savings groups, building on the long tradition of \emph{tontines} (rotating savings associations, known locally as \emph{njangi} in the North West and South West).

\begin{propriete}[Social collateral in microcredit]
Microcredit typically relies on \cle{group guarantees} and \cle{social collateral} instead of physical collateral: borrowers form small solidarity groups whose members guarantee one another's loans, so that peer pressure and community reputation replace land titles and salary slips as security for repayment.
\end{propriete}

\begin{savaistu}[The MC\textsuperscript{2} model]
Cameroon developed its own celebrated microfinance model: the \cle{MC\textsuperscript{2}} network (\emph{Mutuelles Communautaires de Croissance}), born in the Bamiléké highlands of the West Region. In an MC\textsuperscript{2}, members of a community pool their savings in a community-owned fund and lend to one another for small projects. The model became so successful that it was studied and replicated in other African countries — a fine example of Cameroonian financial innovation.
\end{savaistu}

\courssub{Categories of Microfinance Institutions in Cameroon}

Microfinance institutions in Cameroon take several organisational forms:

\begin{itemize}
\item \textbf{Credit unions and savings and credit cooperatives (SACCOs):} member-owned cooperatives that collect savings from members and grant them loans. The largest network is \cle{CAMCCUL} (Cameroon Cooperative Credit Union League), which regroups hundreds of credit unions across the country.
\item \textbf{Village banks and community funds:} small, community-based institutions (such as the MC\textsuperscript{2}) serving a single village or neighbourhood, often with strong support from NGOs or development partners.
\item \textbf{NGOs and specialised microfinance companies:} non-governmental organisations and private companies that provide microcredit and related services, sometimes without collecting deposits.
\end{itemize}

\begin{definition}[Credit union]
A \cle{credit union} is a member-owned cooperative financial institution that pools the savings of its members in order to grant them loans at moderate interest rates. Members are simultaneously owners and customers: they elect the management committee and share in any surplus according to cooperative principles.
\end{definition}

\begin{propriete}[The COBAC classification of MFIs]
In the CEMAC zone, MFIs are classified into three categories and supervised by the \cle{Ministry of Finance} and the \cle{COBAC} (Commission Bancaire de l'Afrique Centrale):
\begin{enumerate}
\item \textbf{First category:} cooperatives that \emph{both} collect savings from members \emph{and} grant them credit (e.g., credit unions affiliated to CAMCCUL).
\item \textbf{Second category:} institutions that grant \emph{credit only}, financing their lending from their own funds or borrowing (they may not collect deposits from the public).
\item \textbf{Third category:} institutions that collect \emph{savings only}, without granting credit.
\end{enumerate}
\end{propriete}

\begin{attention}[Common mistake: "MFIs are unregulated"]
A frequent error in GCE essays is to claim that microfinance institutions operate outside any legal framework. This is false. In Cameroon, MFIs are classified into the three COBAC categories above, must be licensed, and are supervised by the Ministry of Finance and the COBAC. What is true — and worth writing — is that \emph{some unlicensed or fraudulent institutions} operate illegally alongside the regulated sector, which is a problem of enforcement, not an absence of regulation.
\end{attention}

\courssub{Services Offered by Microfinance Institutions}

Modern MFIs in Cameroon offer a growing range of services tailored to low-income clients:

\begin{itemize}
\item \textbf{Micro-savings:} small, flexible deposit accounts with very low minimum balances, allowing market traders and farmers to save daily or weekly surpluses safely.
\item \textbf{Microcredit:} small loans (from a few thousand to a few hundred thousand FCFA) for petty trade, agriculture, artisanal production or school fees, often disbursed through solidarity groups.
\item \textbf{Micro-insurance:} cheap insurance products covering illness, death or crop failure, protecting poor households from shocks that would otherwise push them deeper into poverty.
\item \textbf{Money transfers and mobile banking linkages:} many MFIs now partner with mobile money operators (MTN Mobile Money, Orange Money) so that members can deposit, withdraw and repay loans by telephone, extending services deep into rural areas.
\end{itemize}

\courssub{Importance of Microfinance in Cameroon}

Microfinance plays a strategic role in the Cameroonian economy:

\begin{enumerate}
\item \textbf{Financial inclusion:} MFIs reach rural populations, women and informal-sector operators whom commercial banks consider too risky or too costly to serve. Women, who dominate petty trade and food-crop farming, form the majority of members of many credit unions and \emph{njangi} groups.
\item \textbf{Financing of small businesses and agriculture:} microcredit finances market stalls, poultry projects, cocoa and food-crop farms, tailoring workshops and transport activities — the backbone of the informal economy that employs most Cameroonians.
\item \textbf{Job creation and poverty reduction:} by financing income-generating activities, microfinance helps households create self-employment, smooth consumption between harvests, pay school fees and accumulate small assets.
\item \textbf{Mobilisation of domestic savings:} MFIs transform scattered rural savings into loanable funds, deepening financial intermediation in areas where no commercial bank branch exists.
\end{enumerate}

\begin{exemplebox}[A microcredit chain in Bafoussam]
Mama Rose, a trader in Bafoussam main market, joins a solidarity group of five women at a local credit union. She borrows 100,000 FCFA for six months at a simple monthly interest rate of 2\%, guaranteed by her group. The interest due is $100{,}000 \times 0.02 \times 6 = 12{,}000$ FCFA, so she repays 112,000 FCFA in total. She buys a stock of second-hand clothing, resells it over the season, earns a trading profit of about 60,000 FCFA, repays the loan with interest, and keeps roughly $60{,}000 - 12{,}000 = 48{,}000$ FCFA as net gain. On the strength of her repayment record, she qualifies for a larger loan of 200,000 FCFA the following season. This graduated lending — small first loan, larger loans after successful repayment — is the classic microcredit mechanism.
\end{exemplebox}

\courssub{Problems of Microfinance Institutions in Cameroon}

Despite their achievements, MFIs in Cameroon face serious difficulties:

\begin{itemize}
\item \textbf{Weak management and governance:} many cooperatives are run by elected committees with limited training in accounting, credit analysis and risk management, leading to poor decisions and sometimes mismanagement of members' funds.
\item \textbf{High default rates:} borrowers may divert loans to consumption, suffer business failure, or simply refuse to repay, especially where group solidarity weakens; defaults erode the capital of small MFIs.
\item \textbf{Insufficient capital:} MFIs depend on members' savings and occasional donor support; their small capital base limits the size and number of loans they can grant and makes them fragile in a crisis.
\item \textbf{Fraud and failures of unlicensed institutions:} fraudulent "wonder banks" and unlicensed operators have collected deposits and disappeared, destroying public confidence — the collapse of dubious savings schemes in the past made many Cameroonians wary of all MFIs.
\item \textbf{High interest rates:} because small loans are costly to administer and default risk is high, MFI interest rates are often much higher than commercial bank rates, which can burden borrowers and limit the poverty-reduction impact.
\end{itemize}

\courssub{Commercial Banks versus Microfinance Institutions}

\begin{popfigure}
\begin{center}
\begin{tikzpicture}[scale=1]
\fill[popblueL] (0,0) rectangle (6.4,1);
\fill[popgreenL] (6.4,0) rectangle (12.8,1);
\node[font=\bfseries] at (3.2,0.5) {Commercial banks};
\node[font=\bfseries] at (9.6,0.5) {Microfinance institutions};
\node[align=left, text width=5.8cm, anchor=north west, font=\small] at (0.2,-0.2)
{\textbullet~Large firms, salaried workers\\[2pt]
 \textbullet~Large loans (millions of FCFA)\\[2pt]
 \textbullet~Physical collateral required\\[2pt]
 \textbullet~Lower interest rates\\[2pt]
 \textbullet~Supervised by COBAC / BEAC};
\node[align=left, text width=5.8cm, anchor=north west, font=\small] at (6.6,-0.2)
{\textbullet~Low-income, informal sector\\[2pt]
 \textbullet~Small loans (thousands of FCFA)\\[2pt]
 \textbullet~Group / social collateral\\[2pt]
 \textbullet~Higher interest rates\\[2pt]
 \textbullet~Min. of Finance / COBAC categories};
\draw[popline, thick] (0,0) rectangle (12.8,1);
\draw[popline, thick] (6.4,0) -- (6.4,1);
\draw[popline, thick] (0,-4.6) rectangle (12.8,0);
\draw[popline, thick] (6.4,-4.6) -- (6.4,0);
\end{tikzpicture}
\end{center}
\legende{Comparison of commercial banks and microfinance institutions across key criteria}
\end{popfigure}

The two types of institution are \emph{complementary} rather than rivals: commercial banks finance the formal sector and large projects, while MFIs bank the unbanked at the base of the economy. A successful micro-entrepreneur may eventually "graduate" from a credit union to a commercial bank.

\begin{popfigure}
\begin{center}
\begin{tikzpicture}
\begin{axis}[
    width=9cm, height=6.5cm,
    ybar, bar width=22pt,
    ymin=0, ymax=45,
    ylabel={Share of MFI clientele (\%)},
    symbolic x coords={Traders, Farmers, Women groups, Salaried},
    xtick=data,
    x tick label style={font=\small},
    nodes near coords, nodes near coords style={font=\small},
    every axis plot/.append style={fill=poporange, draw=popdark},
]
\addplot coordinates {(Traders,38) (Farmers,27) (Women groups,22) (Salaried,13)};
\end{axis}
\end{tikzpicture}
\end{center}
\legende{Illustrative distribution of MFI clientele in Cameroon (traders 38\%, farmers 27\%, women groups 22\%, salaried workers 13\%) — indicative shares for classroom discussion, not official statistics}
\end{popfigure}

\begin{exoresolu}[Analysing the role of microfinance]
\textbf{Statement.} A credit union in Bamenda has 2,000 members. Each member saves on average 25,000 FCFA per year. The credit union lends out 80\% of its deposits as microloans at an average interest rate of 18\% per year, and pays 4\% per year on members' savings. (a) Calculate the total annual savings mobilised. (b) Calculate the amount lent out as microloans. (c) Calculate the credit union's gross interest margin on the lent funds. (d) State two reasons why this margin must be higher than that of a commercial bank.
\tcblower
\textbf{Solution.}
\begin{enumerate}
\item Total savings $= 2{,}000 \times 25{,}000 = 50{,}000{,}000$ FCFA (50 million FCFA).
\item Microloans granted $= 0.80 \times 50{,}000{,}000 = 40{,}000{,}000$ FCFA.
\item Interest received on loans $= 0.18 \times 40{,}000{,}000 = 7{,}200{,}000$ FCFA. Interest paid on savings $= 0.04 \times 50{,}000{,}000 = 2{,}000{,}000$ FCFA. Gross interest margin $= 7{,}200{,}000 - 2{,}000{,}000 = 5{,}200{,}000$ FCFA.
\item The margin must be higher because (i) administering many very small loans is costly per franc lent (staff visits, group meetings, monitoring), and (ii) default risk is high since loans are granted without physical collateral, so the credit union must build reserves against losses.
\end{enumerate}
\end{exoresolu}

\begin{methode}[GCE exam tip: essays on financial inclusion]
When answering an essay on microfinance or financial inclusion:
\begin{enumerate}
\item Define microfinance and microcredit precisely.
\item Name concrete Cameroonian institutions: \cle{CAMCCUL}, the \cle{MC\textsuperscript{2}} model, \emph{njangi}/tontine traditions, mobile money linkages.
\item Use the COBAC three-category classification to show knowledge of the regulatory framework.
\item Balance the answer: importance (inclusion, jobs, rural finance) \emph{against} problems (governance, defaults, fraud, high rates).
\item Conclude with a judgement, e.g., the need for stronger supervision and capacity-building.
\end{enumerate}
\end{methode}

\begin{aretenir}[Key points of Section 4]
\begin{itemize}
\item Microfinance provides small-scale savings, credit, insurance and transfer services to low-income people excluded from formal banking; microcredit is its lending component.
\item Microcredit substitutes \cle{group guarantees} and social collateral for physical collateral.
\item Cameroonian MFIs include credit unions (CAMCCUL), village banks and MC\textsuperscript{2}, SACCOs, NGOs and specialised companies.
\item COBAC classifies MFIs into three categories: savings-and-credit cooperatives, credit-only institutions, and savings-only institutions, supervised by the Ministry of Finance and COBAC.
\item Microfinance promotes financial inclusion, finances small businesses and agriculture, and creates jobs, but suffers from weak governance, defaults, insufficient capital, fraud and high interest rates.
\end{itemize}
\end{aretenir}

\begin{exercice}[Essay practice]
(a) Distinguish between microfinance and microcredit. (b) Describe the three COBAC categories of microfinance institutions in Cameroon. (c) Discuss four contributions of microfinance to the Cameroonian economy and four problems faced by MFIs, using named examples such as CAMCCUL and the MC\textsuperscript{2} model.
\end{exercice}

\begin{corrige}[Exercise 1]
(a) Microfinance is the whole range of small-scale financial services (savings, credit, insurance, transfers) for low-income people; microcredit is only the small-loan component. (b) First category: cooperatives collecting savings and granting credit; second category: credit-only institutions; third category: savings-only institutions. (c) Contributions: financial inclusion of rural people, women and informal operators; financing of petty trade and agriculture; job creation and poverty reduction; mobilisation of rural savings. Problems: weak management and governance; high default rates; insufficient capital; fraud by unlicensed operators; high interest rates. A strong answer names CAMCCUL credit unions and the MC\textsuperscript{2} community funds, and concludes with the need for better supervision and training.
\end{corrige}

\coursec{Section 5: Appraisal of a Loan File by a Credit Union (TU 35 — Guided Work 6) and Chapter Synthesis}

In the previous section, we explored the vibrant world of \textbf{microfinance institutions (MFIs)} in Cameroon — Credit Unions (COOPECs), microfinance banks, and financial NGOs. We saw how they mobilise savings from low-income households and provide financial services to the \emph{unbanked} and \emph{underbanked}. However, the heart of any financial intermediary — whether a giant commercial bank like Afriland First Bank or a village Credit Union in Bafut — is \textbf{lending}. Lending transforms idle savings into productive investment. But every loan carries a risk: the borrower may not repay. For a Credit Union, where the capital belongs to the \emph{members themselves} (teachers, farmers, market women), a single bad loan can erode trust and threaten the institution's survival. This is why \textbf{loan appraisal} — the rigorous, systematic evaluation of a borrower's file — is not a bureaucratic formality but the \emph{cornerstone of sustainability}. In this final section, we master the art and science of credit appraisal, apply it to a realistic Cameroonian case, and then synthesise the entire financial system.

\courssub{Why Loan Appraisal Matters: Protecting Members' Savings}

\begin{definition}[Loan Appraisal]
\textbf{Loan appraisal} (or credit appraisal) is the systematic process by which a financial institution evaluates a loan application to determine the borrower's \textbf{creditworthiness} — their ability and willingness to repay the loan according to the agreed terms — before deciding to approve, reject, or modify the request.
\end{definition}

In a \textbf{Credit Union (COOPEC)}, the stakes are uniquely high:
\begin{itemize}
    \item \textbf{Ownership structure:} The capital comes from \emph{members' savings} (shares and deposits). A default directly reduces the value of members' shares or delays their withdrawals.
    \item \textbf{Social mission:} MFIs serve clients who often lack formal collateral (land titles, salaried payslips). Appraisal must rely on \emph{character} and \emph{cash-flow analysis} rather than just asset backing.
    \item \textbf{Regulatory requirement:} COBAC (Central African Banking Commission) regulations require MFIs to have written credit policies and to classify loans by risk. Poor appraisal leads to high \textbf{Non-Performing Loans (NPLs)}, triggering supervisory sanctions.
\end{itemize}

\begin{aretenir}[The Three Pillars of Sound Appraisal]
\begin{enumerate}
    \item \textbf{Protecting savers:} Ensures that the money entrusted by a market woman in Douala or a cocoa farmer in Muyuka is lent only to viable projects.
    \item \textbf{Reducing default risk:} Identifies red flags early (inflated sales figures, hidden debts, character issues).
    \item \textbf{Ensuring sustainability:} Keeps the \textbf{Portfolio at Risk (PAR > 30 days)} low, allowing the MFI to cover its costs and pay dividends on shares.
\end{enumerate}
\end{aretenir}

\courssub{The 5 Cs of Credit Appraisal: A Universal Framework}

The \textbf{5 Cs} is the globally accepted framework for credit analysis. In the Cameroonian microfinance context, each ``C'' takes on a specific flavour.

\begin{methode}[The 5 Cs Checklist for a Cameroonian Borrower]
\begin{enumerate}
    \item \textbf{Character (Reputation \& Willingness to Pay)} \\
          \textit{Question:} ``Is this person honest and known to honour commitments?'' \\
          \textit{Evidence:} Membership duration, savings regularity, testimonials from group leaders (in solidarity groups), credit bureau check (if available via \emph{GIMAC} or MFI internal database), community reputation. A member who has saved 5\,000 FCFA weekly for 3 years without a single withdrawal shows stronger character than a newcomer demanding 500\,000 FCFA.
    \item \textbf{Capacity (Repayment Ability)} \\
          \textit{Question:} ``Does the business generate enough \textbf{net cash flow} to service the debt?'' \\
          \textit{Analysis:} \textbf{Instalment-to-Income Ratio (IIR)}. Monthly instalment $\le$ \textbf{30--35\%} of net monthly household/business income. For a farmer, capacity is seasonal — repayment schedule \emph{must} align with harvest sales (e.g., cocoa main crop Oct--Jan).
    \item \textbf{Capital (Own Contribution / Skin in the Game)} \\
          \textit{Question:} ``How much of their own money is the borrower risking?'' \\
          \textit{Rule:} Most MFIs require \textbf{10--30\% own contribution} (equity). If a trader wants 1\,000\,000 FCFA to buy a stock of frozen fish, they should contribute 200\,000 FCFA from savings. This reduces moral hazard.
    \item \textbf{Collateral (Security / Last Resort)} \\
          \textit{Question:} ``If the business fails, what can we recover?'' \\
          \textit{Reality in Cameroon:} Many MFI clients lack land titles (\emph{titre foncier}). Acceptable collateral: \textbf{group guarantee (solidarity)}, \textbf{personal guarantors} (salaried civil servants), \textbf{chattel mortgage} on equipment/vehicle (\emph{nantissement}), \textbf{blocked savings} in the MFI (\emph{épargne bloquée}). \textbf{Warning: Collateral is a fallback, not the source of repayment.}
    \item \textbf{Conditions (Economic \& Policy Environment)} \\
          \textit{Question:} ``What external factors could derail repayment?'' \\
          \textit{Examples:} Cocoa price crash on world market, fuel price hike increasing transport costs, insecurity in the Northwest/Southwest disrupting market access, new tax on mobile money affecting a POS agent's margins.
\end{enumerate}
\end{methode}

\begin{popfigure}
\begin{tikzpicture}[
    node distance=2.5cm and 1.5cm,
    cnode/.style={circle, draw=popblue, thick, fill=popblueL, minimum width=2.2cm, text width=2cm, align=center, font=\small\bfseries},
    centerNode/.style={circle, draw=poporange, thick, fill=poporange!20, minimum width=2.8cm, text width=2.5cm, align=center, font=\small\bfseries},
    arrow/.style={->, thick, popdark}
]
\node[centerNode, align=center] (officer){Loan Officer\\ \small(Appraiser)};
\node[cnode, align=center] (char) [above left=of officer]{Character\\ \small(Reputation)};
\node[cnode, align=center] (cap) [above=of officer]{Capacity\\ \small(Cash Flow)};
\node[cnode, align=center] (capit) [above right=of officer]{Capital\\ \small(Own Funds)};
\node[cnode, align=center] (coll) [below right=of officer]{Collateral\\ \small(Security)};
\node[cnode, align=center] (cond) [below left=of officer]{Conditions\\ \small(Environment)};

\foreach \n in {char,cap,capit,coll,cond}
    \draw[arrow] (\n) -- (officer);
\foreach \n/\txt in {char/Integrity, cap/Affordability, capit/Commitment, coll/Recovery, cond/Context}
    \node[font=\tiny, text=popmuted] at ($(\n)!0.5!(officer)$) {\txt};
\end{tikzpicture}
\legende{The 5 Cs of Credit Appraisal centred on the Loan Officer. Each ``C'' answers a distinct question; weakness in one cannot be fully compensated by strength in another.}
\end{popfigure}

\courssub{Elements of a Complete Loan File}

A loan file (\emph{dossier de crédit}) is the evidence base. Incomplete files are the \emph{number one cause of delays}. The Credit Union's credit policy should specify the mandatory documents:

\begin{center}
\begin{tabularx}{\linewidth}{|l|X|}\hline
\rowcolor{popblueL}
\textbf{Document} & \textbf{Purpose / What the Appraiser Checks} \\ \hline
\textbf{1. Application Letter} & Formal request: amount, purpose, duration, proposed repayment schedule. Signed by applicant. \\ \hline
\textbf{2. Identification} & Valid CNI (National ID) or passport; proof of residence (utility bill, quarter head attestation). \\ \hline
\textbf{3. Membership \& Savings History} & Membership card, passbook or statement (min. 6--12 months). \textbf{Key:} Regularity, volume, voluntary vs. compulsory savings. \\ \hline
\textbf{4. Business Plan / Purpose Statement} & For business loans: description of activity, location, suppliers, customers, seasonality. For agriculture: crop calendar, acreage, input costs, expected yield. \\ \hline
\textbf{5. Income \& Cash-Flow Evidence} & \textbf{Crucial.} Daily sales notebook (\emph{cahier de caisse}), bank/MFI statements, purchase/sale invoices, mobile money (MoMo) transaction history. Appraiser reconstructs \textbf{Net Operating Income}. \\ \hline
\textbf{6. Guarantors' Files} & CNI, proof of income (payslip for civil servants), signed guarantee letter (\emph{lettre d'engagement de caution}). Guarantor's own debt burden checked. \\ \hline
\textbf{7. Collateral Documents} & Land title (\emph{titre foncier}) or \emph{certificat foncier}, vehicle registration (\emph{carte grise}), chattel mortgage deed (\emph{acte de nantissement}), blocked savings agreement. \\ \hline
\textbf{8. Field Visit Report} & Loan officer's on-site notes: business exists? Stock matches claims? Premises condition? Neighbour testimony? GPS coordinates/photos. \\ \hline
\end{tabularx}
\end{center}

\courssub{The Appraisal Process: Step-by-Step Flowchart}

The process must be standardised to avoid arbitrariness and ensure audit trails.

\begin{popfigure}
\begin{tikzpicture}[
    node distance=1.8cm and 0.5cm,
    proc/.style={rectangle, draw=popblue, thick, fill=popblueL, rounded corners, minimum height=1.2cm, text width=3.2cm, align=center, font=\small},
    dec/.style={diamond, draw=poporange, thick, fill=poporange!20, aspect=2, minimum height=1.2cm, text width=3.5cm, align=center, font=\small\bfseries},
    term/.style={rectangle, draw=popgreen, thick, fill=popgreenL, rounded corners, minimum height=1cm, text width=3cm, align=center, font=\small\bfseries},
    arrow/.style={->, thick, popdark},
    label/.style={font=\tiny, text=popmuted, midway, sloped}
]
\node[proc, align=center] (start){1. Reception \&\\ Registration of File\\ \small(Completeness Check)};
\node[proc, align=center] (verify) [below=of start]{2. Verification \&\\ Field Visit\\ \small(Physical + Document Cross-check)};
\node[proc, align=center] (analyse) [below=of verify]{3. Financial Analysis\\ \small(Cash Flow, Ratios, 5 Cs)};
\node[dec, align=center] (decide) [below=of analyse]{4. Credit Committee\\ Decision};
\node[proc, align=center] (cond) [below left=of decide]{4a. Conditional\\ Approval\\ \small(More docs / Guarantor)};
\node[proc, align=center] (reject) [below right=of decide]{4b. Rejection\\ \small(Written reasons)};
\node[term, align=center] (disburse) [below=of cond]{5. Contract Signing\\ \& Disbursement};
\node[proc, align=center] (monitor) [below=of disburse]{6. Monitoring \&\\ Follow-up\\ \small(Repayment tracking, visits)};
\node[term, align=center] (close) [below=of monitor]{7. Loan Closure\\ \small(Final receipt, release of collateral)};

\draw[arrow] (start) -- (verify);
\draw[arrow] (verify) -- (analyse);
\draw[arrow] (analyse) -- (decide);
\draw[arrow] (decide) -| node[label, above, pos=0.25] {Approve} (cond);
\draw[arrow] (decide) -| node[label, above, pos=0.25] {Reject} (reject);
\draw[arrow] (cond) -- (disburse);
\draw[arrow] (disburse) -- (monitor);
\draw[arrow] (monitor) -- (close);
\draw[arrow, dashed] (reject) |- ++(-4,0) |- (start.west) node[label, near start, left, align=center]{Re-apply after\\improvement};
\end{tikzpicture}
\legende{Standard Loan Appraisal Process in a Cameroonian Credit Union. The Credit Committee (Step 4) is the governance safeguard — no single officer decides alone.}
\end{popfigure}

\courssub{Worked Example: Appraising a Sample Loan File (Guided Work 6)}

\begin{experience}[Case Study: Appraising Mme Ngo's Loan Application]
\textbf{Context:} \emph{COOPEC ``Espoir''} in Bafoussam. Mme \textbf{Ngo Sylvie}, 42, member for 4 years, runs a \textbf{wholesale/retail frozen fish business} at Marché B. She applies for \textbf{2\,000\,000 FCFA} over \textbf{12 months} to buy a second-hand freezer (800\,000 FCFA) and increase stock (1\,200\,000 FCFA) ahead of the Christmas peak.

\textbf{File Contents:}
\begin{itemize}
    \item \textbf{Savings:} 350\,000 FCFA (voluntary), regular weekly deposits of 5\,000--10\,000 FCFA. No withdrawals in 18 months.
    \item \textbf{Business Records:} Daily cash book for last 6 months. Average monthly \textbf{sales}: 1\,800\,000 FCFA. Average monthly \textbf{purchases}: 1\,350\,000 FCFA. Rent: 50\,000 FCFA/month. Transport/Market fees: 80\,000 FCFA/month. Household expenses (est.): 200\,000 FCFA/month.
    \item \textbf{Collateral:} Husband (civil servant, Grade B2) offers personal guarantee + blocks 500\,000 FCFA of his own savings in the COOPEC. Freezer to be purchased serves as chattel mortgage.
    \item \textbf{Guarantors:} Two fellow market traders (members >5 years) + husband.
    \item \textbf{Conditions:} Christmas season (high demand). Risk: Power cuts (freezer spoilage), fuel price hike (transport cost).
\end{itemize}

\textbf{Task:} As the Loan Officer, perform the appraisal using the 5 Cs and the \textbf{Instalment-to-Income Ratio (IIR)}. Assume a monthly interest rate of 1.5\% on declining balance (effective ~1.5\%/month, ~18\% p.a.). Monthly instalment $\approx$ 185\,000 FCFA. Write your recommendation: \textbf{Approve / Approve with Conditions / Reject}.
\end{experience}

\begin{exoresolu}[Detailed Appraisal of Mme Ngo's File]
\textbf{Step 1: Character} \\
\textbf{Strong.} 4-year member, disciplined saver (350\,000 FCFA built from small weekly deposits), no withdrawal discipline. Market reputation verified by field visit: neighbours confirm she is ``sérieuse''. No past loans (first application) — clean slate.

\textbf{Step 2: Capacity (The Quantitative Core)} \\
\begin{align*}
\text{Average Monthly Sales} &= 1\,800\,000 \\
\text{Less: Cost of Goods Sold (Purchases)} &= 1\,350\,000 \\
\text{Gross Margin} &= 450\,000 \\
\text{Less: Operating Expenses (Rent + Transport)} &= 130\,000 \\
\text{Net Business Income} &= 320\,000 \\
\text{Less: Household Expenses} &= 200\,000 \\
\text{Net Disposable Income (NDI)} &= \mathbf{120\,000\ FCFA/month}
\end{align*}
\textbf{Proposed Monthly Instalment} = 185\,000 FCFA. \\
\textbf{Instalment-to-Income Ratio (IIR)} = $\frac{185\,000}{120\,000} \times 100 = \mathbf{154\%}$.

\cle{RED FLAG!} The instalment \textbf{exceeds} her net disposable income. She cannot repay from current cash flow without cutting household consumption drastically or the business growing significantly.

\textbf{Step 3: Capital (Own Contribution)} \\
Total project cost = 2\,000\,000 FCFA. Loan requested = 2\,000\,000 FCFA (100\%). \textbf{Zero own contribution.} Violates typical MFI policy (min 10--20\%). She is asking the MFI to bear all the risk.

\textbf{Step 4: Collateral} \\
\textbf{Adequate on paper.} Husband's blocked savings (500\,000 FCFA) + chattel mortgage on freezer (800\,000 FCFA resale value) + two member guarantors. Covers 1\,300\,000 FCFA of 2\,000\,000 FCFA (65\%). But recall: \textbf{collateral is not repayment source}. Seizing a freezer in Bafoussam market takes months and yields low recovery.

\textbf{Step 5: Conditions} \\
\textbf{Favourable but risky.} Christmas peak (Dec) boosts sales 2--3x. However, power cuts (ENEO) are frequent — freezer failure = total stock loss. No backup generator in plan.

\textbf{DECISION: \textcolor{poporange}{\textbf{APPROVE WITH CONDITIONS (Restructure)}}} \\
\textbf{Justification:} Character is excellent, collateral acceptable, but \textbf{Capacity fails} (IIR > 100\%) and \textbf{Capital contribution is zero}. Granting as requested sets her up for certain default.

\textbf{Proposed Restructuring:}
\begin{enumerate}
    \item Reduce loan to \textbf{1\,200\,000 FCFA} (finances freezer + partial stock). She contributes 300\,000 FCFA from her savings (15\% own contribution) and uses 500\,000 FCFA from husband's blocked savings as \emph{internal collateral} (not spent).
    \item New instalment $\approx$ 110\,000 FCFA/month. IIR = $110\,000 / 120\,000 = \mathbf{92\%}$ — still high but manageable \textbf{if} Christmas sales materialise.
    \item \textbf{Seasonal Repayment Schedule:} Interest-only (or very low principal) Jan--Oct; balloon principal payments Nov--Dec from Christmas peak revenue.
    \item Condition: Purchase voltage stabiliser for freezer (cost 50\,000 FCFA included in loan).
    \item Mandatory monthly field visit by loan officer during first 6 months.
\end{enumerate}
\textbf{Credit Committee votes:} 4 in favour (with conditions), 1 abstention. \textbf{File approved.}
\end{exoresolu}

\begin{attention}[Common Appraisal Mistakes — GCE Traps]
\begin{itemize}
    \item \textbf{Mistake 1: ``Collateral substitutes for Capacity.''} \textbf{Wrong.} A loan officer who approves because ``the husband blocked savings'' ignores that the \emph{business} must generate cash. Seizing collateral destroys the member relationship and yields 30--50\% of value.
    \item \textbf{Mistake 2: ``Ignoring seasonality.''} \textbf{Wrong.} A cocoa farmer with 0 FCFA income Apr--Sep cannot pay equal monthly instalments. The schedule \emph{must} match the cash-flow cycle.
    \item \textbf{Mistake 3: ``Accepting inflated sales figures.''} \textbf{Wrong.} Always cross-check: \emph{stock turnover}, \emph{supplier invoices}, \emph{MoMo history}, \emph{physical stock count}. ``Oral accounting'' is not evidence.
    \item \textbf{Mistake 4: ``Skipping the field visit.''} \textbf{Wrong.} Desk appraisal alone misses: business closed, different activity, ghost employee, location risk.
\end{itemize}
\end{attention}

\courssub{Common Causes of Loan Default in Cameroonian Credit Unions \& Prevention}

\begin{center}
\begin{tabularx}{\linewidth}{|l|X|X|}\hline
\rowcolor{popblueL}
\textbf{Cause of Default} & \textbf{Mechanism} & \textbf{How Good Appraisal Prevents It} \\ \hline
\textbf{1. Over-indebtedness / Multiple borrowing} & Borrower takes loans from 2--3 MFIs + informal lenders; total debt service > income. & \textbf{Check credit bureau (GIMAC), ask for declaration of other debts, verify bank/MoMo statements for hidden repayments.} \\ \hline
\textbf{2. Diversion of funds} & Loan for ``stock'' used for wedding, funeral, house repair, or lent to a friend. & \textbf{Direct disbursement to supplier (cheque/transfer), field visit to verify stock purchase, post-disbursement monitoring visit at 2 weeks.} \\ \hline
\textbf{3. Business failure due to external shock} & Price collapse (cocoa, coffee), insecurity, fuel hike, pandemic. & \textbf{Conditions analysis:} Stress-test cash flow at -20\% revenue. \textbf{Insurance linkage} (crop/index insurance). Flexible restructuring clauses. \\ \hline
\textbf{4. Poor character / Moral hazard} & Willful refusal to pay despite ability; ``strategic default''. & \textbf{Character check:} Long membership, savings discipline, community reputation, guarantor quality. \textbf{Group lending} (peer pressure). \\ \hline
\textbf{5. Weak appraisal / ``Tick-box'' compliance} & Officer rushes file to meet targets; committee rubber-stamps. & \textbf{Institutional safeguards:} Independent credit committee, mandatory field visit sign-off, portfolio quality audits, officer incentives tied to \emph{repayment rate} not \emph{volume}. \\ \hline
\end{tabularx}
\end{center}

\courssub{Chapter Synthesis: The Cameroonian Financial System — A Unified Map}

We have journeyed through the architecture of money and banking in Cameroon:
\begin{itemize}
    \item \textbf{Section 1 (TU 29):} \textbf{Commercial Banks} — balance sheets, credit creation, the engine of broad money ($M_2/M_3$) expansion.
    \item \textbf{Section 2 (TU 30--31):} \textbf{Credit Creation Mechanics} — the money multiplier, reserve requirements, and the \emph{real-world frictions} (excess reserves, leakages, NPLs) that limit it in Cameroon.
    \item \textbf{Section 3 (TU 32--33):} \textbf{The Central Bank (BEAC)} — the conductor: monetary policy (TIAO, reserve requirements, standing facilities), banking supervision (via COBAC), lender of last resort, guardian of the CFA franc peg.
    \item \textbf{Section 4 (TU 34):} \textbf{Microfinance (MFIs)} — the inclusive frontier: COOPECs, MF banks, reaching the \emph{last mile} with solidarity lending, mobile money integration.
    \item \textbf{Section 5 (TU 35):} \textbf{Credit Appraisal} — the micro-foundation of asset quality; without it, the whole edifice crumbles.
\end{itemize}

These are not isolated silos. They form an \textbf{interconnected ecosystem}.

\begin{popfigure}
\begin{tikzpicture}[
    node distance=2.2cm and 1.5cm,
    central/.style={ellipse, draw=poppurple, very thick, fill=poppurple!10, minimum width=3.5cm, minimum height=1.8cm, align=center, font=\bfseries\small},
    inst/.style={rectangle, draw=popblue, thick, fill=popblueL, rounded corners, minimum width=3cm, minimum height=1.4cm, align=center, font=\small},
    client/.style={rectangle, draw=popgreen, thick, fill=popgreenL, rounded corners, minimum width=2.8cm, minimum height=1.2cm, align=center, font=\small},
    arrow/.style={->, thick, popdark},
    dasharrow/.style={->, dashed, thick, popmuted},
    bidir/.style={<->, thick, poporange}
]
% Central
\node[central, align=center] (beac){BEAC\\ \small(Central Bank of CEMAC)\\ \footnotesize Monetary Policy, Supervision,\\ Currency Issuer, Lender of Last Resort};

% Intermediaries - Left: Commercial Banks
\node[inst, align=center] (cbanks) [above left=of beac]{Commercial Banks\\(Afriland, BICEC, UBA,\\ SCB, CBC, etc.)\\ \footnotesize Deposit-taking, Corporate/SME\\ lending, Payment Systems, FX};
% Intermediaries - Right: MFIs
\node[inst, align=center] (mfis) [above right=of beac]{Microfinance Institutions\\(COOPECs, MF Banks, NGOs)\\ \footnotesize Financial Inclusion, Small loans,\\ Savings, Mobile Money Agents};
% Intermediaries - Bottom: Other (Postal, Fintech)
\node[inst, align=center] (other) [below=of beac]{Other Players\\(CAMPOST, Fintechs,\\ Insurance, Pension Funds)\\ \footnotesize Complementary Services};

% Clients - Households
\node[client, align=center] (hh) [left=of cbanks]{Households\\ \footnotesize Salaried, Informal,\\ Farmers, Diaspora};
% Clients - Firms
\node[client, align=center] (firms) [right=of mfis]{Enterprises\\ \footnotesize Large Corps, SMEs,\\ Agribusiness, Startups};
% Clients - Govt
\node[client, align=center] (gov) [below left=of other]{Government\\ \footnotesize Treasury, Projects,\\ Securities (BTA/OAT)};

% Flows: BEAC -> Intermediaries
\draw[arrow] (beac) -- node[label, sloped, align=center]{Refinancing (TIAO)\\ Reserve Requirements\\ Supervision} (cbanks);
\draw[arrow] (beac) -- node[label, sloped, align=center]{Refinancing window\\ COBAC Regulation\\ Financial Inclusion} (mfis);
\draw[arrow] (beac) -- node[label, sloped, align=center]{Treasury operations\\ Securities settlement} (other);

% Flows: Intermediaries -> Clients
\draw[bidir] (cbanks) -- node[label, above, align=center]{Deposits $\leftrightarrow$ Loans\\(Corporate, Mortgage,\\ Consumption, Trade Finance)} (hh);
\draw[bidir] (cbanks) -- node[label, above, align=center]{Deposits $\leftrightarrow$ Loans\\(Investment, Working Cap,\\ Trade, FX)} (firms);
\draw[bidir] (cbanks) -- node[label, left, align=center]{T-bills/Bonds\\(Primary/Secondary Mkt)} (gov);

\draw[bidir] (mfis) -- node[label, above, align=center]{Shares/Deposits $\leftrightarrow$ Micro-loans\\(Solidarity, Individual,\\ Agricultural, Housing)} (hh);
\draw[bidir] (mfis) -- node[label, above, align=center]{Micro/SME Loans\\(Equipment, Stock,\\ Value-chain finance)} (firms);

\draw[bidir] (other) -- node[label, left, align=center]{Postal Savings,\\ Mobile Money,\\ Remittances} (hh);
\draw[bidir] (other) -- node[label, right, align=center]{Insurance, Leasing,\\ Fintech Credit} (firms);

% Inter-bank / Inter-MFI
\draw[dasharrow] (cbanks) -- node[label, above, align=center]{Interbank Market\\(BEA, Repo)} (mfis);
\draw[dasharrow] (cbanks) -- node[label, right] {Correspondent Banking} (other);

% External
\node[client, draw=popgold, fill=popgold!10, align=center] (ext) [above=of beac, yshift=1.5cm]{External Partners\\ \footnotesize IMF, World Bank, AfDB,\\ BCEAO, International Markets};
\draw[arrow] (ext) -- node[label, left, align=center]{Policy Support,\\ Reserves, Technical Asst} (beac);

% Mobile Money overlay (dotted)
\draw[dashed, thick, poppink] (hh) -- node[label, left, text=poppink, align=center]{Mobile Money\\ (MoMo, Orange Money)} (mfis);
\draw[dashed, thick, poppink] (hh) -- (cbanks);
\draw[dashed, thick, poppink] (hh) -- (other);
\end{tikzpicture}
\legende{The Cameroonian Financial System: An Integrated Map. BEAC sits at the centre, steering monetary policy and supervising intermediaries. Commercial banks dominate corporate finance; MFIs drive inclusion. Households and firms access multiple channels. Mobile money (pink dashed) cuts across all layers.}
\end{popfigure}

\begin{aretenir}[Synthesis: Key Linkages for the GCE Exam]
\begin{enumerate}
    \item \textbf{Monetary Policy Transmission:} BEAC raises TIAO $\rightarrow$ Commercial banks' refinancing cost $\uparrow$ $\rightarrow$ Lending rates $\uparrow$ $\rightarrow$ Credit to economy $\downarrow$ $\rightarrow$ Inflation $\downarrow$ (but also investment $\downarrow$). MFIs often lag in transmission due to fixed-rate methodologies.
    \item \textbf{Financial Inclusion Gap:} Commercial banks require collateral/history $\rightarrow$ exclude 60\%+ adults. MFIs fill gap with \emph{character-based} lending. \textbf{Mobile money} is the new bridge (digital savings/credit scoring).
    \item \textbf{Systemic Risk:} A crisis in the cocoa sector $\rightarrow$ Farmer defaults at MFI $\rightarrow$ MFI liquidity crunch $\rightarrow$ MFI draws on BEAC refinancing or commercial bank line $\rightarrow$ Contagion risk. \textbf{COBAC supervision} of both banks and MFIs is the firewall.
    \item \textbf{Money Creation:} Only \textbf{commercial banks} create \emph{scriptural money} (demand deposits) via lending. MFIs create \emph{credit} but their liabilities (shares/deposits) are not fully transactional money (no cheque book for COOPEC shares). BEAC creates \emph{base money} (fiduciary + reserves).
    \item \textbf{Exam Essay Structure:} If asked ``Describe the Cameroonian financial system'', structure: \textbf{1. Apex (BEAC/COBAC)} $\rightarrow$ \textbf{2. Deposit-taking Intermediaries (Banks vs MFIs)} $\rightarrow$ \textbf{3. Non-bank Intermediaries} $\rightarrow$ \textbf{4. Markets (Money, Bond, FX)} $\rightarrow$ \textbf{5. Payment Systems} $\rightarrow$ \textbf{6. Challenges/Reforms (Inclusion, Digital, AfCFTA)}.
\end{enumerate}
\end{aretenir}

\begin{exercice}[GCE-Style Practice Questions]
\begin{enumerate}
    \item \textbf{Case Study:} Mr. Atanga, a cocoa farmer in Mamfe, applies for 500\,000 FCFA at his COOPEC for fertiliser. He has 50\,000 FCFA savings, no land title, but a salaried brother as guarantor. The loan officer approves without a farm visit. The cocoa price crashes 30\%. Mr. Atanga defaults. \textbf{Identify three appraisal failures} and \textbf{propose two policy changes} for the COOPEC.
    \item \textbf{Essay:} ``The Central Bank is the pivot of the financial system.'' Discuss this statement with reference to BEAC's roles in monetary policy, banking supervision, and financial stability in the CEMAC zone.
    \item \textbf{Quantitative:} A Credit Union has total loans of 500 million FCFA. Loans overdue > 30 days = 45 million FCFA. Loans overdue > 90 days = 20 million FCFA. Calculate \textbf{PAR 30} and \textbf{PAR 90}. If the provisioning policy is 25\% for 30--89 days and 100\% for 90+ days, calculate the \textbf{loan loss provision}. Comment on the health of the portfolio.
    \item \textbf{Short Answer:} Distinguish between \textbf{character-based lending} (microfinance) and \textbf{collateral-based lending} (traditional banking). Give one advantage and one limitation of each in the Cameroonian context.
\end{enumerate}
\end{exercice}

\begin{corrige}[Exercise 3: PAR \& Provisioning Calculation]
\textbf{Data:} Gross Loan Portfolio (GLP) = 500 million FCFA. \\
PAR 30 = 45 million. PAR 90 = 20 million.

\textbf{Calculations:}
\begin{align*}
\text{PAR 30 Ratio} &= \frac{45}{500} \times 100 = \mathbf{9\%} \\
\text{PAR 90 Ratio} &= \frac{20}{500} \times 100 = \mathbf{4\%}
\end{align*}

\textbf{Provisioning:}
\begin{itemize}
    \item Loans 30--89 days overdue = PAR 30 $-$ PAR 90 = 45 $-$ 20 = 25 million.
    \item Provision on 30--89 days = 25\% $\times$ 25 = 6.25 million.
    \item Provision on 90+ days = 100\% $\times$ 20 = 20 million.
    \item \textbf{Total Loan Loss Provision} = 6.25 + 20 = \textbf{26.25 million FCFA}.
\end{itemize}

\textbf{Comment:} PAR 30 at 9\% is \textbf{above the COBAC prudential threshold} (typically $\le$ 5--10\% depending on MFI category; generally $>$ 10\% is critical). PAR 90 at 4\% signals a \textbf{significant stock of chronic bad debts}. The provision coverage (26.25 / 45 = 58\%) is partial. The MFI must intensify recovery on the 90+ bucket and review its appraisal standards to stem new NPL formation.
\end{corrige}

\begin{savaistu}[Did You Know? The ``Caisse d'Épargne'' Legacy]
Cameroon's first savings institution was the \textbf{Caisse d'Épargne du Cameroun}, created in 1928 under French mandate. It was the ancestor of today's \textbf{CAMEC (Caisse d'Épargne et de Crédit)} and the postal savings network (CAMPOST). The modern \textbf{COOPEC movement} took off in the 1960s–70s inspired by the \emph{Caisses Populaires} of Quebec (Desjardins model) and German \emph{Genossenschaftsbanken} (Raiffeisen). The first COOPEC, \textbf{COOPEC ``La Réussite''}, was founded in 1963 in Njinikom (Northwest) by Dutch missionaries. Today, the \emph{Réseau des Caisses Populaires du Cameroun (RECAP)} and \emph{CamCCUL} are among Africa's largest MFI networks, serving over 2 million members — a testament to the power of \textbf{member-owned, character-based finance}.
\end{savaistu}

\begin{aretenir}[GCE Advanced Level: Final Exam Strategy for This Chapter]
\begin{itemize}
    \item \textbf{Definitions} (2--3 marks): Learn \emph{verbatim}: Money, Commercial Bank Balance Sheet items, Money Multiplier, TIAO, Required Reserves, Credit Union, Loan Appraisal, 5 Cs, PAR 30/90.
    \item \textbf{Mechanisms} (5--7 marks): \textbf{Credit Creation} step-by-step (T-account style). \textbf{Monetary Policy Transmission} chain (TIAO $\rightarrow$ Rates $\rightarrow$ Credit $\rightarrow$ AD $\rightarrow$ P/Y). \textbf{BEAC Tools} (Open Market Ops, Standing Facilities, Reserve Req).
    \item \textbf{Calculations} (4--6 marks): Money Multiplier ($1/rr$), Excess Reserves, PAR Ratios, Provisioning, Instalment-to-Income Ratio, Effective Interest Rate (flat vs declining).
    \item \textbf{Essay/Case Study} (10--15 marks): Structure: \textbf{Introduction (Definitions + Context)} $\rightarrow$ \textbf{Body: 3--4 argued points with Cameroonian examples (Cocoa, BEAC, COOPEC, MoMo)} $\rightarrow$ \textbf{Evaluation (Limits, Challenges, Policy Recommendations)} $\rightarrow$ \textbf{Conclusion}.
    \item \textbf{Keywords to sprinkle:} \cle{CEMAC}, \cle{COBAC}, \cle{CFA Franc}, \cle{Fixed Exchange Rate}, \cle{Financial Inclusion}, \cle{Non-Performing Loans}, \cle{Mobile Money}, \cle{AfCFTA}, \cle{Monetary Policy Committee (MPC)}.
\end{itemize}
\end{aretenir}

\textbf{End of Chapter ECO04: Money and Banking (Continue).} You have now covered the full syllabus from commercial bank balance sheets to the appraisal desk of a village Credit Union, with the Central Bank orchestrating in between. Master the diagrams, practise the calculations, and always ground your answers in the \textbf{Cameroonian reality}. Good luck for the GCE!

\newpage

\coursec{Integration activity}

\courssub{Aims of the activity}

This integration activity brings together in a single, coherent case study all the notions studied in this chapter: the \cle{balance sheet} of a commercial bank, the mechanism of \cle{credit creation} and the \cle{credit multiplier}, the role of the \cle{central bank} and the \cle{transmission mechanism} of monetary policy, the place of \cle{microfinance} in Cameroon, and the \cle{appraisal of a loan file} using the 5~Cs method.

By the end of the activity, you should be able to:

\begin{itemize}
\item construct and interpret the simplified T-account of a commercial bank after a new deposit, distinguishing \cle{required reserves} from \cle{excess reserves};
\item calculate the maximum credit creation of the banking system from a given reserve requirement, using the credit multiplier $k = \dfrac{1}{r}$;
\item explain, step by step, how an increase in the key rate of the BEAC is transmitted to households, firms and the general price level in Cameroon;
\item appraise a real loan file at a credit union using the 5~Cs (Character, Capacity, Capital, Collateral, Conditions) and defend an accept/reject decision with arguments;
\item produce a short, well-structured written report combining calculation, economic analysis and a justified decision --- exactly the skill required in the GCE Advanced Level essay and data-response papers.
\end{itemize}

\begin{aretenir}[Key formulae used in this activity]
\begin{itemize}
\item Required reserves $=$ reserve ratio $\times$ deposits, i.e. $R = r \times D$.
\item Excess (loanable) reserves $=$ total new reserves $-$ required reserves.
\item Credit multiplier: $k = \dfrac{1}{r}$, where $r$ is the reserve requirement expressed as a \emph{decimal}.
\item Maximum total deposits $= k \times$ initial new deposit; maximum new credit $=$ maximum total deposits $-$ initial deposit.
\item Flat-rate monthly instalment $= \dfrac{\text{principal}}{\text{number of months}} + (\text{principal} \times \text{monthly flat rate})$.
\end{itemize}
\end{aretenir}

\courssub{Situation}

\textbf{From Mama Put's Deposit to a Taxi Loan.}\par

Madame Eyenga runs a successful \emph{mama put} (cooked-food) restaurant near the Mokolo market in Yaoundé. After several good months, she has accumulated \textbf{500,000~FCFA} in cash, which she keeps at home. Worried about theft, she deposits the whole amount on 3~January in a current account at the Mokolo branch of \textbf{Afriland First Bank}.

The Bank of Central African States (\textbf{BEAC}), the central bank common to the six CEMAC countries (Cameroon, Chad, Central African Republic, Congo, Equatorial Guinea and Gabon), imposes on commercial banks a \textbf{reserve requirement of 13\,\%}: every bank must keep 13\,\% of its deposits as reserves (at the central bank or in vault cash) and may lend out the rest. Assume that:

\begin{itemize}
\item banks lend out \emph{all} their excess reserves (no voluntary idle reserves);
\item the public holds no cash: every loan is spent and re-deposited in full in the banking system;
\item the banking system is initially in equilibrium.
\end{itemize}

Three months later, worried about inflationary pressures in the CEMAC zone (imported food and fuel prices are rising), the \textbf{Monetary Policy Committee of the BEAC raises its key policy rate (the TIAO, the interest rate on its main refinancing operations) by 50 basis points}, from 5.00\,\% to 5.50\,\%. Recall that \textbf{one basis point} is one hundredth of a percentage point: 50 basis points $= 0.50$ percentage point.

Meanwhile, Ndifor, a 27-year-old motorcycle-taxi (\emph{bendskin}) rider in Bamenda, wants to buy a second motorcycle costing \textbf{800,000~FCFA} to rent to another rider. He applied to a commercial bank but was refused: he has no payslip, no land title and no formal accounts. He therefore applies for an \textbf{800,000~FCFA loan over 24 months at 1.5\,\% per month (flat) at a CAMCCUL-affiliated credit union}. His loan file contains:

\begin{center}
\begin{tabularx}{\linewidth}{|l|X|}
\hline
\rowcolor{popblueL} \textbf{Item} & \textbf{Details} \\
\hline
Occupation & Motorcycle-taxi rider, 4 years' experience, known at his stage \\
\hline
Daily net income & About 6,000~FCFA/day, worked 26 days/month ($\approx$ 156,000~FCFA/month) \\
\hline
Monthly expenses & Food, rent, fuel, maintenance: $\approx$ 100,000~FCFA \\
\hline
Savings history & Member of the credit union for 3 years; regular savings of 120,000~FCFA; one previous loan of 300,000~FCFA fully repaid on time \\
\hline
Guarantors & Two members of the union, each with savings above 200,000~FCFA \\
\hline
Collateral & No land title; offers the motorcycle itself (logbook deposited) \\
\hline
Business plan & Rent the new bike at 2,000~FCFA/day; expected $\approx$ 52,000~FCFA/month extra income \\
\hline
Monthly instalment & Principal $+$ interest $\approx$ 45,333~FCFA/month for 24 months \\
\hline
\end{tabularx}
\end{center}

Each group of 3--4 students acts successively as (i) the bank's accounts officer, (ii) an economist at the BEAC, and (iii) the credit committee of the credit union. Each group submits a \textbf{one-page report} combining the calculation, the policy analysis and the appraisal decision.

\courssub{Tasks}

\begin{enumerate}
\item \textbf{The bank's T-account.} Record Madame Eyenga's deposit of 500,000~FCFA in the simplified balance sheet (T-account) of Afriland First Bank, distinguishing required reserves from excess reserves. How much can the bank lend immediately?
\item \textbf{Credit creation.} Assuming the loan is spent and re-deposited in the banking system, calculate the credit multiplier $k = \dfrac{1}{r}$ with $r = 13\,\%$. Deduce (a) the total amount of deposits the banking system can ultimately support, and (b) the total new credit created. Show the first three rounds of the process in a small table.
\item \textbf{Limits of the multiplier.} Give two reasons why, in the real Cameroonian economy, actual credit creation is smaller than the theoretical maximum.
\item \textbf{Monetary policy transmission.} The BEAC raises its key rate by 50 basis points. Explain the transmission mechanism step by step: effect on commercial banks' refinancing cost, on lending rates, on Madame Eyenga (saver), on borrowers like Ndifor, on investment and consumption, and finally on inflation in Cameroon. State one reason why this transmission may be weak in Cameroon.
\item \textbf{Loan appraisal.} Using the 5~Cs, appraise Ndifor's loan file. Compute his monthly repayment capacity (income minus expenses, before and after the new business income) and compare it with the monthly instalment of 45,333~FCFA.
\item \textbf{Decision.} As the credit committee, decide: accept, accept with conditions, or reject. Justify your decision in at most ten lines, proposing any guarantees or conditions you judge necessary.
\item \textbf{Synthesis.} In five lines, explain how this case study shows the complementarity between commercial banks, the central bank and microfinance institutions in financing the Cameroonian economy.
\end{enumerate}

\courssub{Model answer}

\textbf{1. The bank's T-account.}\par
The deposit creates a liability (the current account owed to Madame Eyenga) and an asset (cash/reserves held by the bank). Required reserves are
\[ R = r \times D = 0.13 \times 500{,}000 = 65{,}000\ \text{FCFA}; \]
excess reserves are $500{,}000 - 65{,}000 = 435{,}000$~FCFA, which the bank can lend immediately.

\begin{center}
\begin{tabular}{|l r|l r|}
\hline
\rowcolor{popblueL} \multicolumn{2}{|c|}{\textbf{Assets (FCFA)}} & \multicolumn{2}{c|}{\textbf{Liabilities (FCFA)}} \\
\hline
Required reserves & 65,000 & Demand deposit (Mme Eyenga) & 500,000 \\
Excess reserves / loanable & 435,000 & & \\
\hline
\textbf{Total} & 500,000 & \textbf{Total} & 500,000 \\
\hline
\end{tabular}
\end{center}

Note that the balance sheet always balances: total assets $=$ total liabilities $= 500{,}000$~FCFA.

\textbf{2. Credit creation.}\par
The credit multiplier is
\[ k = \frac{1}{r} = \frac{1}{0.13} \approx 7.69. \]
(a) Total deposits the system can ultimately support:
\[ 500{,}000 \times 7.69 \approx 3{,}846{,}154\ \text{FCFA}. \]
(b) Total new credit created:
\[ 3{,}846{,}154 - 500{,}000 \approx 3{,}346{,}154\ \text{FCFA}, \]
equivalently $435{,}000 \times 7.69 \approx 3{,}346{,}154$~FCFA.

\begin{center}
\begin{tabular}{|l|r|r|r|}
\hline
\rowcolor{popblueL} \textbf{Round} & \textbf{New deposit} & \textbf{Required reserves} & \textbf{New loan} \\
\hline
1 (Afriland) & 500,000 & 65,000 & 435,000 \\
2 & 435,000 & 56,550 & 378,450 \\
3 & 378,450 & 49,199 & 329,252 \\
\hline
$\cdots$ Total & 3,846,154 & 500,000 & 3,346,154 \\
\hline
\end{tabular}
\end{center}

Each round, the new loan of one bank becomes the new deposit of the next; 13\,\% is immobilised as reserves and 87\,\% is re-lent. The process converges because each round is smaller than the previous one (it is a geometric progression with ratio $1 - r = 0.87$).

\textbf{3. Limits of the multiplier.}\par
(a) \emph{Cash leakages}: in Cameroon many agents keep part of their money in cash (large informal sector, cash-based retail trade), so loans are not fully re-deposited and each round leaks out of the banking system. (b) \emph{Prudence of banks and lack of creditworthy borrowers}: banks voluntarily hold excess reserves as a precaution, or cannot find solvent borrowers with acceptable guarantees --- especially small informal businesses without land titles or payslips, like Ndifor. Both leakages reduce the effective multiplier well below 7.69.

\textbf{4. Transmission of the 50-basis-point increase.}\par
When the BEAC raises its key rate, commercial banks that refinance at the central bank pay more for the funds they borrow; their cost of funds rises, so they increase their lending rates and may tighten credit conditions (stricter files, more collateral demanded). For \textbf{Madame Eyenga} (a net saver), deposit rates may rise slightly --- a gain. For \textbf{borrowers like Ndifor}, bank credit becomes dearer and harder to obtain, pushing them further towards microfinance. Firms face costlier investment credit, so investment falls; households reduce consumption financed by credit. Aggregate demand contracts, which eases demand-pull pressure on prices and helps bring \textbf{inflation} back towards the CEMAC convergence criterion (an inflation rate of at most 3\,\%).

\begin{popfigure}
\begin{tikzpicture}[
box/.style={draw=popblue, fill=popblueL, rounded corners=2pt, align=center, text width=3.1cm, minimum height=1.0cm, font=\small},
arr/.style={-{Stealth[length=2.5mm]}, thick, popdark}]
\node[box] (a) at (0,0) {BEAC raises key rate (TIAO)};
\node[box] (b) at (4.6,0) {Banks' refinancing cost rises};
\node[box] (c) at (9.2,0) {Lending rates rise, credit tightens};
\node[box] (d) at (9.2,-2.2) {Investment and consumption fall};
\node[box] (e) at (4.6,-2.2) {Aggregate demand contracts};
\node[box] (f) at (0,-2.2) {Inflationary pressure eases};
\draw[arr] (a) -- (b);
\draw[arr] (b) -- (c);
\draw[arr] (c) -- (d);
\draw[arr] (d) -- (e);
\draw[arr] (e) -- (f);
\end{tikzpicture}
\legende{The transmission mechanism of a restrictive monetary policy (interest-rate channel).}
\end{popfigure}

\emph{Weakness of transmission in Cameroon}: a large share of economic activity is informal and financed outside the banking system (tontines, cash holdings, microfinance), and bank credit to the private sector remains limited, so a change in the key rate directly affects only the banked part of the economy. Moreover, when inflation is imported (fuel, food), a demand-side instrument has limited power over it.

\textbf{5. Appraisal with the 5~Cs.}\par
First, check the instalment: with a flat rate of 1.5\,\% per month on 800,000~FCFA over 24 months,
\[ \text{monthly instalment} = \frac{800{,}000}{24} + (800{,}000 \times 0.015) = 33{,}333 + 12{,}000 = 45{,}333\ \text{FCFA}, \]
which matches the figure in the file. Then appraise:
\begin{itemize}
\item \textbf{Character}: excellent --- 3 years' membership, regular savings, previous loan of 300,000~FCFA repaid on time.
\item \textbf{Capacity}: current disposable surplus $\approx 156{,}000 - 100{,}000 = 56{,}000$~FCFA/month, which already covers the instalment of 45,333~FCFA (a margin of about 10,667~FCFA); with the rental income ($\approx 52{,}000$~FCFA/month), the surplus rises to about 108,000~FCFA, giving a comfortable safety margin even if the second bike earns nothing for some months.
\item \textbf{Capital}: savings of 120,000~FCFA, i.e. 15\,\% of the loan --- an acceptable personal contribution showing commitment.
\item \textbf{Collateral}: no land title, but the motorcycle's logbook is deposited and two solvent guarantors co-sign --- typical microfinance substitute guarantees.
\item \textbf{Conditions}: the bendskin business is in strong demand in Bamenda; risks exist (accidents, breakdowns, insecurity), which justify insurance and maintenance provisions.
\end{itemize}

\textbf{6. Decision.}\par
\emph{Accept with conditions.} The file shows good character, sufficient repayment capacity even without the new income, and acceptable guarantees. Conditions: (i) comprehensive insurance on the motorcycle; (ii) deposit of the logbook at the union; (iii) the two guarantors sign the bond; (iv) monthly instalments deducted first from his savings account; (v) disbursement directly to the motorcycle dealer to ensure the loan is used as planned.

\textbf{7. Synthesis.}\par
The case shows a complete financial circuit: the trader's deposit feeds a commercial bank, which multiplies credit under the reserve rules fixed by the BEAC; the central bank steers the cost and quantity of money through its key rate; and microfinance institutions like CAMCCUL credit unions reach clients excluded from banks, using savings history and solidarity guarantees instead of classical collateral. Together, these three levels channel savings towards investment and support growth in Cameroon.

\begin{attention}[Common mistakes to avoid]
Do not confuse the \emph{first} loan (435,000~FCFA) with the \emph{total} credit created ($\approx$ 3,346,154~FCFA). In the multiplier formula, use $r$ as a decimal (0.13, not 13). Remember that 50 basis points means 0.50 percentage point, not 50\,\%. Finally, in the appraisal, never decide on collateral alone: in microfinance, \cle{character} and \cle{capacity} come first.
\end{attention}

\begin{exercice}[For consolidation]
Suppose the BEAC lowers the reserve requirement from 13\,\% to 10\,\% and Madame Eyenga's deposit remains 500,000~FCFA, all other assumptions unchanged.
\begin{enumerate}
\item Recompute the credit multiplier, the maximum total deposits and the maximum new credit.
\item Is this measure expansionary or restrictive? Explain its likely effect on credit and economic activity in Cameroon.
\item A credit union lends 600,000~FCFA over 12 months at 2\,\% per month (flat). Compute the monthly instalment and the total interest paid.
\end{enumerate}
\end{exercice}

\begin{corrige}[Exercise: For consolidation]
\begin{enumerate}
\item $k = \dfrac{1}{0.10} = 10$. Maximum total deposits $= 500{,}000 \times 10 = 5{,}000{,}000$~FCFA; maximum new credit $= 5{,}000{,}000 - 500{,}000 = 4{,}500{,}000$~FCFA.
\item Expansionary: lowering the reserve requirement frees excess reserves, raises the multiplier and allows banks to grant more credit, which stimulates investment, consumption and growth (but may fuel inflation if demand outpaces supply).
\item Monthly instalment $= \dfrac{600{,}000}{12} + (600{,}000 \times 0.02) = 50{,}000 + 12{,}000 = 62{,}000$~FCFA. Total interest $= 12{,}000 \times 12 = 144{,}000$~FCFA.
\end{enumerate}
\end{corrige}

\newpage

\coursec{Methods and worked examples}

\coursec{Money and Banking (Continued)}

\courssub{The Balance Sheet of a Commercial Bank}

\begin{definition}[Commercial Bank Balance Sheet]
A \cle{balance sheet} is a snapshot of a bank's financial position at a given date. It records \cle{assets} (uses of funds) on the left and \cle{liabilities} (sources of funds) on the right. The fundamental identity holds:
\[
\text{Total Assets} \equiv \text{Total Liabilities} + \text{Shareholders' Equity}.
\]
Assets are ordered by decreasing liquidity; liabilities by increasing urgency of repayment.
\end{definition}

\begin{propriete}[Main Items in a Cameroonian Commercial Bank's Balance Sheet]
\begin{itemize}
\item \textbf{Assets (Emplois)}: Cash in vault (\emph{valeurs en caisse}); Reserves at BEAC (\emph{réserves obligatoires et excédentaires}); Treasury bills and government securities (\emph{titres publics}); Loans and overdrafts to customers (\emph{crédits à la clientèle}); Bills discounted (\emph{effets escomptés}); Fixed assets (\emph{immobilisations corporelles et incorporelles}).
\item \textbf{Liabilities (Ressources)}: Sight deposits (\emph{dépôts à vue}); Savings deposits (\emph{dépôts d'épargne}); Time deposits (\emph{dépôts à terme}); Borrowings from BEAC and other banks (\emph{emprunts auprès de la BEAC et banques}); Share capital (\emph{capital social}); Reserves and retained earnings (\emph{réserves et reports à nouveau}); Provisions.
\end{itemize}
\end{propriete}

\begin{methode}[Reading and Analysing a Bank Balance Sheet]
\begin{enumerate}
\item Draw the two-column framework: \textbf{Assets} (left) and \textbf{Liabilities + Equity} (right).
\item Classify every item: deposits are liabilities (the bank owes them); loans, securities, cash are assets (the bank owns claims).
\item Order assets from most liquid (cash, BEAC reserves) to least liquid (fixed assets, long-term loans). Order liabilities from most stable (equity, time deposits) to most repayable on demand (sight deposits).
\item Compute key prudential ratios:
\begin{itemize}
\item \cle{Liquidity ratio} (COBAC standard): $\dfrac{\text{Liquid assets (cash + BEAC reserves + Treasury bills)}}{\text{Short-term liabilities (sight + savings + time deposits)}} \ge 100\%$ (simplified).
\item \cle{Solvency / Capital adequacy ratio} (Basel/COBAC): $\dfrac{\text{Regulatory capital (Tier 1 + Tier 2)}}{\text{Risk-weighted assets}} \ge 8\%$ (or 10\% in CEMAC).
\end{itemize}
\item Interpret: high loans/assets $\rightarrow$ profitability but liquidity risk; high cash/assets $\rightarrow$ safety but low return.
\end{enumerate}
\end{methode}

\begin{exemplebox}[T-Account Representation of a Deposit]
When a customer deposits 1\,000\,000 FCFA in cash at a bank:
\begin{center}
\begin{tabular}{|l|r|l|r|}
\hline
\multicolumn{2}{|c|}{\textbf{Assets}} & \multicolumn{2}{c|}{\textbf{Liabilities}} \\
\hline
Cash in vault & +1\,000\,000 & Sight deposits & +1\,000\,000 \\
\hline
\end{tabular}
\end{center}
The bank's liquidity increases; its obligation to the customer increases by the same amount.
\end{exemplebox}

\begin{exoresolu}[Balance Sheet of Sanaga Commercial Bank]
The following items (in millions of FCFA) belong to Sanaga Commercial Bank, Douala: sight deposits 40\,000; cash in vault 3\,000; loans to customers 45\,000; savings deposits 25\,000; reserves at BEAC 7\,000; shareholders' capital 10\,000; treasury bills held 8\,000; time deposits 15\,000; buildings and equipment 5\,000; reserves (retained profits) 12\,000. Present the balance sheet in regulatory order and comment on liquidity and solvency.
\tcblower
\textbf{Step 1 — Classify and order.}
\begin{itemize}
\item \textbf{Assets (by liquidity)}: Cash in vault (3\,000); Reserves at BEAC (7\,000); Treasury bills (8\,000); Loans to customers (45\,000); Buildings \& equipment (5\,000).
\item \textbf{Liabilities + Equity (by stability)}: Share capital (10\,000); Reserves/retained profits (12\,000); Time deposits (15\,000); Savings deposits (25\,000); Sight deposits (40\,000).
\end{itemize}

\textbf{Step 2 — Present the balance sheet.}
\begin{center}
\begin{tabularx}{\linewidth}{|X|r|X|r|}
\hline
\rowcolor{popblueL} \textbf{Assets (Emplois)} & \textbf{M FCFA} & \textbf{Liabilities + Equity (Ressources)} & \textbf{M FCFA} \\
\hline
Cash in vault (Caisse) & 3\,000 & Share capital (Capital social) & 10\,000 \\
Reserves at BEAC (Réserves BEAC) & 7\,000 & Reserves \& retained profits (Réserves \& RAN) & 12\,000 \\
Treasury bills (Titres publics) & 8\,000 & Time deposits (Dépôts à terme) & 15\,000 \\
Loans to customers (Crédits clientèle) & 45\,000 & Savings deposits (Dépôts d'épargne) & 25\,000 \\
Buildings \& equipment (Immobilisations) & 5\,000 & Sight deposits (Dépôts à vue) & 40\,000 \\
\hline
\textbf{Total Assets} & \textbf{68\,000} & \textbf{Total Liabilities + Equity} & \textbf{68\,000} \\
\hline
\end{tabularx}
\end{center}

\textbf{Step 3 — Compute ratios and interpret.}
\begin{itemize}
\item \textbf{Liquidity ratio (simplified)}: Liquid assets = 3\,000 + 7\,000 + 8\,000 = 18\,000. Short-term liabilities = 40\,000 + 25\,000 + 15\,000 = 80\,000. Ratio = 18\,000 / 80\,000 = 22.5\%. \emph{Below the COBAC 100\% norm} — the bank relies on the interbank market or BEAC refinancing to meet withdrawals.
\item \textbf{Loans / Total assets} = 45\,000 / 68\,000 $\approx$ 66.2\%. High transformation ratio: profitable but vulnerable to a run.
\item \textbf{Equity / Total assets} = (10\,000 + 12\,000) / 68\,000 $\approx$ 32.4\%. Strong capital base, well above the 8--10\% risk-weighted requirement.
\end{itemize}
\textbf{Conclusion:} The bank is solvent but structurally illiquid — a typical profile for a Cameroonian commercial bank engaged in maturity transformation.
\end{exoresolu}

\courssub{Credit Creation by Commercial Banks}

\begin{definition}[Credit Creation]
\cle{Credit creation} is the process by which the banking system, starting from an initial deposit, grants loans that become new deposits in the system, thereby multiplying the money supply. It rests on \cle{fractional reserve banking}: banks keep only a fraction $r$ (reserve ratio) of deposits as reserves and lend out the rest.
\end{definition}

\begin{propriete}[The Money (Credit) Multiplier]
Under the assumptions: (i) no cash drain (public holds no currency), (ii) banks hold zero excess reserves, (iii) all loans are re-deposited in the banking system, (iv) single reserve ratio $r$ for all deposits,
\[
\text{Money multiplier } k = \frac{1}{r}, \qquad
\text{Total deposits created } D = k \times D_0 = \frac{D_0}{r}, \qquad
\text{Total credit created } C = D - D_0 = D_0\left(\frac{1}{r} - 1\right).
\]
With a \cle{cash-drain ratio} $c$ (fraction of loans the public keeps as cash) and/or \cle{excess reserve ratio} $e$, the effective multiplier becomes
\[
k_{\text{eff}} = \frac{1}{r + c + e}.
\]
\end{propriete}

\begin{methode}[Step-by-Step Credit Creation Calculation]
\begin{enumerate}
\item Identify the \cle{initial new deposit} $D_0$ (or initial excess reserves) and the \cle{required reserve ratio} $r$ set by BEAC.
\item State the assumptions (no leakages, or specify $c$ and $e$).
\item Compute the multiplier $k = 1/(r+c+e)$.
\item Compute total deposits $D = k D_0$ and total credit created $C = D - D_0$.
\item (Optional) Show the round-by-round process in a table to illustrate the geometric series.
\end{enumerate}
\end{methode}

\begin{experience}[Simulating Credit Creation Round by Round]
Consider $D_0 = 500$ million FCFA, $r = 10\%$, no leakages.
\begin{center}
\begin{tabular}{|c|c|c|c|c|}
\hline
\rowcolor{popblueL} \textbf{Round} & \textbf{New Deposit} & \textbf{Reserves (10\%)} & \textbf{New Loan} & \textbf{Cumulative Deposits} \\
\hline
Initial & 500 & 50 & 450 & 500 \\
1 & 450 & 45 & 405 & 950 \\
2 & 405 & 40.5 & 364.5 & 1\,355 \\
3 & 364.5 & 36.45 & 328.05 & 1\,719.5 \\
$\vdots$ & $\vdots$ & $\vdots$ & $\vdots$ & $\vdots$ \\
\hline
\textbf{Limit} & \textbf{5\,000} & \textbf{500} & \textbf{4\,500} & \textbf{5\,000} \\
\hline
\end{tabular}
\end{center}
The geometric series $500(1 + 0.9 + 0.9^2 + \cdots) = 500/(1-0.9) = 5\,000$.
\end{experience}

\begin{exoresolu}[Credit Creation in the Cameroonian Banking System]
An initial new deposit of 500 million FCFA is made at a commercial bank in Yaoundé. The required reserve ratio fixed by BEAC is 10\%.
\begin{enumerate}
\item Calculate the money multiplier, total deposits and total credit created assuming no leakages.
\item Repeat with a cash-drain ratio of 5\% (public keeps 5\% of each loan in cash) and zero excess reserves.
\item Explain why the actual credit creation in Cameroon is often lower than the theoretical maximum.
\end{enumerate}
\tcblower
\textbf{Part 1 — No leakages.}
\[
k = \frac{1}{0.10} = 10, \quad
D = 10 \times 500 = 5\,000 \text{ million FCFA}, \quad
C = 5\,000 - 500 = 4\,500 \text{ million FCFA}.
\]

\textbf{Part 2 — With cash drain $c = 0.05$.}
\[
k' = \frac{1}{0.10 + 0.05} = \frac{1}{0.15} \approx 6.667, \quad
D' \approx 6.667 \times 500 = 3\,333 \text{ million FCFA}, \quad
C' \approx 2\,833 \text{ million FCFA}.
\]

\textbf{Part 3 — Why actual creation is lower in Cameroon.}
\begin{itemize}
\item \textbf{High cash-drain ratio:} a large informal sector means much currency circulates outside banks ($c$ high).
\item \textbf{Excess reserves:} banks often hold reserves above the minimum due to uncertainty, poor interbank market, or lack of creditworthy borrowers ($e > 0$).
\item \textbf{Non-bank financial intermediaries:} microfinance institutions, tontines (\emph{njangui}) capture savings that do not enter the commercial bank multiplier.
\item \textbf{BEAC prudential limits:} concentration risk limits, sectoral credit ceilings, and the COBAC liquidity ratio constrain lending.
\item \textbf{Structural factors:} weak collateral registries, high NPLs (non-performing loans), and information asymmetry ration credit.
\end{itemize}
\end{exoresolu}

\begin{attention}[Common Mistakes on Credit Creation]
\begin{itemize}
\item Confusing \emph{total deposits created} ($D$) with \emph{credit created} ($C = D - D_0$).
\item Forgetting that the multiplier applies to the \emph{banking system as a whole}, not to a single bank (a single bank can only lend its excess reserves).
\item Using the reserve ratio as a percentage (10) instead of a decimal (0.10) in the formula.
\item Ignoring leakages: in Cameroon, the cash-drain ratio is historically high (often 30--40\%), drastically reducing the real multiplier.
\end{itemize}
\end{attention}

\courssub{The Central Bank: BEAC and Monetary Policy in CEMAC}

\begin{definition}[Central Bank]
A \cle{central bank} is the apex monetary institution entrusted with: (1) issuing legal tender; (2) acting as banker to commercial banks (lender of last resort); (3) acting as banker and fiscal agent to the State; (4) managing foreign exchange reserves; (5) supervising the financial system; (6) conducting monetary policy to achieve price stability and support growth.
\end{definition}

\begin{propriete}[BEAC — Banque des États de l'Afrique Centrale]
\begin{itemize}
\item \textbf{Status:} Supranational central bank for the 6 CEMAC member States (Cameroon, Central African Republic, Chad, Congo, Equatorial Guinea, Gabon).
\item \textbf{Headquarters:} Yaoundé, Cameroon. National Directorates in each member State.
\item \textbf{Currency:} FCFA (Franc de la Coopération Financière en Afrique Centrale), pegged to the euro at 1 EUR = 655.957 FCFA (fixed since 1999).
\item \textbf{Key governance:} Monetary Policy Committee (CPM) sets the key rates; COBAC (Commission Bancaire de l'Afrique Centrale) supervises banks and MFIs.
\item \textbf{Convergence criteria (surveillance):} Inflation $\le$ 3\%; Fiscal deficit $\le$ 1.5\% GDP; Public debt $\le$ 70\% GDP; Non-accumulation of arrears.
\end{itemize}
\end{propriete}

\begin{methode}[Analysing BEAC's Monetary Policy Instruments]
\begin{enumerate}
\item \textbf{Diagnose the macroeconomic situation:} inflation above 3\% $\rightarrow$ restrictive policy; recession/unemployment $\rightarrow$ expansionary policy.
\item \textbf{Choose the instrument and direction:}
\begin{center}
\begin{tabularx}{\linewidth}{|l|X|X|}
\hline
\rowcolor{popblueL} \textbf{Instrument} & \textbf{Restrictive (Anti-inflation)} & \textbf{Expansionary (Anti-recession)} \\
\hline
Key policy rate (TIAO — Taux d'Intérêt des Appels d'Offres) & \textbf{Raise} & \textbf{Lower} \\
Marginal lending facility rate (TIP — Taux d'Intérêt du Prêt marginal) & \textbf{Raise} & \textbf{Lower} \\
Reserve requirement ratio (Taux de réserves obligatoires) & \textbf{Raise} (reduces multiplier) & \textbf{Lower} \\
Open Market Operations (OMO) — BEAC securities & \textbf{Sell} (absorb liquidity) & \textbf{Buy} (inject liquidity) \\
\hline
\end{tabularx}
\end{center}
\item \textbf{Trace the transmission mechanism:}
\[
\text{Policy rate} \rightarrow \text{Interbank rate} \rightarrow \text{Bank lending rates} \rightarrow \text{Credit volume} \rightarrow \text{Investment/Consumption} \rightarrow \text{AD} \rightarrow \text{Prices/Output}.
\]
\item \textbf{Evaluate limits:} time lags (6--18 months); imported inflation (oil, food) insensitive to rates; single policy for 6 heterogeneous economies; weak interest-rate pass-through in shallow markets; fiscal dominance.
\end{enumerate}
\end{methode}

\begin{exoresolu}[BEAC Reacts to an Inflationary Shock in CEMAC]
Following a surge in global food and fuel prices, CEMAC inflation rises to 5\%, breaching the 3\% convergence criterion. The CPM decides to tighten policy.
\begin{enumerate}
\item Describe the sequence of instruments BEAC can activate.
\item Trace the transmission to a Cameroonian importer in Douala.
\item Discuss two major limits of this policy in the CEMAC context.
\end{enumerate}
\tcblower
\textbf{1. Instruments activated.}
\begin{itemize}
\item \textbf{Raise TIAO} (e.g. from 3.50\% to 4.50\%): makes weekly liquidity injections more expensive for banks.
\item \textbf{Raise TIP} (marginal lending window): raises the ceiling for overnight borrowing.
\item \textbf{Raise reserve requirement} (e.g. from 7\% to 9\%): directly cuts the credit multiplier $k = 1/r$.
\item \textbf{Open market sales} of BEAC bills: mop up excess liquidity permanently.
\end{itemize}

\textbf{2. Transmission to a Douala importer.}
Higher TIAO $\rightarrow$ banks bid less at auctions $\rightarrow$ interbank rate rises $\rightarrow$ banks' marginal cost of funds rises $\rightarrow$ prime lending rate (e.g. from 7.5\% to 9\%) $\rightarrow$ the importer's cost of financing letters of credit rises $\rightarrow$ import volumes contract $\rightarrow$ demand-pull pressure on prices eases.

\textbf{3. Limits.}
\begin{itemize}
\item \textbf{Supply-side inflation:} the shock comes from global commodity prices and poor domestic logistics (bad roads, port delays). Monetary tightening cannot lower the world price of wheat or fuel.
\item \textbf{Time lag and fiscal dominance:} the policy takes months to bite, while the government may simultaneously run large deficits financed by BEAC advances (statutory advances), injecting liquidity that offsets the tightening.
\item \textbf{One size fits six:} a rate hike suited to Gabon (oil-rich, low inflation) may deepen recession in CAR or Chad.
\end{itemize}
\end{exoresolu}

\begin{savaistu}[BEAC and the Euro Peg]
The FCFA/euro peg is guaranteed by the French Treasury under a 2019 reform of the cooperation agreements. BEAC must keep at least 50\% of its foreign reserves in an operations account at the Banque de France. In return, France guarantees unlimited convertibility. This arrangement delivers exchange-rate stability and low inflation but limits monetary sovereignty: BEAC cannot act as a true lender of last resort in foreign currency, and the fixed parity may become misaligned with fundamentals.
\end{savaistu}

\courssub{Microfinance Activities in Cameroon}

\begin{definition}[Microfinance]
\cle{Microfinance} refers to the provision of financial services — small loans (\cle{microcredit}), savings, insurance, payment services — to low-income individuals and micro-enterprises excluded from the formal banking sector due to lack of collateral, credit history, or formal employment.
\end{definition}

\begin{propriete}[Microfinance Landscape in Cameroon (COBAC Regulation)]
\begin{itemize}
\item \textbf{Category 1:} Networks of credit unions / cooperatives (e.g. \cle{CAMCCUL} — Cameroon Cooperative Credit Union League, \cle{MC$^2$} — Mutuelles Communautaires de Croissance). Member-owned, not-for-profit, mobilise savings from members only.
\item \textbf{Category 2:} Microfinance institutions (MFIs) authorised to collect savings from the public and grant credit (e.g. \cle{UCCGN}, \cle{CEC}, \cle{AFG}). For-profit or non-profit, stricter prudential norms.
\item \textbf{Category 3:} MFIs authorised to grant credit only (no public savings collection), often funded by donors or wholesale lenders.
\item \textbf{Informal sector:} \cle{Tontines} (\emph{njangui}, \emph{esusu}) — rotating savings and credit associations (ROSCAs) and accumulating savings and credit associations (ASCAs). Not supervised by COBAC but socially crucial.
\end{itemize}
\end{propriete}

\begin{methode}[Analysing the Role and Mechanism of Microfinance]
\begin{enumerate}
\item \textbf{Mechanism:} Group lending / solidarity groups $\rightarrow$ \cle{social collateral} replaces physical collateral; small, frequent repayments (weekly/bi-weekly); compulsory savings build track record; progressive lending (larger loans after successful repayment).
\item \textbf{Advantages:} Financial inclusion (women, youth, rural); poverty reduction; enterprise development; savings mobilisation; women empowerment; social cohesion.
\item \textbf{Limits/Risks:} High effective interest rates (administrative cost of small loans); over-indebtedness (multiple borrowing); weak governance/fraud in some MFIs; mission drift (consumer loans vs productive loans); limited loan sizes constrain growth; vulnerability to covariate shocks (crop failure, pandemic).
\end{enumerate}
\end{methode}

\begin{exoresolu}[Microfinance and a Tomato Trader in Bafoussam]
Mama Ngo, a tomato retailer in Bafoussam market, needs 150\,000 FCFA to buy stock. Commercial banks refuse her (no payslip, no land title). Explain how a Category-1 credit union can help her, using the specific mechanisms of microfinance.
\tcblower
\textbf{Mechanism at a CAMCCUL-affiliated credit union:}
\begin{enumerate}
\item \textbf{Membership \& savings:} Mama Ngo joins the union, buys a share (e.g. 10\,000 FCFA), and starts compulsory weekly savings of 5\,000 FCFA. After 3 months, she has 60\,000 FCFA savings — this creates a \emph{credit history} and \emph{internal collateral}.
\item \textbf{Loan application:} She applies for 150\,000 FCFA. The credit committee assesses her \emph{character} (market reputation), \emph{capacity} (daily cash flow from tomato sales), \emph{capital} (her savings), \emph{collateral} (her savings pledged + two co-signatories from her solidarity group), \emph{conditions} (working capital, 6-month term).
\item \textbf{Disbursement \& repayment:} Loan disbursed directly to her supplier (avoids diversion). Weekly instalments of ~6\,500 FCFA deducted from daily sales. Group members exert peer pressure for repayment.
\end{enumerate}

\textbf{Benefits:} (i) Access despite no formal collateral; (ii) working capital turnover increases $\rightarrow$ higher income; (iii) savings habit + financial literacy; (iv) profits reinvested in union $\rightarrow$ community development.

\textbf{Risks:} (i) Effective rate $\approx$ 24--30\% p.a. (flat rate $\approx$ 1.5--2\%/month) due to high operating costs; (ii) if tomato prices crash (seasonal glut), repayment capacity falls; (iii) risk of borrowing from another MFI simultaneously (multiple indebtedness); (iv) if the credit union fails (governance crisis), her savings are at risk (though COBAC deposit insurance covers up to a limit).

\textbf{Conclusion:} Microfinance fills the \emph{missing middle} in Cameroon's financial system. Effective supervision (COBAC), credit bureaus (to prevent over-indebtedness), and financial education are essential for sustainability.
\end{exoresolu}

\courssub{Appraisal of a Loan File by a Credit Union}

\begin{definition}[Loan Appraisal / Credit Risk Assessment]
\cle{Loan appraisal} is the systematic evaluation of a borrower's ability and willingness to repay a loan. In microfinance, it combines \emph{quantitative} (cash flow, ratios) and \emph{qualitative} (character, social capital) analysis.
\end{definition}

\begin{methode}[The 5 C's of Credit Appraisal (Microfinance Adaptation)]
A credit officer must verify each \cle{C} before recommending approval:
\begin{enumerate}
\item \textbf{Character (Intégrité):} Honesty, reputation in community, credit history (Credit Bureau / \emph{centrale des risques}), willingness to repay. \emph{Most important C.}
\item \textbf{Capacity (Capacité de remboursement):} Ability to generate cash flow to service debt. Compute \cle{Debt Service Ratio (DSR)} = $\dfrac{\text{Monthly instalment}}{\text{Monthly net cash flow}}$. \textbf{Threshold: DSR $\le$ 30--35\%} (COBAC guideline). Stress-test with income shocks.
\item \textbf{Capital (Apport personnel):} Borrower's own funds in the project (skin in the game). Reduces moral hazard. Typically $\ge$ 10--20\% of project cost.
\item \textbf{Collateral (Garanties):} \emph{Primary:} pledged savings, guarantors (caution solidaire), pledge of financed asset. \emph{Secondary:} moral guarantee, group solidarity. In microfinance, social collateral often substitutes physical collateral.
\item \textbf{Conditions (Conditions du prêt):} Purpose (productive vs consumption), amount, duration (matching asset life), interest rate, repayment frequency, economic environment (sector outlook, seasonality).
\end{enumerate}
\end{methode}

\begin{exoresolu}[Appraising Mr. Abena's Loan File at a Credit Union]
Mr. Abena, a carpenter in Buea, applies for 1\,000\,000 FCFA to buy a woodworking machine. File: monthly net income 180\,000 FCFA; personal savings at union 200\,000 FCFA; two salaried guarantors; two previous loans repaid on time. Proposed loan: 24 months, flat total cost giving monthly instalments of 50\,000 FCFA. Appraise using the 5 C's and decide.
\tcblower
\textbf{1. Character:} Two prior loans repaid fully and on time $\rightarrow$ excellent credit history. No record at the \emph{centrale des risques}. Community reputation verified by loan officer visit. \textbf{Positive.}

\textbf{2. Capacity:}
\[
\text{DSR} = \frac{50\,000}{180\,000} = 27.8\% < 33\%.
\]
Cash flow analysis: daily sales $\approx$ 15\,000 FCFA, material costs $\approx$ 6\,000 FCFA, net daily $\approx$ 9\,000 FCFA $\rightarrow$ monthly net $\approx$ 270\,000 FCFA (higher than declared). The machine will increase output by $\sim$30\%. \textbf{Positive.}

\textbf{3. Capital:} Own savings 200\,000 FCFA = 20\% of loan amount. Demonstrates commitment. \textbf{Positive.}

\textbf{4. Collateral:} (a) Savings 200\,000 FCFA pledged (blocked); (b) Two salaried guarantors (net income > 150\,000 FCFA each) signing \emph{caution solidaire}; (c) Machine to be purchased pledged to union (chattel mortgage). \textbf{Adequate.}

\textbf{5. Conditions:} Purpose: productive asset (machine). Duration: 24 months matches useful life. Rate: flat 2\%/month $\approx$ 24\% p.a. effective $\approx$ 27\% — within union's pricing policy. Sector: furniture demand growing in Buea (urbanisation). \textbf{Positive.}

\textbf{Decision:} \textbf{GRANT} the loan with conditions:
\begin{itemize}
\item Disbursement directly to machine supplier (pro-forma invoice required).
\item Machine insured (fire/theft) with union as beneficiary.
\item Quarterly monitoring visits to workshop.
\item No further borrowing from other MFIs without union consent (cross-check credit bureau).
\end{itemize}
\textbf{Key learning:} The 5 C's protect both the borrower (avoids over-indebtedness) and the credit union (preserves members' savings).
\end{exoresolu}

\begin{attention}[Examination Traps in This Chapter]
\begin{itemize}
\item \textbf{Balance sheet:} A loan granted is an \textbf{asset} (claim on borrower); a deposit received is a \textbf{liability} (obligation to depositor). Never invert them.
\item \textbf{Credit creation:} Credit created = Total deposits $-$ Initial deposit. The multiplier $k = 1/r$ applies to the \emph{system}, not a single bank.
\item \textbf{BEAC:} The FCFA/euro peg is fixed by \textbf{treaty with France}, not unilaterally by BEAC. BEAC manages the peg but cannot devalue alone.
\item \textbf{Microfinance categories:} Category 1 = member-based credit unions (savings from members only); Category 2 = MFIs taking public savings; Category 3 = credit-only MFIs. Know the difference.
\item \textbf{Loan appraisal:} \textbf{Character comes before Collateral.} A borrower with perfect collateral but bad character will default. The 33\% DSR is a guideline — always stress-test.
\item \textbf{Terminology:} \emph{Réserves obligatoires} = required reserves (liability side of BEAC, asset side of commercial bank). \emph{Réserves (comptables)} = retained earnings (equity). Do not confuse.
\end{itemize}
\end{attention}

\courssub{Chapter Synthesis: Money and Banking in Cameroon — Key Linkages}

\begin{aretenir}[Essential Linkages for the GCE Exam]
\begin{itemize}
\item \textbf{Commercial banks} create credit within the framework set by \textbf{BEAC} (reserve ratio, key rate, liquidity ratio). Their balance sheets reflect this constraint.
\item \textbf{BEAC's monetary policy} (TIAO, reserves, OMO) influences the \textbf{cost and volume of credit} $\rightarrow$ affects investment, consumption, inflation, and the external balance of the CEMAC zone.
\item \textbf{Microfinance} complements commercial banks by serving the \emph{excluded majority} (informal sector, women, rural). It is supervised by \textbf{COBAC} (same supervisor as banks) but operates with different technologies (social collateral, group lending).
\item \textbf{Loan appraisal} (5 C's) is the micro-foundation of financial stability: poor appraisal $\rightarrow$ high NPLs $\rightarrow$ bank/MFI failure $\rightarrow$ systemic risk $\rightarrow$ BEAC/COBAC intervention.
\item \textbf{Structural challenges in Cameroon:} High cash-drain ratio, shallow interbank market, dominant informal sector, fiscal dominance, weak collateral registry, low financial literacy — all reduce the effectiveness of monetary policy and the reach of formal finance.
\end{itemize}
\end{aretenir}

\begin{exercice}[GCE-Style Practice Questions]
\begin{enumerate}
\item \textbf{Structured question:} The balance sheet of Ngounso Bank shows: Cash 2\,000; BEAC reserves 5\,000; Treasury bills 3\,000; Loans 40\,000; Fixed assets 10\,000; Sight deposits 30\,000; Savings deposits 20\,000; Time deposits 5\,000; Capital 5\,000; Reserves 2\,000. (All in millions FCFA).
\begin{enumerate}
\item Present the balance sheet in regulatory format.
\item Calculate the liquidity ratio (liquid assets / short-term deposits) and comment on compliance with COBAC norms.
\item If BEAC raises the reserve requirement from 7\% to 10\%, explain the impact on this bank's lending capacity.
\end{enumerate}
\item \textbf{Essay question:} "Monetary policy in the CEMAC zone is ineffective in controlling inflation because inflation is mainly imported and supply-driven." Discuss this statement with reference to BEAC's instruments and the transmission mechanism.
\item \textbf{Case study:} A credit union in Bamenda receives a loan application from a group of 10 women farmers for 5\,000\,000 FCFA to buy a maize sheller. Outline the appraisal process using the 5 C's, highlighting the specificities of group lending.
\item \textbf{Calculation:} Initial deposit 1\,000 million FCFA, reserve ratio 8\%, cash-drain ratio 12\%, excess reserve ratio 3\%. Compute the effective money multiplier, total deposits, and total credit created. Compare with the no-leakage case.
\end{enumerate}
\end{exercice}

\begin{corrige}[Exercise 1]
\textbf{(a) Balance Sheet of Ngounso Bank (millions FCFA)}
\begin{center}
\begin{tabularx}{\linewidth}{|X|r|X|r|}
\hline
\rowcolor{popblueL} \textbf{Assets} & & \textbf{Liabilities + Equity} & \\
\hline
Cash in vault & 2\,000 & Share capital & 5\,000 \\
Reserves at BEAC & 5\,000 & Reserves (retained profits) & 2\,000 \\
Treasury bills & 3\,000 & Time deposits & 5\,000 \\
Loans to customers & 40\,000 & Savings deposits & 20\,000 \\
Fixed assets & 10\,000 & Sight deposits & 30\,000 \\
\hline
\textbf{Total} & \textbf{60\,000} & \textbf{Total} & \textbf{60\,000} \\
\hline
\end{tabularx}
\end{center}

\textbf{(b) Liquidity ratio:} Liquid assets = 2\,000 + 5\,000 + 3\,000 = 10\,000. Short-term deposits = 30\,000 + 20\,000 + 5\,000 = 55\,000. Ratio = 10\,000 / 55\,000 = 18.2\%. \textbf{Non-compliant} with COBAC 100\% norm. The bank is highly illiquid and depends on BEAC refinancing.

\textbf{(c) Impact of reserve ratio increase:} Required reserves rise from 7\% $\times$ 55\,000 = 3\,850 to 10\% $\times$ 55\,000 = 5\,500. The bank currently holds 5\,000 at BEAC. It must either borrow 500 on the interbank market, sell Treasury bills, or reduce new lending. Its lending capacity shrinks.
\end{corrige}

\begin{corrige}[Exercise 4]
No leakages: $k = 1/0.08 = 12.5$; $D = 12\,500$; $C = 11\,500$.
With leakages: $k_{\text{eff}} = 1/(0.08 + 0.12 + 0.03) = 1/0.23 \approx 4.348$; $D' \approx 4\,348$; $C' \approx 3\,348$.
Leakages reduce the multiplier by a factor of $\approx 2.87$, illustrating the huge impact of the informal cash economy and banks' caution in Cameroon.
\end{corrige}

\newpage

\coursec{Exercises}

\courssub{I check my knowledge}

\begin{exercice}[MCQ on Money and Banking Fundamentals]
Answer the following multiple choice questions by selecting the correct option.

\begin{enumerate}
\item Which of the following is a \textbf{liability} on a commercial bank's balance sheet?
\begin{itemize}
\item[A.] Cash in vault
\item[B.] Loans granted to customers
\item[C.] Customer deposits
\item[D.] Treasury bills held
\end{itemize}

\item The credit multiplier ($K$) is correctly given by which formula, where $r$ is the reserve ratio?
\begin{itemize}
\item[A.] $K = r$
\item[B.] $K = \frac{1}{r}$
\item[C.] $K = 1 - r$
\item[D.] $K = \frac{r}{1-r}$
\end{itemize}

\item The Bank of Central African States (BEAC) acts as the central bank for which monetary zone?
\begin{itemize}
\item[A.] WAEMU (UEMOA)
\item[B.] CEMAC
\item[C.] ECOWAS (CEDEAO)
\item[D.] SADC
\end{itemize}

\item Which institution is responsible for the supervision and classification of Microfinance Institutions (MFIs) in the CEMAC zone?
\begin{itemize}
\item[A.] BEAC
\item[B.] COBAC
\item[C.] IMF
\item[D.] Ministry of Finance
\end{itemize}

\item In the \textbf{5 Cs of credit appraisal}, which ``C'' refers to the borrower's reputation and track record of repaying debts?
\begin{itemize}
\item[A.] Capacity
\item[B.] Capital
\item[C.] Character
\item[D.] Collateral
\end{itemize}
\end{enumerate}
\end{exercice}

\begin{exercice}[True/False with Justification]
State whether each statement is True or False. \textbf{Justify your answer} in one or two sentences.

\begin{enumerate}
\item Commercial banks create money by lending out their excess reserves.
\item A higher reserve ratio ($r$) leads to a higher credit multiplier and more money creation.
\item The TIAO (Taux d'Intérêt des Appels d'Offres) is the rate at which BEAC lends to commercial banks through open market operations.
\item Microfinance Institutions in Cameroon are allowed to collect savings from the public only if they are classified in Category 1 by COBAC.
\item ``Collateral'' in the 5 Cs framework refers exclusively to land titles owned by the borrower.
\end{enumerate}
\end{exercice}

\begin{exercice}[Key Definitions]
Define the following terms precisely, as required in the GCE Advanced Level Economics examination.

\begin{enumerate}
\item Balance sheet of a commercial bank
\item Credit creation (money multiplier process)
\item Legal reserve requirement (reserve ratio)
\item Central Bank (with reference to BEAC)
\item Microfinance Institution (MFI) -- Category 2 (COBAC classification)
\item The 5 Cs of credit appraisal
\end{enumerate}
\end{exercice}

\begin{exercice}[Direct Calculations: Credit Multiplier and Money Creation]
Assume the banking system in Cameroon has an initial deposit of \SI{2\,000\,000}{FCFA} and a legal reserve ratio of $10\%$.

\begin{enumerate}
\item Calculate the credit multiplier ($K$).
\item Calculate the maximum total deposits that can be created by the banking system.
\item Calculate the maximum total new credit (loans) that can be generated.
\item Recalculate (i), (ii) and (iii) if the reserve ratio is increased to $20\%$.
\item Briefly comment on the effect of changing the reserve ratio on the money supply.
\end{enumerate}
\end{exercice}

\courssub{I practise}

\begin{exercice}[Constructing a Commercial Bank Balance Sheet]
The following items (in millions of FCFA) are extracted from the books of \textbf{Société Générale des Banques au Cameroun (SGBC)} at the end of the financial year:

\begin{center}
\begin{tabular}{lr}
Cash in vault & 1\,200 \\
Reserves at BEAC & 3\,500 \\
Treasury Bills (short-term) & 2\,800 \\
Loans granted to customers & 15\,000 \\
Overdrafts & 4\,200 \\
Buildings and equipment & 5\,000 \\
Customer deposits: Current accounts & 12\,000 \\
Customer deposits: Savings accounts & 8\,000 \\
Customer deposits: Term deposits & 5\,000 \\
Borrowed funds (from BEAC/other banks) & 3\,000 \\
Equity (Share capital + Reserves) & 3\,700 \\
\end{tabular}
\end{center}

\begin{enumerate}
\item Construct the complete Balance Sheet of SGBC in the standard format (Assets on the left, Liabilities \& Capital on the right). Ensure it balances.
\item Calculate the \textbf{Liquidity Ratio} defined as:
\[
\text{Liquidity Ratio} = \frac{\text{Liquid Assets (Cash + Reserves at BEAC + Treasury Bills)}}{\text{Total Customer Deposits}} \times 100\%
\]
\item Comment on the liquidity position of SGBC if the regulatory minimum liquidity ratio set by COBAC is $30\%$.
\end{enumerate}
\end{exercice}

\begin{exercice}[Credit Creation Process: Numerical Application]
Consider a simplified banking system with only one commercial bank (or a system where all banks operate identically). The legal reserve ratio ($r$) is $10\%$. A customer deposits \SI{5\,000\,000}{FCFA} in cash at the bank.

\begin{enumerate}
\item Show the changes in the bank's balance sheet \textbf{after the initial deposit} (before any lending).
\item The bank lends out its maximum excess reserves. Show the balance sheet \textbf{after the first loan} is granted and the borrower writes a cheque that is deposited in the same bank (or another bank in the system).
\item Continue the process for \textbf{two more rounds} of lending and re-depositing. Present a summary table showing: Round, New Deposit, Required Reserves, Excess Reserves (New Loans), Cumulative Deposits, Cumulative Loans.
\item Using the credit multiplier formula, state the theoretical maximum total deposits, total loans, and total money supply created from the initial \SI{5\,000\,000}{FCFA}.
\item Explain why the actual money creation in Cameroon might be lower than this theoretical maximum. Give \textbf{three} reasons.
\end{enumerate}
\end{exercice}

\begin{exercice}[Microfinance Institutions vs Commercial Banks]
Answer the following short-answer questions.

\begin{enumerate}
\item Distinguish between a \textbf{Commercial Bank} and a \textbf{Microfinance Institution (MFI)} in Cameroon using \textbf{four} criteria: (a) Legal framework/Supervisory body, (b) Target clientele, (c) Range of services, (d) Capital requirements.
\item State \textbf{three} specific financial services typically offered by MFIs in Cameroon (e.g., UCC, MC2) that are tailored to the informal sector or low-income households.
\item Explain \textbf{why} the BEAC raising the TIAO (Taux d'Intérêt des Appels d'Offres) tends to reduce inflation in the CEMAC zone. Use the transmission mechanism: TIAO $\rightarrow$ Bank lending rates $\rightarrow$ Investment/Consumption $\rightarrow$ Aggregate Demand $\rightarrow$ Price Level.
\end{enumerate}
\end{exercice}

\begin{exercice}[Loan Appraisal using the 5 Cs -- Practical Case]
A farmer, Mr. Nfor, applies for a loan of \SI{600\,000}{FCFA} from a Credit Union (Category 2 MFI) to purchase a cocoa dryer and improve his fermentation process. He submits the following information:

\begin{itemize}
\item Monthly net income from cocoa sales: \SI{120\,000}{FCFA} (seasonal, high Oct--Dec, low Jan--Sep).
\item Monthly household expenses: \SI{70\,000}{FCFA}.
\item He owns the cocoa farm (5 hectares) but has \textbf{no land title} (customary land tenure).
\item He provides two guarantors: a civil servant (salary \SI{250\,000}{FCFA}/month) and a trader (estimated income \SI{180\,000}{FCFA}/month).
\item He has been a member of the Credit Union for 3 years, with savings of \SI{150\,000}{FCFA} and a previous loan of \SI{200\,000}{FCFA} repaid on schedule.
\item The dryer costs \SI{550\,000}{FCFA}; installation and training \SI{50\,000}{FCFA}.
\end{itemize}

Appraise this loan file using the \textbf{5 Cs framework} (Character, Capacity, Capital, Collateral, Conditions). For each ``C'', state your assessment (Strong / Adequate / Weak) and justify with data from the case. Conclude with a recommendation: \textbf{Approve, Approve with conditions, or Reject}.
\end{exercice}

\courssub{I prepare for the GCE}

\begin{exercice}[Essay: Functions of BEAC as Central Bank -- 25 marks]
\textbf{``Explain the functions of a central bank, illustrating each with the activities of the Bank of Central African States (BEAC) in the CEMAC zone.''}

\textit{Guidance for structuring your answer (not part of the exam question):}
\begin{itemize}
\item Introduction: Define central bank; state BEAC's role in CEMAC (6 member states, common currency CFA Franc).
\item Function 1: Monetary Policy / Price Stability (TIAO, reserve requirements, open market operations).
\item Function 2: Banker to the Government (managing public debt, Treasury accounts -- note CEMAC rules on advances).
\item Function 3: Banker to Commercial Banks (lender of last resort, clearing house, holding reserves).
\item Function 4: Issuer of Currency (monopoly on banknotes/coins, managing foreign exchange reserves).
\item Function 5: Supervision \& Regulation (with COBAC -- prudential ratios, licensing).
\item Function 6: Manager of Foreign Exchange Reserves (operations account with French Treasury, exchange rate peg).
\item Conclusion: Summary of BEAC's effectiveness / challenges (e.g., asymmetric shocks, limited policy autonomy).
\end{itemize}
\end{exercice}

\begin{exercice}[Essay: Problems of Commercial Banks in Cameroon -- 25 marks]
\textbf{``Discuss the problems faced by commercial banks in Cameroon and suggest possible solutions.''}

\textit{Guidance for structuring your answer (not part of the exam question):}
\begin{itemize}
\item Introduction: Importance of commercial banks in financial intermediation; context of Cameroonian banking sector (oligopoly, foreign-owned dominance).
\item Problem 1: High Non-Performing Loans (NPLs) -- causes (poor appraisal, weak legal enforcement, economic shocks).
\item Problem 2: Excess Liquidity / Lack of viable projects -- paradox of high liquidity coexisting with credit rationing for SMEs.
\item Problem 3: Weak Judicial System / Collateral enforcement difficulties (land title issues, lengthy court procedures).
\item Problem 4: High Operating Costs / Infrastructure deficits (electricity, transport, ICT) reducing profitability.
\item Problem 5: Regulatory Burden / Compliance costs (COBAC Basel II/III implementation, AML/CFT).
\item Problem 6: Competition from Mobile Money / Fintech / MFIs eroding traditional deposit base.
\item Solutions: For each problem, propose a concrete policy or bank-level solution (e.g., credit bureaus, movable collateral registry, specialized SME desks, digital transformation, public-private partnerships).
\item Conclusion: Balanced assessment -- progress made (e.g., CEMAC banking reforms) vs persistent structural issues.
\end{itemize}
\end{exercice}

\begin{exercice}[Case Study: Credit Union Loan Appraisal -- 20 marks]
\textbf{Case Study:} Mrs. Aisha runs a small food processing cooperative (cassava flour, ``garri'') in a rural area of the East Region. She applies for a loan of \SI{1\,500\,000}{FCFA} from a Category 2 MFI (Credit Union) to buy a mechanical grater, a hydraulic press, and build a drying shed. The total project cost is \SI{1\,800\,000}{FCFA}; she contributes \SI{300\,000}{FCFA} from own savings.

\textbf{File Summary:}
\begin{itemize}
\item \textbf{Character:} Member for 5 years; savings \SI{400\,000}{FCFA}; two previous loans (\SI{500\,000}{FCFA} and \SI{800\,000}{FCFA}) repaid early. Treasurer of her church women's group. No credit bureau record (rural area).
\item \textbf{Capacity:} Cooperative monthly revenue \SI{1\,200\,000}{FCFA}; operating costs \SI{700\,000}{FCFA} (raw material, labour, transport); net surplus \SI{500\,000}{FCFA}/month. Loan repayment proposed: \SI{100\,000}{FCFA}/month over 18 months. Household expenses \SI{150\,000}{FCFA}/month.
\item \textbf{Capital:} Own contribution \SI{300\,000}{FCFA} (16.7\% of project cost). Cooperative assets (equipment, stock) valued at \SI{2\,000\,000}{FCFA}.
\item \textbf{Collateral:} No land title for the shed location (community land). Proposes: (a) Chattel mortgage on the new equipment (grater, press) valued at \SI{1\,200\,000}{FCFA}; (b) Personal guarantee from the Cooperative President (civil servant, salary \SI{300\,000}{FCFA}); (c) Pledge of her savings \SI{400\,000}{FCFA} in the Credit Union.
\item \textbf{Conditions:} Growing local demand for packaged garri; new rural market opening 5 km away; input prices (cassava) stable; no major competitors. Risk: Seasonal cassava shortage (Jan--Mar), poor road access during rains.
\end{itemize}

\textbf{Tasks:}
\begin{enumerate}
\item Appraise the loan file using the \textbf{5 Cs}. For each C, provide a rating (Strong / Adequate / Weak) and a detailed justification referencing the data above. [10 marks]
\item Calculate the \textbf{Debt Service Coverage Ratio (DSCR)} for the cooperative: $\frac{\text{Net Surplus}}{\text{Proposed Monthly Repayment}}$. Interpret the result. [3 marks]
\item As the Credit Committee, state your final decision: \textbf{Approve / Approve with Conditions / Reject}. If ``Approve with Conditions'', specify \textbf{three} binding conditions (covenants) to mitigate the identified risks. [7 marks]
\end{enumerate}
\end{exercice}

\newpage

\coursec{Solutions to the exercises}

\begin{corrige}[Exercise 1 — MCQ on Money and Banking Fundamentals]

\begin{enumerate}
\item \textbf{Correct answer: C — Customer deposits.}\par
\textbf{Justification:} On a commercial bank's balance sheet, the \emph{liabilities} side records everything the bank \emph{owes}. Customer deposits are debts of the bank towards its clients, who can withdraw them on demand (current accounts) or after notice (savings and term accounts). Options A, B and D are \emph{assets} (uses of funds): cash in vault is a liquid asset, loans granted are claims on customers, and Treasury bills are financial investments.

\item \textbf{Correct answer: B — $K = \dfrac{1}{r}$.}\par
\textbf{Justification:} In the simplified credit-creation model (no cash leakage, banks lend all excess reserves), the credit multiplier is the inverse of the legal reserve ratio. With $r = 10\%$, $K = 10$: each franc of new deposit can support ten francs of total deposits in the system. Option A confuses the multiplier with the ratio itself; C and D have no economic meaning here.

\item \textbf{Correct answer: B — CEMAC.}\par
\textbf{Justification:} The BEAC (Bank of Central African States) is the common central bank of the six member states of the \emph{Communauté Économique et Monétaire de l'Afrique Centrale}: Cameroon, Chad, Central African Republic, Congo, Equatorial Guinea and Gabon, which share the CFA Franc (XAF). WAEMU is served by the BCEAO; ECOWAS and SADC are not monetary unions with a single central bank of this type.

\item \textbf{Correct answer: B — COBAC.}\par
\textbf{Justification:} The \emph{Commission Bancaire de l'Afrique Centrale} (COBAC) is the CEMAC banking supervisory authority. It licenses, classifies (Categories 1, 2 and 3) and supervises microfinance institutions, in addition to supervising commercial banks. The BEAC conducts monetary policy; the Ministry of Finance handles State finances; the IMF is an external institution.

\item \textbf{Correct answer: C — Character.}\par
\textbf{Justification:} \emph{Character} assesses the borrower's honesty, reputation and past repayment behaviour (credit history, membership record in the Credit Union). \emph{Capacity} is the ability to repay from income; \emph{Capital} is the borrower's own financial contribution; \emph{Collateral} is the security pledged against the loan; \emph{Conditions} refer to the purpose of the loan and the economic environment.
\end{enumerate}

\begin{attention}[Typical MCQ traps]
Candidates frequently confuse \emph{assets} and \emph{liabilities} (deposits ``belong'' to customers, hence they are a liability \emph{for the bank}), and confuse BEAC (monetary policy) with COBAC (banking supervision). Always ask: ``who owes what to whom?'' and ``who does what?''.
\end{attention}
\end{corrige}

\begin{corrige}[Exercise 2 — True/False with Justification]

\begin{enumerate}
\item \textbf{TRUE.} When a bank's actual reserves exceed its legal (required) reserves, it can grant new loans by simply crediting the borrower's account; this new deposit is new (scriptural) money. Money is destroyed symmetrically when loans are repaid.

\item \textbf{FALSE.} The credit multiplier is $K = \dfrac{1}{r}$, so it varies \emph{inversely} with the reserve ratio. A \emph{higher} $r$ forces banks to immobilise more reserves per franc of deposit, leaving less to lend: the multiplier \emph{falls} (e.g. $r = 10\% \Rightarrow K = 10$; $r = 20\% \Rightarrow K = 5$).

\item \textbf{TRUE.} The TIAO (\emph{Taux d'Intérêt des Appels d'Offres}) is BEAC's key policy rate: the rate at which commercial banks obtain central bank refinancing through weekly tenders (open market operations). Changes in the TIAO are transmitted to banks' lending rates and hence to economic activity and prices.

\item \textbf{TRUE.} Under the COBAC classification, only \textbf{Category 1} MFIs (savings and credit cooperatives such as Credit Unions and MC\textsuperscript{2}) are authorised to collect savings/deposits from their members \emph{and} from the public within regulatory limits. Category 2 institutions grant credit only (no deposit collection), and Category 3 covers other limited operators.

\item \textbf{FALSE.} Collateral covers \emph{any} asset pledged to secure a loan: land titles and buildings (mortgages), but also vehicles, equipment (chattel mortgage), stocks of goods, savings pledged in the institution, and even personal guarantees (guarantors). Restricting collateral to land titles is precisely one of the obstacles to credit access in Cameroon, where many borrowers hold land under customary tenure without formal titles.
\end{enumerate}
\end{corrige}

\begin{corrige}[Exercise 3 — Key Definitions]

\begin{enumerate}
\item \textbf{Balance sheet of a commercial bank:} A financial statement, at a given date, showing the bank's \emph{assets} (uses of funds: cash, reserves at the central bank, securities, loans and overdrafts, fixed assets) on the left and its \emph{liabilities and equity} (sources of funds: customer deposits, borrowed funds, share capital and reserves) on the right, with the two sides always equal: Total Assets $=$ Total Liabilities $+$ Equity.

\item \textbf{Credit creation (money multiplier process):} The process by which the banking system as a whole creates scriptural money by granting loans: an initial deposit generates excess reserves, which are lent and re-deposited repeatedly, so that total deposits equal the initial deposit multiplied by the credit multiplier $K = \dfrac{1}{r}$, where $r$ is the reserve ratio.

\item \textbf{Legal reserve requirement (reserve ratio):} The minimum fraction $r$ of customer deposits that every commercial bank is obliged by the monetary authority (BEAC/COBAC) to hold as reserves (in vault cash or on its account at the central bank) rather than lend out. It is an instrument of monetary policy and of banking security.

\item \textbf{Central Bank (with reference to BEAC):} The public monetary institution that holds the monopoly of currency issue, conducts monetary policy, acts as banker to the State and to commercial banks (lender of last resort), manages foreign exchange reserves and (with COBAC) supervises the banking system. The \textbf{BEAC} fulfils these functions for the six CEMAC states, issuing the CFA Franc (XAF) from its headquarters in Yaoundé.

\item \textbf{Microfinance Institution — Category 2 (COBAC):} A financial institution, licensed and supervised by COBAC, authorised to \emph{grant microcredit} to low-income clients and the informal sector but \emph{not authorised to collect savings/deposits from the public}; it finances its lending from own funds and borrowing. (Category 1 = savings and credit cooperatives collecting deposits; Category 3 = other limited operators.)

\item \textbf{The 5 Cs of credit appraisal:} The five criteria used by a lender to evaluate a loan file: \textbf{Character} (honesty, repayment track record), \textbf{Capacity} (ability to repay from income/cash flow), \textbf{Capital} (borrower's own contribution and net worth), \textbf{Collateral} (security pledged) and \textbf{Conditions} (purpose of the loan and the economic environment of the activity).
\end{enumerate}

\begin{methode}[Writing a good GCE definition]
A full-mark definition contains: (i) the \emph{genus} (what kind of thing it is — a statement, a process, a ratio, an institution); (ii) the \emph{essential characteristics} (what distinguishes it); (iii) where useful, a formula or an example. One or two precise sentences suffice — do not write a paragraph of vague description.
\end{methode}
\end{corrige}

\begin{corrige}[Exercise 4 — Direct Calculations: Credit Multiplier and Money Creation]

\textbf{Given:} initial deposit $D_0 = \SI{2\,000\,000}{FCFA}$; reserve ratio $r = 10\% = 0.10$.

\begin{enumerate}
\item \textbf{Credit multiplier:}
\[ K = \frac{1}{r} = \frac{1}{0.10} = 10 \]

\item \textbf{Maximum total deposits:}
\[ D_{\max} = K \times D_0 = 10 \times 2\,000\,000 = \SI{20\,000\,000}{FCFA} \]

\item \textbf{Maximum total new credit (loans):} Total loans equal total deposits minus the initial deposit (which is ``old'' money, not created):
\[ C_{\max} = D_{\max} - D_0 = 20\,000\,000 - 2\,000\,000 = \SI{18\,000\,000}{FCFA} \]
(Equivalently, $C_{\max} = D_0 \times \frac{1-r}{r} = 2\,000\,000 \times 9$.)

\item \textbf{With $r = 20\% = 0.20$:}
\begin{align*}
K &= \frac{1}{0.20} = 5 \\
D_{\max} &= 5 \times 2\,000\,000 = \SI{10\,000\,000}{FCFA} \\
C_{\max} &= 10\,000\,000 - 2\,000\,000 = \SI{8\,000\,000}{FCFA}
\end{align*}

\item \textbf{Comment:} Doubling the reserve ratio from $10\%$ to $20\%$ halves the multiplier (from 10 to 5) and therefore halves the money-creation capacity of the banking system (deposits fall from 20 to 10 million FCFA; new credit from 18 to 8 million FCFA). The reserve ratio is thus a powerful \emph{contractionary/expansionary} instrument: raising $r$ reduces the money supply (fighting inflation), while lowering $r$ expands credit (stimulating activity).
\end{enumerate}

\begin{attention}[Common mistake]
Do not confuse \emph{total deposits} ($K \times D_0$) with \emph{total new credit} ($K \times D_0 - D_0$). The initial deposit is not ``created'' money; examiners deduct marks when candidates present 20 million as ``credit created''.
\end{attention}
\end{corrige}

\begin{corrige}[Exercise 5 — Constructing a Commercial Bank Balance Sheet]

\textbf{1. Balance sheet of SGBC (millions of FCFA).} Classification: cash, reserves at BEAC, Treasury bills, loans, overdrafts and buildings are \emph{assets}; deposits and borrowed funds are \emph{liabilities}; equity is \emph{capital}.

\begin{center}
\begin{tabularx}{\linewidth}{|X r|X r|}
\hline
\rowcolor{popblueL} \textbf{ASSETS} & \textbf{Amount} & \textbf{LIABILITIES \& CAPITAL} & \textbf{Amount} \\
\hline
Cash in vault & 1\,200 & Current account deposits & 12\,000 \\
\hline
Reserves at BEAC & 3\,500 & Savings account deposits & 8\,000 \\
\hline
Treasury Bills & 2\,800 & Term deposits & 5\,000 \\
\hline
Loans to customers & 15\,000 & Borrowed funds (BEAC/banks) & 3\,000 \\
\hline
Overdrafts & 4\,200 & Equity (capital + reserves) & 3\,700 \\
\hline
Buildings and equipment & 5\,000 & & \\
\hline
\rowcolor{popgreenL} \textbf{TOTAL ASSETS} & \textbf{31\,700} & \textbf{TOTAL LIAB. + CAPITAL} & \textbf{31\,700} \\
\hline
\end{tabularx}
\end{center}

\textbf{Check:} Assets: $1\,200 + 3\,500 + 2\,800 + 15\,000 + 4\,200 + 5\,000 = 31\,700$. Liabilities + capital: $12\,000 + 8\,000 + 5\,000 + 3\,000 + 3\,700 = 31\,700$. The balance sheet \emph{balances}. $\blacksquare$

\textbf{2. Liquidity Ratio:}
\[
\text{Liquid assets} = 1\,200 + 3\,500 + 2\,800 = 7\,500 \text{ (millions FCFA)}
\]
\[
\text{Total customer deposits} = 12\,000 + 8\,000 + 5\,000 = 25\,000 \text{ (millions FCFA)}
\]
\[
\text{Liquidity Ratio} = \frac{7\,500}{25\,000} \times 100\% = 30\%
\]

\textbf{3. Comment:} SGBC's liquidity ratio ($30\%$) is exactly \emph{equal} to the COBAC regulatory minimum of $30\%$. The bank is therefore \emph{just compliant}: it holds enough liquid assets to meet the prudential requirement, but with \textbf{no safety margin}. Any unexpected deposit withdrawal, or any fall in the value of its Treasury bills, would push the bank \emph{below} the minimum and expose it to COBAC sanctions or forced refinancing at BEAC (at the TIAO or penalty rate). Prudent management would target a ratio \emph{above} the minimum (e.g. $33$--$35\%$), even though holding more low-yield liquid assets reduces profitability — the classic \emph{liquidity--profitability trade-off}.
\end{corrige}

\begin{corrige}[Exercise 6 — Credit Creation Process: Numerical Application]

\textbf{Given:} $r = 10\%$; initial cash deposit $D_0 = \SI{5\,000\,000}{FCFA}$.

\textbf{1. Balance sheet after the initial deposit (before lending):}

\begin{center}
\begin{tabularx}{\linewidth}{|X r|X r|}
\hline
\rowcolor{popblueL} \textbf{ASSETS} & \textbf{FCFA} & \textbf{LIABILITIES} & \textbf{FCFA} \\
\hline
Cash / Reserves & 5\,000\,000 & Demand deposit (Customer A) & 5\,000\,000 \\
\hline
\rowcolor{popgreenL} \textbf{Total} & \textbf{5\,000\,000} & \textbf{Total} & \textbf{5\,000\,000} \\
\hline
\end{tabularx}
\end{center}

Required reserves: $0.10 \times 5\,000\,000 = 500\,000$; \textbf{excess reserves} $= 5\,000\,000 - 500\,000 = 4\,500\,000$.

\textbf{2. After the first loan (4\,500\,000 lent, then re-deposited by the cheque recipient):}

\begin{center}
\begin{tabularx}{\linewidth}{|X r|X r|}
\hline
\rowcolor{popblueL} \textbf{ASSETS} & \textbf{FCFA} & \textbf{LIABILITIES} & \textbf{FCFA} \\
\hline
Reserves & 5\,000\,000 & Deposit A & 5\,000\,000 \\
\hline
Loan to B & 4\,500\,000 & Deposit (recipient of B's cheque) & 4\,500\,000 \\
\hline
\rowcolor{popgreenL} \textbf{Total} & \textbf{9\,500\,000} & \textbf{Total} & \textbf{9\,500\,000} \\
\hline
\end{tabularx}
\end{center}

New required reserves: $0.10 \times 9\,500\,000 = 950\,000$; excess reserves $= 5\,000\,000 - 950\,000 = 4\,050\,000$, which can be lent in round 3.

\textbf{3. Summary table (rounds 1 to 4):}

\begin{center}
\begin{tabularx}{\linewidth}{|c|X|X|X|X|X|}
\hline
\rowcolor{popblueL} \textbf{Round} & \textbf{New deposit} & \textbf{Required reserves (10\%)} & \textbf{New loan (excess)} & \textbf{Cumul. deposits} & \textbf{Cumul. loans} \\
\hline
1 & 5\,000\,000 & 500\,000 & 4\,500\,000 & 5\,000\,000 & 0 \\
\hline
2 & 4\,500\,000 & 450\,000 & 4\,050\,000 & 9\,500\,000 & 4\,500\,000 \\
\hline
3 & 4\,050\,000 & 405\,000 & 3\,645\,000 & 13\,550\,000 & 8\,550\,000 \\
\hline
4 & 3\,645\,000 & 364\,500 & 3\,280\,500 & 17\,195\,000 & 12\,195\,000 \\
\hline
\end{tabularx}
\end{center}

Each new deposit is $90\%$ of the previous one: the process is a geometric progression of ratio $(1-r) = 0.9$.

\textbf{4. Theoretical maximum (multiplier formula):}
\begin{align*}
K &= \frac{1}{r} = \frac{1}{0.10} = 10 \\
D_{\max} &= 10 \times 5\,000\,000 = \SI{50\,000\,000}{FCFA} \text{ (total deposits)} \\
C_{\max} &= 50\,000\,000 - 5\,000\,000 = \SI{45\,000\,000}{FCFA} \text{ (total loans)} \\
\text{Money created} &= \SI{45\,000\,000}{FCFA} \text{ of new scriptural money.}
\end{align*}

\textbf{5. Why actual money creation in Cameroon is lower — three reasons:}
\begin{itemize}
\item \textbf{Cash leakages:} in a largely cash-based, informal economy, borrowers withdraw a significant part of loans in banknotes instead of re-depositing them; each leakage withdraws reserves from the banking circuit and shrinks the multiplier.
\item \textbf{Banks do not lend all excess reserves:} faced with high non-performing loans and few ``bankable'' projects, Cameroonian banks voluntarily hold \emph{excess liquidity} (or buy Treasury bills) rather than lend, especially to SMEs.
\item \textbf{Limited demand for credit / bankable demand:} many potential borrowers lack land titles, guarantors or formal accounts, so they cannot obtain loans; without new loans, the multiplication process stops. (One may also cite: prudential liquidity ratios imposed by COBAC above the legal reserve, and imports/payments abroad that drain reserves out of the domestic system.)
\end{itemize}
\end{corrige}

\begin{corrige}[Exercise 7 — Microfinance Institutions vs Commercial Banks]

\textbf{1. Four distinguishing criteria:}

\begin{center}
\begin{tabularx}{\linewidth}{|l|X|X|}
\hline
\rowcolor{popblueL} \textbf{Criterion} & \textbf{Commercial Bank} & \textbf{Microfinance Institution} \\
\hline
(a) Legal framework / supervision & Banking law of CEMAC; licensed and supervised by COBAC as \emph{credit establishments}; large minimum capital (billions of FCFA). & Specific CEMAC microfinance regulation (COBAC classification, Categories 1--3); often cooperative law (OHADA / COBAC); much lower minimum capital. \\
\hline
(b) Target clientele & Formal sector: companies, civil servants, salaried workers, traders with guarantees and land titles. & Informal sector and low-income households: farmers, petty traders, artisans, women's groups, rural populations. \\
\hline
(c) Range of services & Full range: current accounts, cheques, overdrafts, international transfers, trade finance, cards, foreign exchange. & Narrow range: member savings, small loans (microcredit), group lending, sometimes money transfer; no cheques or foreign operations. \\
\hline
(d) Capital requirements & High minimum share capital fixed by COBAC/CEMAC banking regulation. & Low minimum capital adapted to cooperative size; funds come mainly from members' shares and savings. \\
\hline
\end{tabularx}
\end{center}

\textbf{2. Three services tailored to the informal sector (e.g. UCC, MC\textsuperscript{2}):}
\begin{itemize}
\item \textbf{Microcredit / small seasonal loans} without land-title collateral (e.g. loans to farmers for inputs, repaid at harvest), often with guarantors or group guarantee instead.
\item \textbf{Mobilisation of small, flexible savings} (daily or weekly deposits, ``susu''-type collection), accessible in neighbourhoods and villages where banks have no branches.
\item \textbf{Group (solidarity) lending and financial education:} loans to solidarity groups (e.g. \emph{njangui}/tontine members) where members guarantee each other, plus training in basic management of income and repayment discipline.
\end{itemize}

\textbf{3. Transmission mechanism of a TIAO increase:}
\begin{itemize}
\item \textbf{TIAO $\uparrow$:} refinancing at BEAC becomes more expensive for commercial banks (their cost of funds rises).
\item \textbf{$\rightarrow$ Bank lending rates $\uparrow$:} banks pass the higher cost onto customers; credit also becomes scarcer as banks' liquidity tightens.
\item \textbf{$\rightarrow$ Investment and consumption $\downarrow$:} firms borrow less for investment; households reduce credit-financed consumption.
\item \textbf{$\rightarrow$ Aggregate demand $\downarrow$:} total spending in the CEMAC economy contracts relative to supply.
\item \textbf{$\rightarrow$ Price level:} demand-pull inflationary pressure eases, so the general price level stabilises (inflation falls).
\end{itemize}
This is the standard \emph{restrictive monetary policy} channel through which BEAC defends price stability and the internal/external value of the CFA Franc.
\end{corrige}

\begin{corrige}[Exercise 8 — Loan Appraisal using the 5 Cs (Mr. Nfor)]

\textbf{Loan requested:} \SI{600\,000}{FCFA} (dryer 550\,000 + installation/training 50\,000 — the budget is coherent and fully costed).

\begin{enumerate}
\item \textbf{Character — STRONG.} Member for 3 years; savings of \SI{150\,000}{FCFA}; a previous loan of \SI{200\,000}{FCFA} repaid \emph{on schedule}. This demonstrates reliability and attachment to the Credit Union.

\item \textbf{Capacity — ADEQUATE (with reservation).} Net monthly income \SI{120\,000}{FCFA} minus household expenses \SI{70\,000}{FCFA} leaves a surplus of about \SI{50\,000}{FCFA}/month, i.e. \SI{600\,000}{FCFA}/year — sufficient to repay \SI{600\,000}{FCFA} over roughly 12--18 months with interest. \emph{However}, income is highly \emph{seasonal} (high Oct--Dec, low Jan--Sep): repayment must be scheduled on the cocoa season, not in equal monthly instalments. The dryer should also \emph{raise} income (better-fermented, well-dried cocoa sells at a higher price), improving future capacity.

\item \textbf{Capital — ADEQUATE.} Own savings of \SI{150\,000}{FCFA} represent $25\%$ of the loan amount, and he owns a productive 5-hectare farm. He has a real stake in the project, though his cash contribution to the project itself is modest.

\item \textbf{Collateral — WEAK to ADEQUATE.} No land title (customary tenure) means no mortgage is possible — a classic Cameroonian constraint. Mitigating elements: two guarantors, including a civil servant with a stable salary of \SI{250\,000}{FCFA}/month (a strong guarantee) and a trader (\SI{180\,000}{FCFA}/month); plus the possibility of pledging his \SI{150\,000}{FCFA} savings and taking a chattel mortgage on the dryer. The guarantee package partly compensates the absence of a land title.

\item \textbf{Conditions — STRONG.} The purpose is \emph{productive} (equipment that improves quality and value added, not consumption); cocoa is a strategic Cameroonian export crop with steady demand; the total cost (\SI{600\,000}{FCFA}) matches the quotation exactly.
\end{enumerate}

\textbf{Recommendation: APPROVE WITH CONDITIONS.}
\begin{itemize}
\item Repayment schedule \emph{indexed on the cocoa season} (larger instalments Oct--Dec, token payments off-season).
\item Pledge of the \SI{150\,000}{FCFA} savings + signed guarantees of the two guarantors (civil servant first) + chattel mortgage on the dryer.
\item Disbursement directly to the equipment supplier (not in cash to the borrower) to guarantee the use of funds.
\end{itemize}
\end{corrige}

\begin{corrige}[Exercise 9 — Essay: Functions of BEAC as Central Bank (25 marks)]

\textbf{Indicative mark scheme (examiner expectations):}

\begin{center}
\begin{tabularx}{\linewidth}{|X|c|}
\hline
\rowcolor{popblueL} \textbf{Component} & \textbf{Marks} \\
\hline
Introduction: definition of a central bank; BEAC, CEMAC (6 states), CFA Franc & 3 \\
\hline
Function 1: Monetary policy / price stability (TIAO, reserve requirements, open market) & 4 \\
\hline
Function 2: Banker to the Government (Treasury accounts, public debt; statutory limits on advances) & 3 \\
\hline
Function 3: Banker to commercial banks (refinancing, lender of last resort, clearing, reserves) & 4 \\
\hline
Function 4: Issuer of currency (monopoly of issue of notes and coins) & 3 \\
\hline
Function 5: Supervision and regulation (with COBAC: licensing, prudential ratios) & 3 \\
\hline
Function 6: Management of foreign exchange reserves (peg to the euro, operations account with the French Treasury) & 3 \\
\hline
Conclusion: balanced assessment (effectiveness, challenges) & 2 \\
\hline
\rowcolor{popgreenL} \textbf{Total} & \textbf{25} \\
\hline
\end{tabularx}
\end{center}

\textbf{Model answer (outline developed):}

\textbf{Introduction.} A central bank is the public institution at the apex of a country's monetary and banking system, responsible for issuing currency, conducting monetary policy and safeguarding financial stability. For the six CEMAC states (Cameroon, Chad, Central African Republic, Congo, Equatorial Guinea, Gabon), this role is performed collectively by the \textbf{Bank of Central African States (BEAC)}, headquartered in Yaoundé, which issues the common currency, the CFA Franc (XAF).

\textbf{1. Conduct of monetary policy and price stability.} BEAC's primary mission is price stability in the zone. It fixes its key policy rate, the \textbf{TIAO} (rate of its weekly liquidity tenders), sets \textbf{legal reserve requirements} for commercial banks, and conducts open market operations. Raising the TIAO makes bank refinancing dearer, restrains credit and cools inflation; lowering it stimulates activity.

\textbf{2. Banker to the Government.} BEAC keeps the accounts of the member States' Treasuries, manages their public debt issuance (Treasury bills and bonds on the CEMAC market) and may grant advances \emph{within strict statutory limits} — a discipline designed to prevent the inflationary monetary financing of budget deficits that marked the region in the past.

\textbf{3. Banker to commercial banks.} Commercial banks hold their required reserves at BEAC, obtain refinancing there (BEAC is \emph{lender of last resort} to solvent but illiquid banks), and settle interbank payments through BEAC's clearing and payment systems. BEAC is thus the hub of the zone's banking circuit.

\textbf{4. Issuer of currency.} BEAC holds the \emph{monopoly of issuing} banknotes and coins in CEMAC. It guarantees the quality and confidence in the currency, replaces worn notes and combats counterfeiting.

\textbf{5. Supervision and regulation.} Together with \textbf{COBAC} (the banking commission), BEAC participates in licensing banks and MFIs and enforcing prudential rules (capital adequacy, liquidity ratios, classification of microfinance institutions), protecting depositors and systemic stability.

\textbf{6. Manager of foreign exchange reserves.} BEAC centralises and manages the zone's external reserves, which back the \emph{fixed peg} of the CFA Franc to the euro, guaranteed by the French Treasury through the \emph{operations account} mechanism, in exchange for the pooling of reserves and convertibility rules.

\textbf{Conclusion.} BEAC performs all the classical functions of a central bank, but in a \emph{multinational} framework: this guarantees monetary credibility and a stable, convertible currency, at the cost of limited national policy autonomy and difficulties in responding to asymmetric shocks (e.g. an oil-price fall hitting Congo differently from Cameroon). Recent CEMAC reforms aim to strengthen its independence, transparency and supervisory effectiveness.

\textbf{Examiner expectations:} each function must be (i) \emph{named}, (ii) \emph{explained} generally, and (iii) \emph{illustrated} with a BEAC-specific instrument or fact (TIAO, COBAC, operations account, Yaoundé headquarters). A purely theoretical answer with no BEAC illustration loses the application marks.
\end{corrige}

\begin{corrige}[Exercise 10 — Essay: Problems of Commercial Banks in Cameroon (25 marks)]

\textbf{Indicative mark scheme:}

\begin{center}
\begin{tabularx}{\linewidth}{|X|c|}
\hline
\rowcolor{popblueL} \textbf{Component} & \textbf{Marks} \\
\hline
Introduction: role of banks in intermediation; structure of the Cameroonian sector & 3 \\
\hline
Problem 1: High non-performing loans (causes, consequences) & 3 \\
\hline
Problem 2: Excess liquidity paradox / credit rationing of SMEs & 3 \\
\hline
Problem 3: Weak collateral enforcement and judicial delays & 3 \\
\hline
Problem 4: High operating costs / infrastructure deficits & 2 \\
\hline
Problem 5: Regulatory and compliance burden & 2 \\
\hline
Problem 6: Competition from mobile money, fintech and MFIs & 2 \\
\hline
Solutions matched to problems (at least 4 concrete solutions) & 5 \\
\hline
Conclusion: balanced assessment & 2 \\
\hline
\rowcolor{popgreenL} \textbf{Total} & \textbf{25} \\
\hline
\end{tabularx}
\end{center}

\textbf{Model answer (developed outline):}

\textbf{Introduction.} Commercial banks are the engine of financial intermediation: they collect savings and transform them into credit for investment and growth. In Cameroon the sector is an \emph{oligopoly} dominated by a few large, mostly foreign-owned banks (Afriland, SGBC, BICEC, Ecobank, SGC...), alongside MFIs and mobile money operators. Despite substantial resources, these banks face serious problems that limit their contribution to development.

\textbf{Problem 1 — High non-performing loans (NPLs).} A significant share of loans is not repaid, because of poor ex-ante appraisal, weak borrower information, economic shocks (fall in cocoa/oil prices, crises in the North-West/South-West and Far North) and deliberate default. NPLs erode bank profits and capital, and make banks excessively cautious.

\textbf{Problem 2 — The excess liquidity paradox.} Cameroonian banks are highly liquid, yet SMEs and farmers complain of \emph{credit rationing}. Banks prefer risk-free Treasury bills and lending to the State or large corporations, because small borrowers lack guarantees, formal accounts and viable projects. Liquidity coexists with under-financing of the real economy.

\textbf{Problem 3 — Collateral and judicial weaknesses.} Most rural land is held under \emph{customary tenure} without land titles, so it cannot be mortgaged; and when default occurs, court procedures to seize and sell collateral are slow and uncertain. Risk therefore appears high, and interest rates and guarantee requirements rise accordingly.

\textbf{Problem 4 — High operating costs.} Electricity cuts, poor roads, costly ICT and security, and the need for own generators raise operating costs, compress margins and keep bank charges and lending rates high.

\textbf{Problem 5 — Regulatory and compliance burden.} Implementing COBAC's prudential framework (Basel-type capital rules), anti-money-laundering and counter-terrorism-financing obligations impose heavy compliance costs, especially on smaller institutions.

\textbf{Problem 6 — New competition.} Mobile money (MTN MoMo, Orange Money), fintechs and MFIs capture deposits and payment services, particularly among the young and the informal sector, eroding the traditional deposit base of banks.

\textbf{Solutions.}
\begin{itemize}
\item \emph{Against NPLs and information asymmetry:} operational \textbf{credit bureaus} and a shared database of borrowers; better staff training in the 5 Cs appraisal.
\item \emph{Against the collateral problem:} a \textbf{registry of movable collateral} (equipment, stocks, receivables) under OHADA law; faster commercial courts and arbitration; land-tenure reform to expand titling.
\item \emph{To finance SMEs:} specialised \textbf{SME desks}, partial \textbf{public guarantee funds}, and channelling excess liquidity into productive credit through refinancing windows.
\item \emph{To cut costs and face fintechs:} \textbf{digital transformation} (agency banking, mobile apps, partnerships with mobile money operators rather than confrontation).
\item \emph{Public action:} investment in electricity and transport infrastructure; proportionate, stable regulation to avoid unnecessary compliance costs.
\end{itemize}

\textbf{Conclusion.} CEMAC banking reforms (stronger COBAC supervision, higher capital requirements, modernised payment systems) have stabilised the sector, but structural problems — NPLs, collateral enforcement, SME exclusion — persist. Solving them requires joint action by banks (better appraisal, innovation), the State (justice, land reform, infrastructure) and the regulator (proportionate, incentive-based rules), so that the sector's abundant liquidity finally irrigates the Cameroonian economy.

\textbf{Examiner expectations:} a problem must always be \emph{explained} (cause $\rightarrow$ mechanism $\rightarrow$ consequence), not merely listed; solutions must be \emph{matched} to the problems; Cameroonian/CEMAC illustrations earn the application marks.
\end{corrige}

\begin{corrige}[Exercise 11 — Case Study: Credit Union Loan Appraisal (20 marks)]

\textbf{1. Appraisal with the 5 Cs [10 marks — 2 per C: rating + justification].}

\begin{itemize}
\item \textbf{Character — STRONG.} Five years of membership; \SI{400\,000}{FCFA} of savings; two previous loans (\SI{500\,000}{FCFA} and \SI{800\,000}{FCFA}) repaid \emph{early}; treasurer of her church women's group (community trust). The absence of a credit-bureau record is neutral in a rural area — her internal track record is excellent.

\item \textbf{Capacity — STRONG.} Net cooperative surplus: $1\,200\,000 - 700\,000 = \SI{500\,000}{FCFA}$/month. Even after household expenses (\SI{150\,000}{FCFA}), about \SI{350\,000}{FCFA}/month remains available, against a proposed repayment of \SI{100\,000}{FCFA}/month. Repayment absorbs only $20\%$ of the net surplus — very comfortable.

\item \textbf{Capital — ADEQUATE.} Own contribution of \SI{300\,000}{FCFA}, i.e. $\frac{300\,000}{1\,800\,000} = 16.7\%$ of project cost — acceptable (many MFIs require $10$--$20\%$), though not high. Cooperative assets of \SI{2\,000\,000}{FCFA} exceed the loan, showing a solid patrimony.

\item \textbf{Collateral — ADEQUATE.} No land title (community land) — the classic weakness. But the package is reasonable: chattel mortgage on equipment worth \SI{1\,200\,000}{FCFA} ($80\%$ of the loan), a solvent guarantor (civil servant, \SI{300\,000}{FCFA}/month salary) and pledged savings of \SI{400\,000}{FCFA}. Combined coverage: $1\,200\,000 + 400\,000 = 1\,600\,000 > 1\,500\,000$, plus the guarantee.

\item \textbf{Conditions — STRONG.} Growing demand for packaged garri, a new market 5 km away, stable input prices, no major competitors. Identified risks: seasonal cassava shortage (Jan--Mar) and poor road access in the rainy season — manageable with stocking and scheduling.
\end{itemize}

\textbf{2. Debt Service Coverage Ratio [3 marks].}
\[
\text{DSCR} = \frac{\text{Net Surplus}}{\text{Monthly Repayment}} = \frac{500\,000}{100\,000} = 5
\]
\textbf{Interpretation:} the cooperative generates \emph{five times} the cash needed to service the debt each month. A DSCR above $1.25$--$1.5$ is generally considered safe; $5$ is excellent. Even a severe shock (e.g. surplus halved during the cassava shortage) would leave a DSCR of $2.5$, still comfortable. (1 mark for the formula, 1 for the calculation, 1 for the interpretation.)

\textbf{3. Decision of the Credit Committee [7 marks].}

\textbf{Decision: APPROVE WITH CONDITIONS} (the file is fundamentally sound; conditions target the two identified risks: seasonality and collateral formalisation). Three binding covenants:

\begin{enumerate}
\item \textbf{Formalisation of guarantees before disbursement:} registration of the chattel mortgage on the grater and press, written and signed guarantee of the Cooperative President, and blocking (pledge) of the \SI{400\,000}{FCFA} savings for the duration of the loan.
\item \textbf{Controlled disbursement and use of funds:} payment made \emph{directly to the equipment suppliers} against invoices, and construction of the drying shed verified by a Credit Union officer before release of the final tranche — ensuring the \SI{1\,500\,000}{FCFA} finances the project and nothing else.
\item \textbf{Seasonal risk management:} obligation to build a \emph{cassava stock} (or dried chips) before January and to maintain a minimum savings buffer of one monthly instalment (\SI{100\,000}{FCFA}) in the pledged account, so that repayments continue during the Jan--Mar shortage and the rainy season.
\end{enumerate}

\textbf{Examiner expectations:} marks are awarded for (i) using the data \emph{quantitatively} (16.7\%, DSCR $= 5$, coverage $1\,600\,000$), (ii) balanced judgement (recognising the collateral weakness without rejecting a strong file), and (iii) conditions that are \emph{specific, binding and linked to the identified risks} — vague conditions such as ``she must repay well'' earn no credit.
\end{corrige}

\newpage

\coursec{Summary sheet}

\begin{popfigure}
\begin{center}\tcbox[colback=popink,colframe=popink,coltext=white,arc=9pt,boxsep=4pt,left=6pt,right=6pt]{\bfseries\small Money and Banking (II)}\end{center}
\vspace{-2pt}
\begin{tcbraster}[raster columns=3,raster equal height=rows,raster column skip=6pt,raster row skip=6pt,enhanced,blankest]
\begin{tcolorbox}[enhanced,colback=white,colframe=poporange,boxrule=0.9pt,arc=3pt,title={Bank Balance Sheet},fonttitle=\bfseries\scriptsize,coltitle=white,colbacktitle=poporange,left=4pt,right=4pt,top=2pt,bottom=2pt,boxsep=1pt,toptitle=1.5pt,bottomtitle=1.5pt,halign=left]
\begin{itemize}[nosep,leftmargin=*,itemsep=1pt,font=\scriptsize]
\item Assets: uses of funds
\item Liabilities: sources
\end{itemize}
\end{tcolorbox}
\begin{tcolorbox}[enhanced,colback=white,colframe=popgreen,boxrule=0.9pt,arc=3pt,title={Credit Creation},fonttitle=\bfseries\scriptsize,coltitle=white,colbacktitle=popgreen,left=4pt,right=4pt,top=2pt,bottom=2pt,boxsep=1pt,toptitle=1.5pt,bottomtitle=1.5pt,halign=left]
\begin{itemize}[nosep,leftmargin=*,itemsep=1pt,font=\scriptsize]
\item Credit multiplier $k=1/r$
\item Leakages limit it
\end{itemize}
\end{tcolorbox}
\begin{tcolorbox}[enhanced,colback=white,colframe=poppink,boxrule=0.9pt,arc=3pt,title={Problems of Banks in Cameroon},fonttitle=\bfseries\scriptsize,coltitle=white,colbacktitle=poppink,left=4pt,right=4pt,top=2pt,bottom=2pt,boxsep=1pt,toptitle=1.5pt,bottomtitle=1.5pt,halign=left]
\begin{itemize}[nosep,leftmargin=*,itemsep=1pt,font=\scriptsize]
\item Bad debts, low savings
\end{itemize}
\end{tcolorbox}
\begin{tcolorbox}[enhanced,colback=white,colframe=popgold,boxrule=0.9pt,arc=3pt,title={Central Bank (BEAC)},fonttitle=\bfseries\scriptsize,coltitle=white,colbacktitle=popgold,left=4pt,right=4pt,top=2pt,bottom=2pt,boxsep=1pt,toptitle=1.5pt,bottomtitle=1.5pt,halign=left]
\begin{itemize}[nosep,leftmargin=*,itemsep=1pt,font=\scriptsize]
\item Issue, banker to banks
\item Monetary policy tools
\end{itemize}
\end{tcolorbox}
\begin{tcolorbox}[enhanced,colback=white,colframe=poppurple,boxrule=0.9pt,arc=3pt,title={Microfinance},fonttitle=\bfseries\scriptsize,coltitle=white,colbacktitle=poppurple,left=4pt,right=4pt,top=2pt,bottom=2pt,boxsep=1pt,toptitle=1.5pt,bottomtitle=1.5pt,halign=left]
\begin{itemize}[nosep,leftmargin=*,itemsep=1pt,font=\scriptsize]
\item MC2, credit unions
\item Financial inclusion
\end{itemize}
\end{tcolorbox}
\begin{tcolorbox}[enhanced,colback=white,colframe=popdark,boxrule=0.9pt,arc=3pt,title={Loan Appraisal},fonttitle=\bfseries\scriptsize,coltitle=white,colbacktitle=popdark,left=4pt,right=4pt,top=2pt,bottom=2pt,boxsep=1pt,toptitle=1.5pt,bottomtitle=1.5pt,halign=left]
\begin{itemize}[nosep,leftmargin=*,itemsep=1pt,font=\scriptsize]
\item The 5 Cs of credit
\end{itemize}
\end{tcolorbox}
\end{tcbraster}

\legende{Mind map of the chapter: from the commercial bank's balance sheet to loan appraisal by credit unions.}
\end{popfigure}

\begin{bilan}
\textbf{1. Key definitions}
\begin{itemize}
\item \cle{Balance sheet}: a financial statement showing, at a given date, a bank's \textbf{assets} (uses of funds: cash, loans, investments, premises) and \textbf{liabilities} (sources of funds: deposits, capital, reserves, borrowings). Assets always equal liabilities.
\item \cle{Credit creation}: the process by which commercial banks create \emph{deposit money} by granting loans, since a loan becomes a new deposit elsewhere in the banking system.
\item \cle{Central bank}: the apex monetary institution (e.g. \cle{BEAC} for the CEMAC zone) that issues currency, acts as banker to the government and to commercial banks, and conducts monetary policy.
\item \cle{Microfinance}: the provision of small-scale financial services (savings, microcredit, micro-insurance) to low-income people excluded from traditional banks; in Cameroon, mainly through \cle{credit unions} and MC\textsuperscript{2} institutions.
\item \cle{Loan appraisal}: the systematic assessment of a borrower's file before granting credit, based on the \textbf{5 Cs}: Character, Capacity, Capital, Collateral, Conditions.
\end{itemize}

\textbf{2. Formulas and laws to know}
\begin{itemize}
\item Balance sheet identity: $\text{Total Assets} = \text{Total Liabilities} + \text{Capital}$.
\item Simple credit multiplier: $k = \dfrac{1}{r}$, where $r$ is the reserve (liquidity) ratio.
\item Total deposits created: $\Delta D = \dfrac{\Delta D_0}{r}$; total credit created: $\Delta C = \Delta D - \Delta D_0$.
\item Example: with $r = 20\%$ and an initial deposit of $100\,000$ FCFA, $k = 5$, total deposits $= 500\,000$ FCFA, credit created $= 400\,000$ FCFA.
\item Monetary policy tools of BEAC: \cle{key interest rate}, \cle{reserve requirements}, \cle{open market operations}, refinancing (rediscount) facilities.
\end{itemize}

\textbf{3. Methods}
\begin{itemize}
\item \emph{Presenting a balance sheet}: classify items into assets (ordered by decreasing liquidity) and liabilities (ordered by increasing exigibility); check the equality of both sides.
\item \emph{Computing credit creation}: identify $r$, compute $k = 1/r$, multiply the initial deposit by $k$, then subtract the initial deposit to obtain credit created. Always mention leakages (cash withdrawals, imports, excess reserves) which reduce the multiplier in practice.
\item \emph{Appraising a loan file (5 Cs)}: verify the borrower's \textbf{Character} (reputation, repayment history), \textbf{Capacity} (income vs. repayment instalments), \textbf{Capital} (own contribution), \textbf{Collateral} (guarantees), and \textbf{Conditions} (purpose, economic environment).
\end{itemize}

\textbf{4. Common mistakes to avoid}
\begin{itemize}
\item Confusing \emph{assets} with \emph{liabilities}: deposits are a \textbf{liability} for the bank (money owed to customers), loans granted are an \textbf{asset}.
\item Believing one bank alone creates the full multiplier: credit creation works through the \textbf{banking system as a whole}.
\item Forgetting leakages: the formula $k = 1/r$ gives a \emph{maximum}; real creation is smaller.
\item Confusing BEAC (central bank, CEMAC zone) with commercial banks like Afriland, SGBC or BICEC.
\item Thinking microfinance charges low interest: rates are often \emph{higher} because of risk and small loan sizes.
\end{itemize}

\textbf{5. GCE tips}
\begin{itemize}
\item In essay questions on credit creation, always \textbf{define}, give the \textbf{formula}, a \textbf{numerical illustration}, then \textbf{limits} (leakages, demand for loans, BEAC controls).
\item For BEAC, structure answers around its \textbf{functions} (issue of currency, banker to banks, banker to the State, guardian of the CFA franc, monetary policy) and mention the COBAC supervisory role.
\item Use Cameroonian examples naturally: MC\textsuperscript{2} networks, CamCCUL, village savings groups (\emph{njangui}/tontines) as informal microfinance.
\item For problems of banks in Cameroon, cite: low savings and banking rate, bad debts, poor infrastructure, competition from microfinance and mobile money, liquidity concentration in Douala and Yaound\'e.
\end{itemize}
\end{bilan}

\findecours

\end{document}
